%=====================================================================
%  main.tex —— 范畴表示论与局部几何朗兰兹导论（中文译本）
%  底本：Gurbir Dhillon, arXiv:2205.14578v1, 83 页, 2022
%=====================================================================
\documentclass[UTF8, a4paper, 12pt, fontset=fandol, scheme=chinese]{ctexart}

%=====================================================================
%  preamble.tex —— 中文译本的版式与定理环境设定
%  底本：Gurbir Dhillon, arXiv:2205.14578（amsart, 10pt）
%  目标：在中文 ctexart 12pt 下复现原文的编号体制与版面层次
%=====================================================================

%--------------------- 版面 ---------------------
\usepackage{geometry}
\geometry{a4paper, left=2.7cm, right=2.7cm, top=2.9cm, bottom=2.9cm}

\usepackage{fancyhdr}
\pagestyle{fancy}
\fancyhf{}
\fancyhead[L]{\small\kaishu 范畴表示论与局部几何朗兰兹导论}
\fancyhead[R]{\small\thepage}
\renewcommand{\headrulewidth}{0pt}
\fancypagestyle{plain}{\fancyhf{}\fancyfoot[C]{\small\thepage}%
  \renewcommand{\headrulewidth}{0pt}}

%--------------------- 数学与图 ---------------------
\usepackage{amsmath}
\usepackage{amssymb}
\usepackage{amsthm}
\usepackage{amscd}
\usepackage{mathtools}
\usepackage{array}
\usepackage{graphicx}
\usepackage{tabularx}
\usepackage{float}
\usepackage{xspace}
\usepackage{xcolor}
\usepackage{tocvsec2}
\usepackage[mathscr]{euscript}
\usepackage{scalerel}
\usepackage{aligned-overset}
% xy-pic 用于原文的 \xymatrix 交换图（原文载入 [all,cmtip]）
\usepackage[all,cmtip]{xy}

%--------------------- 编号体制 ---------------------
% 原文：amsart + \numberwithin{equation}{section}；
%       定理类环境一律挂 subsubsection 计数器（印刷为「2.3.10.」形式）
\numberwithin{equation}{section}

%--------------------- 定理环境 ---------------------
% 对应 amsart 的 plain 样式：题头粗体，正文斜体，题头以句点收尾，
% optional 注排在编号之后、句点之前，形如「定理 3.3.13 (Beilinson–Bernstein, [BB81]).」
\newtheoremstyle{zhthm}%
  {3pt}{3pt}% 上下间距
  {\itshape}% 正文
  {}% 缩进
  {\bfseries}% 题头
  {.}% 题头后标点
  {0.5em}% 题头后空隙
  {\thmname{#1}\thmnumber{ #2}\thmnote{ (#3)}}
\newtheoremstyle{zhdefn}%
  {3pt}{3pt}{\normalfont}{}{\bfseries}{.}{0.5em}%
  {\thmname{#1}\thmnumber{ #2}\thmnote{ (#3)}}

% 环境名沿用原文，避免正文大量改名；仅把显示名换成中文。
\theoremstyle{zhthm}
\newtheorem{theo}[subsubsection]{定理}
\newtheorem*{theo*}{定理}
\newtheorem{lem}[subsubsection]{引理}
\newtheorem{lemma}[subsubsection]{引理}
\newtheorem{cor}[subsubsection]{推论}
\newtheorem{pro}[subsubsection]{命题}
\newtheorem{conj}[subsubsection]{猜想}
\newtheorem{re}[subsubsection]{注}
\newtheorem{remark}[subsubsection]{注}
\newtheorem{example}[subsubsection]{例}
\newtheorem{ex}[subsubsection]{例}
\newtheorem{proofw}[subsubsection]{定理的证明}

\theoremstyle{zhdefn}
\newtheorem{defn}[subsubsection]{定义}

\renewcommand{\proofname}{证明}

% 原文用 \subsubsection{} 产生「2.3.5.」式的段首编号：编号不加粗、首行不缩进。
\newcommand{\ssszh}{%
  \par\medskip\noindent\refstepcounter{subsubsection}%
  \thesubsubsection.\enspace\ignorespaces}

%--------------------- 译注 ---------------------
% 原文脚注 37 条照译为普通脚注；译者所加注释另立序列，印刷为「译1」「译2」……
\usepackage{manyfoot}
\DeclareNewFootnote[plain]{trans}
\renewcommand{\thefootnotetrans}{译\arabic{footnotetrans}}
\newcommand{\zy}[1]{\footnotetrans{#1}}

%--------------------- 标题层次 ---------------------
\ctexset{
  section={
    name={,},
    number=\arabic{section},
    aftername={.\enspace},
    format=\Large\bfseries\centering,
    beforeskip=3ex plus 1ex minus .2ex,
    afterskip=2ex plus .2ex,
  },
  subsection={
    aftername={.\enspace},
    format=\large\bfseries\raggedright,
    beforeskip=2.5ex plus .5ex minus .2ex,
    afterskip=1.2ex plus .2ex,
  },
}

% 附录：原文印刷为「Appendix A. From functions on …」，小节自然成为「A.1.」
\newcommand{\zhappendix}{%
  \appendix
  \ctexset{section={name={Appendix\ }, number=\Alph{section}}}}

%--------------------- 中西文混排细节 ---------------------
\xeCJKsetup{CJKecglue={\hskip 0.24em plus 0.05em minus 0.05em}}
\setlength{\parindent}{2em}
\linespread{1.22}
\allowdisplaybreaks[1]

%--------------------- 超链接与书签 ---------------------
\usepackage[colorlinks=true, linkcolor=blue!55!black, citecolor=blue!55!black,
            urlcolor=blue!55!black, bookmarksnumbered=true]{hyperref}

%--------------------- 移植自原文导言区的宏 ---------------------
% 由 工程文件/scripts/port_macros2.py 从原文导言区自动移植，勿手改。
\newcommand{\pt}{\on{pt}}
\newcommand{\IC}{\on{IndCoh}}
\newcommand{\ICn}{\on{IndCoh}_{\on{nilp}}}
\newcommand{\fgc}{\widehat{\fg}_{\kappa_c}}
\newcommand{\cB}{\check{B}}
\newcommand{\cN}{\check{N}}
\newcommand{\fgk}{\widehat{\fg}_\kappa}
\newcommand{\Op}{\mathrm{Op}_{\cG}}
\newcommand{\Fp}{\mathbb{F}_p}
\newcommand{\gkc}{\gk_{\kappa_c}}
\newcommand{\TwoCoh}{\on{2-QCoh}}
\newcommand{\TwoIC}{\on{2-IndCoh}_{\on{nilp}}}
\newcommand{\LS}{\on{LocSys}_{\cG}}
\newcommand{\Fqt}{\Fq(\!(t)\!)}
\newcommand{\uHom}{\underline{\Hom}}
\newcommand{\udHom}{\on{Hom}_{\DGCat}}
\newcommand{\sB}{\on{Fl}_G}
\newcommand{\sBumba}{\on{Fl}_G}
\newcommand{\sM}{\mathscr{M}}
\newcommand{\fg}{\mathfrak{g}}
\newcommand{\fh}{\mathfrak{h}}
\renewcommand{\mod}{\on{-mod}}
\newcommand{\Dmod}{\on{D-mod}}
\newcommand{\Fq}{\mathbb{F}_q}
\newcommand{\Fun}{\on{Fun}}
\newcommand{\Vect}{\on{Vect}}
\newcommand{\rmod}{\on{mod-}}
\newcommand{\fk}{\mathfrak{k}}
\newcommand{\fkk}{\fk_\kappa}
\newcommand{\sL}{\mathscr{L}}
\newcommand{\sA}{\mathscr{A}}
\newcommand{\sT}{\mathscr{T}}
\newcommand{\sV}{\mathscr{V}}
\newcommand{\sU}{U}
\newcommand{\Usi}{\mathscr{U}^{\scaleto{\frac{\infty}{2}}{7pt}}}
\newcommand{\HCh}{\on{HCh}}
\newcommand{\HChl}{\on{HCh}^{\on{tame}}}
\newcommand{\sP}{\mathscr{P}}
\newcommand{\per}{\Pos_{\hspace{.2mm} \kappa}}
\newcommand{\hk}{\widehat{\fh}_\kappa}
\newcommand{\Pos}{\on{Pos}}
\newcommand{\wt}{\widetilde}
\newcommand{\sS}{\mathscr{S}}
\newcommand{\sZ}{\mathscr{Z}}
\newcommand{\starb}{\overset{}{\star}}
\newcommand{\cgk}{\widehat{\check{\mathfrak{g}}}_{\check{\kappa}}}
\newcommand{\Sp}{\on{Spec }}
\newcommand{\aF}{\mathsf{F}}
\newcommand{\HK}{\sH}
\newcommand{\aG}{\mathsf{G}}
\newcommand{\sH}{\mathscr{H}}
\newcommand{\cl}{\mathcal{L}}
\newcommand{\sE}{\mathscr{E}}
\newcommand{\conv}{\operatorname{conv}}
\newcommand{\Mod}{\operatorname{-mod}}
\newcommand{\Ind}{\operatorname{Ind}}
\newcommand{\Res}{\operatorname{Res}}
\newcommand{\cone}{\operatorname{cone}}
\newcommand{\Fin}{\operatorname{Fin}}
\newcommand{\supp}{\operatorname{supp}}
\newcommand{\spa}{\operatorname{span}}
\newcommand{\ch}{\operatorname{ch}}
\newcommand{\relint}{\operatorname{rel int}}
\newcommand{\aff}{\operatorname{aff}}
\newcommand{\St}{\on{s}}
\newcommand{\mult}{\operatorname{mult}}
\newcommand{\End}{\operatorname{End}}
\newcommand{\Hom}{\operatorname{Hom}}
\newcommand{\CC}{\mathbb{C}}
\newcommand{\gr}{\text{gr}}
\newcommand{\sN}{\mathscr{N}}
\newcommand{\g}{\widehat{\mathfrak{g}}}
\newcommand{\n}{\mathfrak{n}}
\newcommand{\p}{\partial}
\newcommand{\fL}{\mathfrak{L}}
\newcommand{\fl}{\mathfrak{l}}
\newcommand{\Gr}{\operatorname{Gr}}
\newcommand{\fu}{\mathfrak{u}}
\newcommand{\fp}{p}
\newcommand{\fs}{\mathfrak{s}}
\newcommand{\halpha}{\check{\alpha}}
\newcommand{\hpi}{\check{\prod}}
\newcommand{\hrho}{\check{\rho}}
\newcommand{\fsl}{\mathfrak{sl}}
\newcommand{\hch}{\fg\operatorname{-mod}^{B, \heartsuit}}
\newcommand{\hc}{\fg\operatorname{-mod}^B}
\newcommand{\Ug}{\hat{U}\gk}
\newcommand{\ft}{\mathfrak{t}}
\newcommand{\x}{{\bf x}}
\newcommand{\fz}{\mathfrak{z}}
\newcommand{\corank}{\text{corank}}
\newcommand{\ca}{\mathcal{A}}
\newcommand{\id}{\text{id}}
\newcommand{\sF}{\mathscr{F}}
\newcommand{\sG}{\mathscr{G}}
\newcommand{\tr}{\operatorname{tr}}
\newcommand{\da}{\dot{a}_{\on{reg}}}
\newcommand{\disp}{\overset{\rightarrow}{\ell}}
\newcommand{\V}{\ensuremath{\mathbb{V}}}
\newcommand{\Bun}{\operatorname{Bun}}
\newcommand{\Fl}{\operatorname{Fl}}
\newcommand{\dy}{\mathbf{Dy}(\gk)}
\newcommand{\bb}{\mathfrak{b}}
\newcommand{\vla}{\ensuremath{\mathbb{V}^\lambda}}
\newcommand{\R}{\ensuremath{\mathbb{R}}}
\newcommand{\C}{\ensuremath{\mathbb{C}}}
\newcommand{\B}{\ensuremath{\mathcal{B}}}
\newcommand{\hoom}{\underline{\operatorname{Hom}}}
\newcommand{\Q}{\ensuremath{\mathbb{Q}}}
\newcommand{\bba}{\ensuremath{\mathbb{A}}}
\newcommand{\scrf}{\ensuremath{\mathscr{F}}}
\newcommand{\M}{\ensuremath{\mathbb{M}}}
\newcommand{\Z}{\ensuremath{\mathbb{Z}}}
\newcommand{\calp}{\ensuremath{\mathcal{P}}\xspace}
\newcommand{\liehr}{\ensuremath{\mathfrak{h}_{\mathbb{R}}}}
\newcommand{\wtvla}[1]{\ensuremath{\wt_{#1} L(\lambda)}}
\newcommand{\lie}[1]{\ensuremath{\mathfrak{#1}\xspace}}
\newcommand{\comment}[1]{}
\newcommand{\ms}{\vspace{.3 cm}}
\newcommand{\OO}{\mathscr{O}}
\newcommand{\KK}{\mathcal{K}}
\newcommand{\Spec}{\text{Spec }}
\newcommand{\Sym}{\operatorname{Sym}}
\newcommand{\vv}{\lvert 0\rangle}
\newcommand{\och}{\underline{\text{ch}}}
\newcommand{\FF}{\mathbb{F}}
\newcommand{\sinfty}{\scaleto{\frac{\infty}{2}}{9pt}}
\newcommand{\fri}{\mathfrak{i}}
\newcommand{\Ad}{\mathbb{A}}
\newcommand{\F}{\mathcal{F}}
\newcommand{\vk}{V_\kappa}
\newcommand{\Ext}{\operatorname{Ext}}
\newcommand{\IW}{\mathsf{IWhit} }
\newcommand{\sca}{\mathscr{A}}
\newcommand{\fb}{\mathfrak{b}}
\newcommand{\gk}{\widehat{\fg}}
\newcommand{\fn}{\mathfrak{n}}
\newcommand{\fm}{\mathfrak{m}}
\newcommand{\ce}{\mathcal{E}}
\newcommand{\bp}{\bar{\p}}
\newcommand{\sD}{\mathscr{D}}
\newcommand{\VV}{\mathbb{V}}
\newcommand{\ski}{\vspace{.2cm}}
\newcommand{\bk}{\backslash}
\newcommand{\A}{\mathbb{A}}
\newcommand{\Ran}{\operatorname{Ran}}
\newcommand{\wOp}{\widehat{Op}}
\newcommand{\scc}{\mathscr{C}}
\newcommand{\hj}{\hspace{.07cm} \widehat{\hspace{-.07cm} j}}
\newcommand{\ftp}{D_\kappa}
\newcommand{\seq}{\hspace{1mm} \simeq \hspace{1mm}}
\newcommand{\h}{\widehat{\fh}}
\newcommand{\hs}{\hspace{.5mm}}
\newcommand{\Hmod}{H\mod^{\on{weak, naive}}}
\newcommand{\rhoc}{\check{\rho}}
\newcommand{\lies}{\operatorname{Lie-}\star}
\newcommand{\normord}[1]{:\mathrel{#1}:}
\newcommand{\fgl}{\mathfrak{gl}}
\newcommand{\fagl}{\widehat{\fgl}}
\newcommand{\cw}{\mathscr{W}}
\newcommand{\cfg}{\check{\fg}}
\newcommand{\RHom}{\operatorname{RHom}}
\newcommand{\Flg}{\operatorname{Fl}_G}
\newcommand{\Wf}{{}^fW}
\newcommand{\flg}{\Flg}
\newcommand{\gmk}{\widehat{\fg}_{-\kappa}}
\newcommand{\ck}{\underline{k}}
\newcommand{\acg}{{{\cfg}_{\ck}}}
\newcommand{\cg}{\check{G}}
\newcommand{\sO}{\mathscr{O}}
\newcommand{\sJ}{\mathscr{J}}
\newcommand{\cG}{\check{G}}
\newcommand{\cC}{\mathscr{C}}
\newcommand{\sC}{\mathscr{C}}
\newcommand{\sHom}{\on{Hom}}
\newcommand{\inv}{\msf{inv}}
\newcommand{\msf}{\operatorname}
\newcommand{\cI}{\check{I}}
\newcommand{\ur}{\overset{\circ}{I}}
\newcommand{\sW}{\EuScript{W}_\kappa}
\newcommand{\cW}{\mathscr{W}}
\newcommand{\Avr}{\on{Av}_{*}^{\psi, r}}
\newcommand{\sI}{\hs \overset I \star \hs }
\newcommand{\crho}{\check{\rho}}
\newcommand{\Qlb}{\overline{\mathbb{Q}}_\ell}
\newcommand{\Oreg}{\g\on{-mod}_{\on{reg}}^I}
\newcommand{\vc}{\mathbb{V}_{\on{crit}}}
\newcommand{\Whit}{\mathsf{Whit}}
\newcommand{\QCoh}{\on{QCoh}}
\newcommand{\Rep}{\mathsf{Rep}}
\newcommand{\xar}[1]{\xrightarrow{#1}}
\newcommand{\DGCat}{\msf{DGCat}_{\on{cont}}}
\newcommand{\on}{\operatorname}



\hypersetup{
  pdftitle={范畴表示论与局部几何朗兰兹导论},
  pdfauthor={Gurbir Dhillon 著；数译译},
  pdfsubject={arXiv:2205.14578 中文译本},
  pdfkeywords={范畴表示论, 局部几何朗兰兹, Hecke 范畴, 仿射 Beilinson--Bernstein 局部化}
}

\begin{document}

%--------------------- 扉页 ---------------------
\begin{titlepage}
\centering
\vspace*{3.2cm}

{\LARGE\bfseries 范畴表示论与局部几何朗兰兹导论\par}

\vspace{0.9cm}
{\large An Informal Introduction to Categorical Representation Theory\\
and the Local Geometric Langlands Program\par}

\vspace{2.2cm}
{\large Gurbir Dhillon\par}
\vspace{0.35cm}
{\normalsize Yale University\par}

\vspace{2.4cm}
{\large 中 文 译 注 本\par}

\vspace{0.5cm}
{\normalsize 底本：arXiv:2205.14578v1（2022 年 5 月 29 日），共 83 页\par}

\vfill
{\normalsize 数译（DSH agent）译\par}
{\normalsize 版本 v1.0\quad 2026 年 11 月\par}
\vspace{1.2cm}
\end{titlepage}

%--------------------- 摘要 ---------------------
\vspace*{0.5cm}
\begin{center}
{\large\bfseries 摘\quad 要}
\end{center}
\vspace{0.3cm}

\noindent 本文以讲清动机的方式，介绍群的范畴作用的理论与局部几何朗兰兹纲领。行文之中，
着重说明这些想法在约化李代数与仿射李代数的通常表示论中的新旧应用。

\vspace{0.8cm}
\noindent{\small 原文日期：2022 年 5 月 31 日。}

\newpage

%--------------------- 译者说明 ---------------------
\section*{译者说明}
\addcontentsline{toc}{section}{译者说明}

\begin{enumerate}
\item \textbf{底本。} 本书译自 arXiv:2205.14578v1（2022 年 5 月 29 日提交，正文 83 页），
  作者 Gurbir Dhillon，原题 \emph{An Informal Introduction to Categorical Representation
  Theory and the Local Geometric Langlands Program}。作者在标题里自称「非正式的导论」，
  书中多以例子、类比与历史线索推进，证明常常只给思路，细节留给出处的文献。译文一仍其旧，
  既不补证，也不删减。

\item \textbf{编号。} 译文沿用原文的编号体制。节、小节、段首编号（形如 2.3.5.）、公式编号
  （形如 (2.1)）与定理类环境的编号（形如 定理 3.2.10.）都与原文一致。原文把定理类环境
  一律挂在段首计数器上，故定理、命题、注、例共用一套三位编号，译文照此办理。所有交叉引用
  都由 \verb|\ref| 生成，没有手写编号。

\item \textbf{定理类环境。} 原文的 Theorem、Proposition、Lemma、Corollary、Conjecture、
  Definition、Remark、Example 依次译作定理、命题、引理、推论、猜想、定义、注、例。原文有时
  把作者名与文献写在环境正文的开头，有时写在环境的可选注内，两种排法都照原样保留。

\item \textbf{译注。} 脚注中不带「译」字的，是原文脚注的照译；带「译」字并另行编号的，是译者
  所加。译者所加的解释性注释一律标明来源：有外部依据的写作「[来源：……]」，只能由上下文
  推定的写作「[译者按（上下文推断，无外部依据）：……]」。脚注之外的补充说明，随文以
  「译者按」标出。

\item \textbf{术语。} 关键术语首次出现时，一般在括号内附原文；全书译名见
  \texttt{资料卡/04-术语文献勘误.md} 与 \texttt{工程文件/work/glossary.tsv}。D-模、反常层、
  Hecke 范畴、带钉、抛孔、仿射 Grassmann 簇等，沿用本丛书已冻结的译名。

\item \textbf{勘误。} 原文有 3 处明显的笔误，例如重复的冠词与拼错的 theory。译文按正确含义
  翻译，并在资料卡中登记，正文不作改动标记。原文个别术语的用法较宽松（如把任意非
  Archimedes 局部域上的取值都称作 $p$-adic group），译文保留原样，必要时加译注说明。

\item \textbf{图表。} 原文没有插图，只有 5 幅交换图与 3 个表格。交换图用 xy-pic 重排，
  版式与原文一致。

\item \textbf{致谢。} 翻译中的疏漏与不妥，概由译者负责。读者若发现误译，欢迎指出。
\end{enumerate}

\vspace{1em}
\noindent 数译（DSH agent）\\
2026 年 11 月

\newpage
\setcounter{tocdepth}{1}
\tableofcontents
\newpage

%=====================================================================
%  ch01 —— 原文 297–496：题名、摘要、§1 引言、致谢
%=====================================================================

\section{引言}

\subsection{} 群的范畴作用理论与局部几何朗兰兹对应，是现代表示论中两个重要的角色。
过去二十年里，它们发展迅速：Beilinson、Drinfeld、Frenkel、Gaitsgory 作出了决定性的贡献，
更晚近则有 Beraldo、Raskin、Yang。

无论从历史看，还是从主要应用看，范畴作用与局部几何朗兰兹都同李代数的表示论深深交织。
这两个联系，本书会反复强调。

一方面，范畴作用的理论为处理约化李代数的表示论提供了高效的框架。这也是进入范畴作用理论的
一条便利路径，本书正是由此入手。

另一方面，也是更重要的方面，范畴作用的形式体系与局部几何朗兰兹的视角，对仿射李代数及相关
顶点代数的表示论发展不可或缺。它们既提供了解决旧猜想的工具，也催生了新的猜想。这些进展将在
本书后半部分讲到。

\subsection{} 这一圈想法有一个美妙之处：尽管表面上用到了现代代数与几何中相当精致的工具，
只要讲法得当，非专家也能迅速吸收其中的基本结果、技巧与视角。我们希望本书能为非专家读者
提供这样一个人口。

为此，我们给自己定下的目标是：以可读的方式综述这一领域，并补足文献中已有的材料。出于这一
考虑，我们选择着重讲类比、动机与历史前身，而不去对所涉材料作完全严格的展开。

\subsection{} 追求这样一个目标，立刻会遇到两个问题。

第一个问题是：假定读者已经知道多少，又该把多少内容讲得细一些。我们假定读者对单李群与单李
代数的表示论、代数几何、同调代数的基础有所了解，但不超过这些科目入门课程的内容。当然，
这些预备也可以换成一定的数学成熟度，或者至少是愿意把若干断言先信下来。特别地，我们不假定
读者熟悉几何表示论的工具，例如 D-模与反常层、高阶代数，以及导出与无穷维代数几何。需要什么，
我们就在用到时发展什么。这样做的同时，我们仍力求以最少的枝节传达基本想法，希望读者借此
获得方向感，再去攻读有关文献。

第二个问题是：强调哪些想法、方法与直觉，必然受篇幅限制，也受综述者本人偏好与局限的限制。
因此，对于许多贡献被略去或强调不足的数学家与物理学家，我们在此致歉。先说一点：本书与三维、
四维超对称规范理论之间极其丰富的联系，只在末尾略作提及。

\subsection{} 我们努力使全书的材料复杂度单调递增。读者可以从头读起，也可以跳到觉得陌生
或有趣之处，从那里往下读。我们希望前几节，特别是第 \ref{s:bbloc} 节，能作为非专家进入
几何表示论的一篇轻量导引。至于最后几节，特别是第 \ref{s:affbbloc} 节，我们也预计，要领会
其中每一个细节，需要一定程度的预备知识。

本书的具体安排如下。第 \ref{s:repthryfingroup} 节回顾有限群表示论的若干基本方面，为后面
遇到的范畴对应物作铺垫。第 \ref{s:bbloc} 节回顾 Beilinson--Bernstein 局部化定理，并说明
由此得到的等变 D-模与该李代数的 Harish-Chandra 模之间的对应，如何引出范畴群作用的想法。
第 \ref{s:catrepgrp} 节讨论有限维代数群的范畴群作用理论的基础，并花一些篇幅解释现代同调代数
中若干必要的工具与构造，好让读者能够上手。要走到局部几何朗兰兹，必须从代数群的情形过渡到
环路群。因此，第 \ref{s:catloop} 节建立环路群的范畴表示论的基础，也就是检讨进入无穷维情形
后哪些地方要变。最后，第 \ref{s:affbbloc} 节讲局部几何朗兰兹对应，较详细地讨论其中已知与
猜想的部分，并着重介绍它在临界水平仿射李代数的局部化理论中的一个重要应用。附录
\ref{s:frmfun} 还为非专家提供了一份简明导引，讲 D-模、Riemann--Hilbert 对应与函数-层对应。

\subsection{} 除书中各处给出的文献指引外，读者还可参阅关于局部几何朗兰兹的两部出色综述
\cite{frenkellocalnotes} 与专著 \cite{frenkelloop}，以及讲义 \cite{paris-notes}；关于整体
图景，\cite{frenkelglobalnotes} 是一份通俗的导引。

\vspace{.5cm}

\noindent{\bf 致谢。} 感谢 Pramod Achar、Dima Arinkin、Dario Beraldo、Roman Bezrukavnikov、
Alexander Braverman、Justin Campbell、Harrison Chen、Thomas Creutzig、Davide Gaiotto、
Tony Feng、Edward Frenkel、Igor Frenkel、Dennis Gaitsgory、Ivan Losev、Sam Raskin、
Daniel Sage、Germ\'an Stefanich、Jonathan Wang、David Yang、Zhiwei Yun、Xinwen Zhu，
与他们的讨论加深了我们对本课题的理解。我们还要深深感谢 Cora Bernard、Apoorva Khare
以及匿名审稿人：他们仔细阅读了本文早先的稿本，并提出了有益的意见。

作者的工作由美国国家科学基金会（NSF）的博士后资助支持，资助号 2103387。

%=====================================================================
%  ch02 —— 原文 497–647：§2 群代数与卷积代数（2.1 概览、2.2 卷积与群代数）
%=====================================================================

\section{群代数与卷积代数的表示}
\label{s:repthryfingroup}
\subsection{概览}

\ssszh 设 $G$ 是整数环上的分裂约化群，\footnote{读者若通篇把 $G$ 取为可逆
$n \times n$ 矩阵之群 $GL_n$，甚至只取 $GL_2$，损失也不大。} 它在有限域、局部域与实数上
的点，其复表示论之间有富于启发的类比。借此可以把构造与想法从一个领域搬到另一个领域。
Harish-Chandra 甚至把这种彼此渗透称作表示论的 Lefschetz 原理
\cite{langlands1985harish}。\footnote{回想通常的 Lefschetz 原理：对任意复代数簇成立的
事实，例如一般光滑性，对任意特征零的代数闭域上的代数簇也成立。}

类似地，范畴理论中有一条基本原理，其源头是 Drinfeld 与 Laumon 关于整体几何朗兰兹对应的
工作，以及 Lusztig 关于特征标层的工作。这条原理说：$G(\mathbb{C})$ 的范畴表示论，其行为
应当与 $G(\mathbb{F}_q)$ 的复表示论（的一部分）相似；$G(\mathbb{C}(\!(t)\!))$ 的范畴
表示论，其行为应当与 $G(\mathbb{F}_q(\!(t)\!))$ 的复表示论相似。

\ssszh 有鉴于此，在进入范畴情形之前，我们先回顾通常有限群表示论的若干要点，然后把讨论
特殊化到 $G(\mathbb{F}_q)$。

\subsection{通过卷积看群代数}

\ssszh 设 $H$ 是有限群，$k$ 是特征零的代数闭域。用 $\on{Rep}(H)^\heartsuit$ 记 $H$ 在
$k$-向量空间上的表示所构成的阿贝尔范畴。

\begin{re}
$\on{Rep}(H)^\heartsuit$ 上的上标，是为了强调我们说的是阿贝尔范畴，即相应导出范畴
$\on{Rep}(H)$ 上标准 $t$-结构的核。类似地，下文中出现的模范畴与层范畴，除另有说明外
都指导出版本；相应的阿贝尔范畴则以上标「$\heartsuit$」标记。
\end{re}

\ssszh 回忆：给 $H$ 在向量空间 $V$ 上的一个表示，就是给一个映射
%
\begin{equation} \label{e:hom}
\alpha: H \rightarrow \on{End}(V).
\end{equation}
%
这里 $\on{End}(V)$ 指 $V$ 的 $k$-线性自同态，而 $\alpha$ 是单位的，并把 $H$ 中的乘法与
$\on{End}(V)$ 中的复合对应起来。更一般地，对任意 $k$-代数 $A$，考虑半群映射的集合
%
\begin{equation}  \label{e:homalg}
H \rightarrow A,
\end{equation}
%
其中 $A$ 以乘法看作半群。一个关键的观察是：由于我们丢掉了 $A$ 上的结构，即加法，
这可以扩充为从由 $H$ 造出的 $k$-代数（即群代数）出发的映射。下面我们换一种说法来讲群代数
的构造，以便适应后文要遇到的范畴变体。

\ssszh \label{ss:opsfinitesets}对有限集 $X$，用 $\on{Fun}(X)$ 记 $X$ 上 $k$-值函数
构成的向量空间。回忆这个对应有两条基本性质。第一，$X$ 分解为乘积 $X = X_1 \times X_2$，
给出其函数空间的一个典范张量积分解
$$ \Fun(X) \simeq \Fun(X_1) \otimes \Fun(X_2).$$
第二，给定有限集之间的映射 $f: X \rightarrow Y$，有拉回与推前映射
%
$$
 f^!: \on{Fun}(Y) \rightarrow \Fun(X) \quad \text{与} \quad
f_*: \Fun(X) \rightarrow \Fun(Y).
$$
%
具体地说，$f^!$ 把映射 $\phi: Y \rightarrow k$ 送到复合 $\phi \circ f$；而 $f_*$ 沿
$f$ 的纤维对 $\Fun(X)$ 中的函数 $\psi$ 求积分，即对 $Y$ 的任意元素 $y$，令
%
\[
(f_*\psi)(y) := \underset{x \in f^{-1}(y)}{\Sigma} \hspace{.5mm} \psi(x).
\]
%

\ssszh 把这两条性质合起来，$H$ 上的群结构就使 $\Fun(H)$ 具有 $k$-代数结构。确切地说，
若记
%
$
m: H \times H \rightarrow H
$
%
为乘法映射，便得到函数上的卷积映射
%
\begin{equation} \label{e:convfun}
\Fun(H) \otimes \Fun(H) \simeq \Fun(H \times H) \xrightarrow{ m_*}
\Fun(H).
\end{equation}
%

具体地说，若对 $H$ 的元素 $h$ 用 $\delta_h$ 记 $\Fun(H)$ 中相应的 delta 函数，即
%
\[
\delta_h(h') := \begin{cases}   1 & \text{若 } h = h', \\ 0 & \text{否则,}
\end{cases}
\]
%
则 $\delta_h$（$h \in H$）构成 $\Fun(H)$ 的一组基，并按规则
%
\[
\delta_{h_1} \star \delta_{h_2} = \delta_{h_1h_2}
\]
%
相乘，其中 $h_1h_2$ 指 $H$ 中的乘积。对应 $h \mapsto \delta_h$ 给出映射
$H \rightarrow \Fun(H)$，其性质是：任何形如 \eqref{e:homalg} 的映射，都唯一地扩充为
$k$-代数映射
%
\[
\Fun(H) \rightarrow A.
\]
%
特别地，这给出 $H$ 的表示与 $\Fun(H)$ 的左模之间熟悉的等价。

显然，$\Fun(H)$ 由 $H$ 的点的线性组合构成；在 $H$ 于向量空间上的表示中，它按相应的
自同态的线性组合作用。

%=====================================================================
%  ch03 —— 原文 648–862：§2.3 主系列表示
%=====================================================================

\subsection{主系列表示}

\ssszh 本节余下部分回顾 $G$ 的 $\mathbb{F}_q$-点的一个基本表示，即其旗流形上的函数空间。
为此，先作若干准备。

\ssszh  \label{ex:actspace}同前，$H$ 记有限群。若 $H$ 作用在有限集 $X$ 上，就诱导出
$H$ 在 $\Fun(X)$ 上的自然表示。确切地说，给定 $H$ 的元素 $h$ 与 $\Fun(X)$ 中的函数
$\phi$，$h$ 对 $\phi$ 的作用给出元素 $h \cdot \phi$，其定义为
%
\[
(h \cdot \phi)(x) := \phi(h^{-1} x), \quad \text{对 } x \in X.
\]
%
相应的左模结构
%
\begin{equation} \label{e:leftmod}
\Fun(H) \otimes \Fun(X) \rightarrow \Fun(X)
\end{equation}
%
可以这样描述。若把作用记为 $a: H \times X \rightarrow X$，则 \eqref{e:leftmod} 由复合
%
\begin{equation} \label{e:actionfun}
\Fun(H) \otimes \Fun(X) \simeq \Fun(H \times X) \xrightarrow{a_*} \Fun(X).
\end{equation}
%
给出。

\ssszh\label{s:funselfdual} 回忆 $\Fun(X)$ 典范地自对偶。确切地说，记
$\Delta: X \rightarrow X \times X$ 为对角映射、$\pi: X \rightarrow \on{pt}$ 为投影，
则有完美配对
%
\[
\Fun(X) \otimes \Fun(X) \simeq \Fun(X \times X) \xrightarrow{\Delta^!}
\Fun(X) \xrightarrow{\pi_*} \Fun(\on{pt}) \simeq k,
\]
%
具体地说，这不过是「$L^2$」内积
%
\[
f \otimes g \mapsto \underset{x}{\Sigma} \hspace{1mm} f(x) \cdot g(x).
\]
%

\ssszh \label{sss:endos}

特别地，给定有限 $H$-集 $X$ 与 $Y$，有典范等同
%
\[
\begin{gathered}
\on{Hom}_{\on{Rep}(H)^\heartsuit}( \Fun(X), \Fun(Y)) \simeq ( \Fun(X)^* \otimes
\Fun(Y))^{H} \\
\simeq \Fun(X \times Y)^H \simeq \Fun(X \times Y/ H),
\end{gathered}
\]
%
其中上标「$H$」记 $H$-不变子空间，而 $X \times Y / H$ 记 $X \times Y$ 关于 $H$ 的对角
作用的轨道集。

具体地说，给定函数 $K \in \Fun(X \times Y)$，对应的同态 $\Fun(X) \rightarrow \Fun(Y)$
由积分变换
%
\[
   \Fun(X) \xrightarrow{\pi_X^!} \Fun(X \times Y) \xrightarrow{ K \cdot - } \Fun(X \times Y) \xrightarrow{\pi_{Y, *}} \Fun(Y)
\]
%
给出，其中 $\pi_X: X \times Y \rightarrow X$ 与 $\pi_Y: X \times Y \rightarrow Y$ 为投影
映射，$K \cdot -$ 指点乘函数。此外，$K$ 是 $H$-不变的，当且仅当相应的积分变换是
$H$-表示的映射。

\ssszh 我们把上述讨论特殊化到 $G(\mathbb{F}_q)$，其中 $\mathbb{F}_q$ 记有 $q$ 个元素的
有限域。

考虑旗流形 $\sB$，它参数化 $G$ 的 Borel 子群。这是整数环上的射影概形。$G$ 通过共轭
Borel 子群自然地作用在 $\sB$ 上，在 $\mathbb{F}_q$-点上给出作用
%
\[
G(\mathbb{F}_q) \times \sB(\mathbb{F}_q) \rightarrow \sB(\mathbb{F}_q).
\]
%
特别地，如第 \ref{ex:actspace} 小节那样，得到幂幺主系列模
%
\[
\Fun(G(\mathbb{F}_q)) \otimes \Fun(\sB(\mathbb{F}_q)) \rightarrow
\Fun(\sB(\mathbb{F}_q)).
\]
%
此外，它的自同态代数，即 {\em Hecke 代数}，按第 \ref{sss:endos} 小节的方式典范地
等同于
%
\[
\Fun( \sB(\mathbb{F}_q) \times \sB(\mathbb{F}_q) / G(\mathbb{F}_q)).
\]
%
我们要强调：这里考虑的不是相应 $\Fq$-簇上的多项式函数代数，而只是这些簇的有理点所成
有限集上 $k$-值函数的代数。

\ssszh 稍具体些：回忆 $\sB$ 参数化 $G$ 的 Borel 子群。这些子群在 $G$ 中彼此共轭，且自
正规化。因此固定一个 Borel 子群 $B$，就把旗流形等同于齐性空间 $G/B$；在 $\Fq$-点上还有
%
\[
\sB(\mathbb{F}_q) \simeq G(\Fq)/B(\Fq).
\]
%
类似地，$(\sB \times \sB) /G$ 等同于 $B \backslash G / B$，于是有同构
%
\[
\Fun( \sB(\mathbb{F}_q) \times \sB(\mathbb{F}_q) / G(\mathbb{F}_q)) \simeq
\Fun( B(\Fq) \backslash G(\Fq) / B(\Fq)).
\]
%
它把左边的代数结构与右边由卷积给出的代数结构对应起来。

\ssszh 为具体起见，我们较详细地讨论最小的非平凡例子。

\begin{ex} 设 $G = GL(V)$，其中 $V$ 是秩为二的自由阿贝尔群。此时 $\sB$ 是射影直线
$\mathbb{P}(V)$，它参数化 $V$ 的秩一商，并带有 $GL(V)$ 的自然作用。要看这一点，取
$V$ 的一组基 $e_1, e_2$，注意投影 $V \rightarrow V / \langle e_2 \rangle$ 的稳定化子
就是上三角矩阵构成的 Borel 子群
%
\[
B =\begin{pmatrix} * & * \\ & * \end{pmatrix},
\]
%
于是得到所要的与 $G/B$ 的等同。特别地，$\sB(\mathbb{F}_q)$ 有 $q+1$ 个点。

模 $\Fun(\sB(\mathbb{F}_q))$ 分解为两个单模，如下。投影 $\pi: \sB \rightarrow \on{pt}$
给出 $G(\mathbb{F}_q)$-模的（分裂）短正合列
%
\[
0 \rightarrow \on{St} \rightarrow \Fun(\sB(\mathbb{F}_q)) \rightarrow
\Fun(\on{pt}(\mathbb{F}_q)) \simeq k \rightarrow 0,
\]
%
其核 $\on{St}$ 即 Steinberg 模，是单模；理由例如是 Hecke 代数为二维。
\end{ex}

\ssszh 本节内容可概括如下。第一，我们回顾了如何把有限群的群代数构造为卷积代数。
第二，对有限群 $G(\Fq)$，我们回顾了它在旗流形 $\sB(\Fq)$ 上函数中的自然作用。

%=====================================================================
%  ch04 —— 原文 863–1207：§3 Beilinson–Bernstein 局部化（3.1 概览、3.2 局部化定理）
%=====================================================================

\section{Beilinson--Bernstein 局部化与隐藏的对称性}
\label{s:bbloc}

\subsection{概览}
\label{s:overview3}

\ssszh 上一节我们考虑的是旗流形的 $\Fq$-点上函数构成的向量空间
%
\begin{equation} \label{e:funflgs}
\Fun(\sB(\Fq)),
\end{equation}
%
它自然地带有 $G(\Fq)$ 的作用。

本节我们改为关注复旗流形上 D-模的导出范畴
%
\begin{equation} \label{e:dmodflg}
\Dmod(\sB(\C)).
\end{equation}
%
我们先按 Beilinson--Bernstein 回顾这个范畴 \eqref{e:dmodflg} 与李代数表示论之间的基本
联系。在回顾这一对应的若干重要性质时，我们将被引向 $G(\Fq)$ 在 $\Fun(\Fl_G(\Fq))$ 上
作用的类似物。由此我们就会遇到群的范畴作用这一想法。

\subsection{局部化定理}

\ssszh 读者或许已从第 \ref{s:overview3} 小节的讨论中看出，人们把 D-模范畴
\eqref{e:dmodflg} 看作 \eqref{e:funflgs} 的某种范畴化。这个类比在几何表示论中是绝对
根本的。为此，我们在附录 \ref{s:frmfun} 中对它作了相当详细的解释，并介绍了 D-模。
非专家读者不妨先翻到那里，掌握了基础再回来读本节。

\ssszh 本节余下部分在特征零的代数闭域 $k$ 上工作。为记号简便，我们重新定义 $\Fl_G$ 为
相应的基变换。也就是说，若用 $\Fl_{G, \mathbb{Z}}$ 记第 \ref{s:repthryfingroup} 节
讨论过的 $\Spec \Z$ 上的版本，则置
%
\[
 \Fl_G := \Fl_{G, \Z} \underset {\Spec \Z} \times \Spec k.
\]
%
当然，读者可以取 $k$ 为复数，但要注意我们并不需要它们的度量拓扑的任何特殊性质。

\ssszh 旗流形上的 D-模范畴
$$\Dmod(\Fl_G)$$
原来还有一副看似不同的面貌：它可以实现为李代数表示的范畴。下面回顾这一孕育了几何表示论
的非凡对应。

\ssszh 回忆光滑代数簇 $X$ 上的 D-模 $\sM$ 是带平坦联络的拟凝聚层。等价地，$\sM$ 带有
微分算子层 $\sD_X$ 的作用，且该作用延拓正则函数 $\OO_X$ 的作用。特别地，其整体截面带有
整体微分算子的作用，即有一个重言函子
%
\begin{equation} \label{e:glob}
\Gamma(X, -): \Dmod(X) \rightarrow \Gamma(X, \sD_X)\on{-mod.}
\end{equation}
%
Beilinson--Bernstein 的一个著名定理说：当 $X = \sB$ 时，这个函子是等价。

\ssszh 我们来说明这为什么出人意料。

首先，若 $X$ 是仿射的，类似的断言成立，但并不出人意料。事实上，此时它来自拟凝聚层
定义的设定方式：在仿射簇上，函子
%
\begin{equation} \label{e:globqcoh}
\Gamma(X, -): \on{QCoh}(X) \rightarrow \Gamma(X, \OO_X)\mod
\end{equation}
%
是等价。

相反，若 $X$ 不是仿射的，例如射影的，则 \eqref{e:globqcoh} 几乎从不是等价。例如，
对 $\sB$ 而言，其整体（导出）函数只有标量，于是 \eqref{e:globqcoh} 退化为映射
%
\[
\Gamma(\sB, -): \on{QCoh}(\sB) \rightarrow \on{Vect,}
\]
%
其中 $\on{Vect}$ 记向量空间的导出范畴。这远非等价，因为 $\sB$ 不是一个点！正因如此，
在进一步给拟凝聚层配上平坦联络之后，\eqref{e:glob} 对旗流形成立为等价，才是非平凡的事。

我们还要强调：对一般射影簇，\eqref{e:glob} 不是等价。例如，熟悉代数曲线的读者可以
验证：使 \eqref{e:glob} 成为等价的光滑射影曲线只有 $\mathbb{P}^1$，即最小的旗流形。

\ssszh 下面转向整体微分算子
%
\[
\Gamma(\sB, \sD_{\sB})
\]
%
的描述。回忆对任意光滑簇 $X$，$\sD_{X}$ 作为代数层由函数 $\OO_X$ 与向量场 $\sT_X$
生成。我们来看在本情形 $X = \sB$ 中，它们对整体截面有何贡献。

如上所述，旗流形上的整体函数只是标量
%
\[
k \xrightarrow{\sim} \Gamma(\sB, \OO_{\sB}).
\]
%
对向量场，$G$ 在 $\sB$ 上的自然作用给出其李代数 $\fg$ 的无穷小作用，即李代数同态
%
\begin{equation}\label{e:vecfields}
\fg \rightarrow \Gamma(\sB, \sT_{\sB}).
\end{equation}
%
为接近最终答案，我们可以乐观地希望这是同构。

为顾及高阶微分算子，我们同样可以要求由包络代数出发的映射
%
\begin{equation}\label{e:globdiff}
U(\fg) \rightarrow \Gamma(\sB, \sD_{\sB})
\end{equation}
%
是同构。这将是「$\sD_{\sB}$ 的截面局部地由函数与向量场生成」的一个强整体类比。

\ssszh 结果，\eqref{e:vecfields} 与 \eqref{e:globdiff} 都不是同构。必须考虑中心来修正。

先从向量场说起。注意 $G$ 的中心 $Z$ 平凡地作用在 $\sB$ 上。因此需要模去 $\fg$ 的中心
$\fz$。这样做之后，确实得到等价
%
\begin{align*}
\fg/\fz &\xrightarrow{\sim} \Gamma(\sB, \sT \hspace{.2mm}).
\intertext{类似地，过渡到所有微分算子时，必须施加中心商。也就是说，若用 $Z(\fg)$
记 $U(\fg)$ 的中心，它在 $\fg$ 的平凡表示上的作用给出一个特征
$\chi: Z(\fg) \rightarrow k.$ 用 $U_0(\fg)$ 记整个包络代数相应的中心商}
U_0(\fg) &:= U(\fg) \underset{Z(\fg)} \otimes  k_\chi.
\intertext{于是与向量场的情形一样，\eqref{e:globdiff} 穿过所得的商，并得到等价}
U_0(\fg) &\xrightarrow{\sim} \Gamma(\sB, \sD).
\end{align*}

\label{s:casa}

\begin{ex} 我们显式描述 $\fg = \fgl_2$ 时 $U_0(\fg)$ 的样子。可把李代数分解为中心
与无迹矩阵
%
\[
\fgl_2 = \fz \oplus \fsl_2.
\]
%
这诱导包络代数及其中心的分解
%
\[
U(\fgl_2) \simeq \on{Sym}(\fz) \otimes U(\fsl_2) \quad \text{与} \quad
Z(\fgl_2) \simeq \on{Sym}(\fz) \otimes Z(\fsl_2).
\]
%
\label{s:cas}
中心 $Z(\fsl_2)$ 是单变量多项式环，由滤过次数为二的 Casimir 元 $\Omega$ 生成。
具体地说，若取 $\fz$ 的基元 $z$ 与 $\fsl_2$ 的标准生成元 $f, h, e$，则
%
\[
Z(\fgl_2) \simeq k[ z, ef + fe + \frac{1}{2}h^2],
\]
%
其中第二个生成元就是 $\Omega$。容易看出两个生成元在平凡表示上都以零作用。特别地，
$U_0(\fg)$ 是 $U(\fgl_2)$ 模去这两个元素所生成的理想得到的商，即
%
\[
U_0(\fgl_2) = U(\fgl_2)/( \fz, \Omega) \simeq U(\fsl_2)/(\Omega).
\]
%
\end{ex}

\ssszh 用 $\fg\on{-mod}_0$ 记 $U_0(\fg)$-模的导出范畴，即中心特征平凡的 $\fg$ 表示。
总结上面的讨论，Beilinson--Bernstein 定理如下。

\begin{theo} \label{t:bbloc}(Beilinson--Bernstein \cite{beilinson-bernstein81})
整体截面函子给出 $t$-正合的等价
%
\[
\Gamma(\sB, -): \Dmod(\sB) \rightarrow \fg\on{-mod}_0.
\]
%
\end{theo}

这个定理对表示论发展的意义，怎么强调都不过分。这里只作几点说明。

\ssszh 第一，由 Quillen 的一个 Schur 引理版本，每个单 $\fg$-模都有中心特征，即中心
$Z(\fg)$ 以标量作用。因此，研究不可约模时，逐个中心特征来做就足够了。

上述定理把中心特征平凡的那些模，与 $\sB$ 上的单 D-模等同起来。这使人得以把这类不可约模
研究中的问题，翻译成几何或拓扑的问题。后者可以借助代数簇拓扑的丰富知识来处理，
特别是 Hodge 理论及其奇异变体。

该定理最初的动机正是这样一种应用。即在所有 $\fg$-模的范畴中，对 $G$ 的任意 Borel
子群 $B$，可联系上 $(\fg, B)$-模范畴，即其 $B$ 的李代数作用可积的 $\fg$-模。

一个自 Verma 1966 年的论文 \cite{verma66} 起受到大量研究的问题，是确定单
$(\fg, B)$-模的特征。在 Jantzen \cite{jantz79} 等人于若干低秩情形定出之后，Kazhdan 与
Lusztig 于 1979 年对中心特征平凡的单 $(\fg,B)$-模提出了一般猜想 \cite{kl79}。提出猜想
时，他们已知道其公式与 $\sB$ 的 Schubert 子簇密切相关，更确切地说，与其相交上同调
D-模的 $!$-茎密切相关。但与 $\fg$-模的联系尚不清楚。

Beilinson--Bernstein 的工作 \cite{bb82}，以及独立的 Brylinski--Kashiwara
\cite{brylin81}，提供了所需的联系。即局部化定理把 $(\fg,B)\mod_0$ 的单对象与上述相交
上同调层互换，特征公式随之而来。

\begin{re} 上面只讨论了平凡中心特征的情形。Beilinson--Bernstein 实际上对任意正则中心
特征证明了类似的陈述。关于包括广义或奇异中心特征的变体，见 \cite{kashloc}、
\cite{beilinsonginzburg}、\cite{bmr}、\cite{krem}、\cite{ahc}。 \end{re}

\ssszh 接下来我们用三个基本例子说明局部化等价的行为。

\begin{ex} 如第 \ref{s:casa} 小节所述，局部化把微分算子代数与包络代数的中心商互换，即
%
\[
\Gamma(\sB, \sD) \simeq U_0(\fg).
\]
%
\end{ex}

\begin{ex}如我们已讨论的，局部化把结构层与平凡表示互换，即
%
\[
\Gamma(\sB, \OO) \simeq k.
\]
%
读者或许知道，这是 Borel--Weil--Bott 定理的最简单情形；该定理把所有不可约 $G$-模构造为
$\sB$ 上线丛的整体截面。非正式地说，Beilinson--Bernstein 的工作把这一点推广，给出
所有不可约 $\fg$-模的类似构造。
\end{ex}

\begin{ex} 最后，对 $\sB$ 的点 $x$，考虑 $\delta_x$ 的像。回忆 $x$ 的稳定化子是 Borel
子群 $B$。于是，用 $\fb$ 记 $B$ 的李代数，delta D-模被送到相应的诱导模
%
\begin{equation}
\Gamma(\Fl_G, \delta_x) \simeq U(\fg) \underset{U(\fb)} \otimes
\det(\fg/\fb).
\end{equation}
%
我们略述其证明，读者不妨自己完整想一遍。为造出映射
%
\begin{equation} \label{e:hwv}
U(\fg) \underset{U(\fb)} \otimes k_{\chi} \rightarrow \Gamma(\Fl_G,
\delta_x),
\end{equation}
%
注意把右边的生成线与 $\det(T_x \sB) = \det( \fg/\fb)$ 等同起来，与 $\fb$-作用是相容的，
参看注 \ref{r:hwdelta}。要看出所得的映射是同构，只需注意两边都由任一与 $\fb$ 横截的
子代数 $\fn^-$（即某个对立 Borel 的幂幺根）的作用，从各自的行列式线自由生成。
\end{ex}

%=====================================================================
%  ch05 —— 原文 1208–1793：§3.3 Equivariance and Harish-Chandra modules
%=====================================================================

\subsection{等变性与 Harish-Chandra 模}

\ssszh 专家们立刻看出，局部化定理所承载的结构远不止范畴的等价。对具体的表示论应用
而言，主要的相容性如下。固定代数子群 $H \subset G$，为简单起见设它连通。

\ssszh 先从李理论一侧开始。与 $H$ 相伴的是 Harish-Chandra 模构成的全子范畴
%
\[
(\fg, H)\mod^\heartsuit \subset \fg\mod^\heartsuit.
\]
%
具体地说，一个 $\fg$-模是 Harish-Chandra 模，是指李代数 $\fh \subset \fg$ 的作用可积分
为 $H$ 的作用。由于我们假设 $H$ 连通，这是一条性质，而不是额外的结构。

限制到中心特征平凡的模，得到类似的全子范畴
%
\[
(\fg, H)\mod_0^\heartsuit \subset \fg\mod_0^\heartsuit.
\]
%

\ssszh 下面举几个简短的例子。

\begin{ex}
若 $\fg$ 是半单的，则中心商 $U_0(\fg)$ 只是平凡子群的 Harish-Chandra 模。
\end{ex}

\begin{ex} 平凡模 $k$ 是 $G$ 自身的 Harish-Chandra 模，从而也是任意子群的
Harish-Chandra 模。
\end{ex}

\begin{ex} Verma 模 $U(\fg) \underset{U(\fb)} \otimes \det(\fg/\fb)$ 是 $(\fg, B)$-模。
更一般地，对 $H$ 的任意表示 $W$，诱导模
%
\[
U(\fg) \underset{U(\fh)} \otimes W
\]
%
是 $(\fg, H)$-模，而每个 $(\fg, H)$-模都是这种模的商。
\end{ex}

\ssszh \label{s:eqdmod}下面转向几何一侧。与 $H$ 相伴的是 $H$-等变 D-模构成的全子范畴
%
\[
\Dmod(\sB)^{H, \heartsuit} \subset \Dmod(\sB)^\heartsuit.
\]
%

我们在任意带 $H$ 作用的代数簇 $X$ 的一般情形下回顾它的定义。

\begin{re} 对我们来说，代数簇按定义是 $k$ 上有限型的、拟紧、拟分离、既约的概形。
此外，本书对代数簇陈述的一切，对 $k$ 上有限型的拟紧拟分离概形同样成立，也就是说，
既约性并不起任何作用。读者若跳过这些技术细节，只考虑例如拟射影簇，也不会有多少损失。
\end{re}

初看之下，一个 D-模 $M$ 是 $H$-等变的，大致是说它沿 $X$ 中的 $H$-轨道为常值。
要得到真正的定义，注意我们因此应当要求：对 $H$ 的任意点 $h$ 与 $X$ 的任意点 $x$，
存在茎的同构
%
\begin{equation}\label{e:stalks}
i_x^! M \simeq i_{hx}^! M.
\end{equation}
%
当然，由于我们的空间并不离散，需要把这样的陈述在 $h$ 与 $x$ 变动时统一成一个族，
并与平行移动相容。为此，考虑作用映射与投影映射
%
\begin{equation}
\alpha, \pi: H \times X \rightarrow X,\quad \quad \alpha(h, x) = hx,
\quad \pi(h, x) = x.
\end{equation}
%

有了这些，$\Dmod(X)^\heartsuit$ 的对象 $M$ 是 $H$-等变的，当且仅当存在同构
%
\begin{equation} \label{e:wholething}
\alpha^! M \simeq \pi^! M.
\end{equation}
%
由于我们只要求这样的同构存在，这显然是 D-模的一条性质，而非额外的结构。

与函数的情形作比较是有启发的。即，固定一个作用在有限集上的有限群，不妨暗示性地
记它们为 $H(\Fq)$ 与 $X(\Fq)$。注意，$\Fun(X(\Fq))$ 中的函数 $m$ 是 $H(\Fq)$-不变的，
即沿轨道为常值，当且仅当它满足等式
%
\[
\alpha^! m = \pi^! m,
\]
%
其中 $\alpha, \pi: H(\Fq) \times X(\Fq) \rightarrow X(\Fq)$ 如上定义，而函数的
$!$-拉回如第 \ref{ss:opsfinitesets} 小节所述。

\ssszh 下面是一对基本而重要的等变 D-模的例子。

\begin{ex} 设 $X$ 由单个 $H$-轨道构成，即选定基点后形如 $H/K$，其中 $K$ 是 $H$ 的
某个子群。此时我们断言，$X$ 上的 $H$-等变 D-模与 $K$ 的连通分支群的表示之间，有一个
典则等价，即
%
\[
\Dmod(X)^{H, \heartsuit} \simeq \on{Rep}(\pi_0(K))^{\heartsuit}.
\]
%
事实上，容易看出 $X$ 上任一等变 D-模必是局部系统，而条件 \eqref{e:wholething}
等价于说沿 $H \rightarrow H/K$ 的拉回可平凡化。另一方面，纤维化
$K \rightarrow H \rightarrow H/K$ 所联系的同伦群长正合列在低次给出正合列
%
\[
\pi_1(H) \rightarrow \pi_1(H/K) \rightarrow \pi_0(K) \rightarrow 1,
\]
%
于是断言得证。 \label{e:oneorbit}
\end{ex}

\begin{ex} 设 $X$ 由有限多条 $H$-轨道构成。此时 $X$ 上的 $H$-等变 D-模范畴可由
各条轨道通过重粘（recollement）「拼接」而成。

其细节对我们并不那么重要。不过要提到下面这条基本推论：$X$ 上不可约的 $H$-等变
D-模，与各条轨道上不可约的 $H$-等变局部系统，处于典则的双射。具体地说，
$\Dmod(X)^{H, \heartsuit}$ 的每个不可约对象都支撑在单条轨道的闭包上，它在该轨道上的
限制是不可约的 $H$-等变局部系统。这样，由上一个例子，单对象的分类就归结为有限群论的
一道习题。

\label{e:bigbrain}
\end{ex}

\ssszh 两侧都已介绍，现在可以陈述所许诺的相容性了：局部化把 Harish-Chandra 模与
等变 D-模互换。

\begin{theo}[Beilinson--Bernstein, \cite{beilinson-bernstein81}]\label{t:heqv}
有交换图
%
\[
\xymatrix{\Dmod(\sB)^{H, \heartsuit} \ar[r]^\sim \ar[d] & (\fg,
H)\mod_0^\heartsuit \ar[d]  \\  \Dmod(\sB)^\heartsuit \ar[r]^\sim &
\fg\mod_0^\heartsuit,}
\]
%
其中竖直箭头是重言的包含函子，底部的水平箭头是局部化。
\end{theo}

\ssszh 事实上，定理 \ref{t:heqv} 在作了必要的调整（mutatis mutandis）之后，对任意的
$H$（不必连通）成立，对适当定义的等变 D-模与 Harish-Chandra 模的导出范畴也成立。
后一断言并不能由前一断言推出，因为这些导出范畴通常并不是其心（heart）的导出范畴，
而是稍为精细的对象。

\begin{re} 建议读者把上述关于导出范畴的微妙之处，看成如下形式的陈述：「要得到正确的答案，就必须稍微聪明一点。」然而，不难看出为什么朴素地用等变层的阿贝尔范畴的导出范畴
去处理会得到错误的答案。这里的定向原则是：$X$ 上的 $H$-等变 D-模，应当就是商叠
$X/H$ 上的 D-模，参看下面的注 \ref{r:introstacks}，因此应当与它的拓扑保持预期的关系。
在 $X = \on{pt}$ 这个最简单的情形，我们尤其应当希望
%
\[
\on{Hom}_{\Dmod(\on{pt})^H}(k, k) \simeq \on{C}^*(\on{pt}/H, k),
\]
%
即：常值层的导出自同态应当是分类空间的上同调。但既然在阿贝尔范畴的层次上
%
\[
\Dmod(\on{pt})^{H, \heartsuit} \simeq \on{Rep}(\pi_0(H)),
\]
%
我们单从阿贝尔范畴本身根本得不到任何高阶 Ext。这一修正归功于 Bernstein--Lunts，
我们将在第 \ref{s:catrepgrp} 节遇到它。
\end{re}

\begin{re} \label{r:introstacks}对不熟悉叠的读者，下面这个粗略的描述已足以应付我们的
需要。代数闭域 $k$ 上的 Artin 叠 $Y$，大致说来就像代数簇，但多了一条性质：$Y$ 的每个
$k$-点 $y$ 都带有一群「内部对称性。」

对我们最相关的一类 Artin 叠来自如下构造。给定代数群 $H$ 作用在代数簇 $X$ 上，可以
构造商叠 $X/H$。$X/H$ 的点由 $X$ 上的 $H$-轨道给出；而对 $X$ 的 $k$-点 $x$，
$X/H$ 中对应点的内部对称群恰是 $x$ 在 $H$ 中的稳定化子 $H_x$。也就是说，这个叠的
$k$-点并不是代数簇 $\pt$，即 $\Spec k$，而是它的商 $\pt/H_x$。一般地，Artin 叠的
$k$-点 $y$ 是某个代数群 $H_y$ 的分类空间 $\pt/H_y$，这个 $H_y$ 就是上面提到的内部
对称群。

作为简单的例子，考虑 $X = \mathbb{A}^1$ 带上由对合 $z \mapsto -z$ 诱导的
$H = \mathbb{Z}/2$ 的作用。若挖掉原点，这就限制为 $H$ 在 $\mathbb{G}_m$ 上的作用。
此时每条 $H$-轨道都是自由的，特别地，商中的点应当像普通代数簇那样具有平凡的内部
对称性。事实上，商 $\mathbb{G}_m/H$ 就是代数簇 $\mathbb{G}_m$，其中的商映射
%
\[
\mathbb{G}_m \rightarrow \mathbb{G}_m / H
\]
%
等同于平方映射 $\mathbb{G}_m \rightarrow \mathbb{G}_m$，$w \mapsto w^2$。
然而，完整的商 $\mathbb{A}^1/H$ 不是代数簇，因为对应原点的那个点已「轨道化」成了一份
$\pt/H$。直观地说，对一般的商叠，内部对称群记录了在从 $X$ 构造轨道空间 $X/H$ 时
发生的折叠。

最后要提到，与代数簇的情形相比，叠上的层的基本新特征是其纤维带有内部对称群的作用。
即，对代数群 $H$，有
%
\[
\QCoh(\pt/H) \simeq \on{Rep}(H) \quad \text{与} \quad \Dmod(\pt/H) \simeq
\on{C}_*(H)\mod,
\]
%
其中 $\on{C}_*(H)$ 记 $H$ 的 de Rham 上同调的对偶，并看作卷积下的 dg 代数。

因此，对 Artin 叠 $Y$ 上的凝聚层，它在 $k$-点 $y$ 处的 $*$-纤维带有 $H_y$ 的作用。
类似地，对 $Y$ 上的 D-模，它在 $y$ 处的 $!$-纤维带有 $\on{C}_*(H_y)$ 的作用，可以把
它想成 $H_y$ 的局部常值作用。特别地，该 $!$-纤维的上同调群带有连通分支群
$\pi_0(H_y)$ 的作用。
\end{re}

\ssszh 下面给出定理 \ref{t:heqv} 如何应用的一个基本例子。

由 \ref{t:heqv} 可知，$(\fg, H)\mod_0$ 中只有有限多个不可约对象，当且仅当 $\sB$ 上
只有有限多条 $H$-轨道。而且此时可以像例 \ref{e:bigbrain} 那样显式列举单对象。

作为具体例子，记 $T \subset GL_2$ 为极大环面，我们来分类
%
\[
(\fgl_2, T)\mod_0^\heartsuit
\]
%
的单对象。为此注意 $T$ 在 $\mathbb{P}^1$ 上有三条轨道——北极 $B_0$、南极
$B_\infty$，以及其余一切。由于 $T$ 在这三条轨道上的稳定化子都连通，可知共有三个
不可约对象。具体地说，它们是最高权为 $-2$ 的单 Verma 模、最低权为 $2$ 的单 Verma 模，
以及平凡模，即
%
\[
      U(\fgl_2) \underset{U(\fb_0)} \otimes \det(\fgl_2/\fb_0), \quad
      U(\fgl_2) \underset{U(\fb_\infty)} \otimes
      \det(\fgl_2/\fb_\infty), \quad \text{与} \quad k.
\]
%
还要注意，由类似的推理，当 $n > 2$ 时，$(\fgl_n, T)\mod_0^\heartsuit$ 中有无穷多个
单对象。

这段讨论相当软性，但我们要强调：更精细的问题，例如确定单模的特征，可以从几何入手，
而且通常翻译为关于轨道闭包拓扑的问题。

\ssszh 最后，就定理 \ref{t:heqv} 为何成立作两点定向性的说明。在下面的普遍性中讨论
会更清楚。\label{s:kandx}

设代数群 $K$ 作用在光滑代数簇 $X$ 上。如前面那样，若用 $\fk$ 记 $K$ 的李代数，
则有整体截面函子
%
\[
   \Gamma(X, -): \Dmod(X)^\heartsuit \rightarrow \fk\mod^\heartsuit.
\]
%
它有左伴随，由局部化给出
%
\[
   \on{Loc}: \fk\mod^\heartsuit \rightarrow \Dmod(X)^\heartsuit, \quad
   \quad M \mapsto \sD_X \underset{U(\fk) } \otimes M.
\]
%
我们断言，对任意代数子群 $H \subset K$，它限制为一个伴随对
%
\begin{equation} \label{e:locadj}
  \on{Loc}:(\fk, H)\mod^\heartsuit \leftrightarrows
  \Dmod(X)^{H, \heartsuit}: \Gamma(X, -).
\end{equation}
%
特别地，取 $K = G$、$X = \sB$，就重新得到定理 \ref{t:heqv}。

\ssszh 要看出 \eqref{e:locadj} 为何成立，需要先重新审视等变 D-模的定义。我们沿用
第 \ref{s:eqdmod} 小节的记号。对一般的群 $H$（不必连通），等变 D-模定义中的部分数据
是一个指定的同构
%
\[
\tau:  \alpha^! M \simeq \pi^! M,
\]
%
它满足单位性，并在 $H \times H \times X$ 上满足上圈条件（即结合性条件）。在茎的
层次上，上圈条件要求：把 $i^!_{h_1 h_2 x} M \simeq i_x^! M$ 认同为 $i_x^! M$ 的
两种方式——或对配对 $(h_1h_2, x)$ 用一次 $\tau$，或先对 $(h_2, x)$、再对
$(h_1, h_2x)$ 用两次 $\tau$——彼此一致。有一个很好的事实：若 $H$ 连通，则这样的
数据只要存在就唯一，而其存在性等价于先验地更弱的陈述 \eqref{e:wholething}。

\begin{re}
对不连通的群，等变性确实是为 D-模附加的进一步结构。例如可以验证：若 $X = \on{pt}$，
则 $\Dmod(\on{pt})^\heartsuit \simeq \Vect^\heartsuit$ 的对象 $W$ 上的
$H$-等变数据，等同于连通分支群 $\pi_0(H)$ 在 $W$ 上的表示。

读者或许从例 \ref{e:oneorbit} 中认出了这类表示。这不是巧合，事实上我们有
%
\[ \Dmod(H/K)^{H, \heartsuit} \simeq \Dmod(H)^{H \times K, \heartsuit}
\simeq \Dmod(\on{pt})^{K, \heartsuit}.  \]
%
用叠上的 D-模来说，这来自恒等式 $H \backslash H/K \simeq \on{pt}/K$。
\end{re}

\label{s:rigideq}

\ssszh \label{s:rigimon}第 \ref{s:rigideq} 小节中这种刚性化的等变性定义，或许会让
读者想起代数 $H$-模（即 $\sO_H$-余模）的定义。后者是一个向量空间 $W$，配有满足
单位性与上圈条件的余作用映射
%
\[
   W \rightarrow \sO_H \otimes W.
\]
%

\begin{re}若读者此前未见过这个定义，在目前的情形下它归结为下面这个或许更熟悉的描述：
$W$ 带有抽象群 $H(k)$ 的作用，是有限维表示 $M_\alpha$ 的并，并且关于这种子表示
$M_\alpha$ 的任意一组基，$H(k)$ 作用的矩阵系数都是 $H$ 上的正则函数。
\end{re}

把 $W$ 看作 $\on{pt}$ 上的拟凝聚层，这可以改述为 $\on{pt}$ 上一个 $H$-等变拟凝聚层
的数据，即沿 $\on{pt} \rightarrow \on{pt}/H$ 的下降数据。

有了这些，我们可以把 $H$-簇 $X$ 上 D-模的 $H$-等变数据重新表述如下。先设 $X$ 仿射，
并对 D-模 $M$ 用同一个字母记相应的 $D_X$-模，即其整体截面。

于是，$M$ 上的 $H$-等变数据，等同于在其底向量空间上给定 $H$ 的代数作用，满足
(i) 它与 $H$ 在 $D_X$ 上的作用相容（即从 $D_X$-模升级为 $H \ltimes D_X$-模），
(ii) 由此诱导的 $\fh$ 在 $M$ 上的无穷小作用，与用向量场的作用一致，即与下述复合一致
%
\begin{equation} \label{e:hch}
\fh \rightarrow T_X \rightarrow D_X.
\end{equation}
%
也就是说，人们把 $M$ 上由联络给出的 $\fh$ 作用，积分成群 $H$ 的作用。

对一般的 $X$（不必仿射），可以给出类似的定义。具体地说，D-模的等变数据是
(i) 其底拟凝聚层的等变数据，它与代数层 $\sD_X$ 上典则的等变数据相容，并满足上面
条件 (ii) 的类似物。

\ssszh 由这个表述可以看出，整体截面限制为一个函子
%
\[
         \Dmod(X)^{H, \heartsuit} \rightarrow (\fk, H)\mod^\heartsuit.
\]
%

我们再来勾勒局部化为何限制为一个函子
%
\[
   (\fk, H)\mod^\heartsuit \rightarrow \Dmod(X)^{H, \heartsuit}.
\]
%
出发点是：$\sD_X$ 自然地是 $H$-等变的拟凝聚代数层——这只是在说，函数与向量场可以
沿 $H$ 作用移动。然而，$\sD_X$ 并\emph{不}是 $H$-等变的 D-模。问题出在条件 (ii)，
即 $\fh$ 的两个无穷小作用必须一致。

要看出这一点，注意由 $\fh$ 在 $\sD_X$ 上的左作用所诱导的 $\fh$ 作用，形如
$h \otimes P \mapsto h \cdot P$，其中 $h$ 是 $\fh$ 的元素，$P$ 是 $\sD_X$ 的局部
截面。另一方面，由 $H$-等变性得到的 $\sD_X$ 上的无穷小 $\fh$ 作用必须由该代数的
导子给出，因此不能与左作用一致。于是它只能是伴随作用，即形如
%
\[
     h \otimes P \mapsto h\cdot P - P\cdot h.
\]
%

既然问题出在右乘的出现，读者就不会惊讶：用右乘与一个 $H$-可积模作张量应当有帮助。
确实，给定 $(\fk, H)$-模 $M$，容易看出
%
\begin{equation*}
\on{Loc}(M) = \sD_X \underset{U(\fk)} \otimes M
\end{equation*}
%
的底拟凝聚层，从 $H$ 在其张量因子上的作用继承了一份 $H$-等变数据。至于条件 (ii)，
注意这个作用的微分形如
%
\[
          h \otimes (P \otimes m) \mapsto h \cdot P \otimes m - P \cdot h
          \otimes m + P \otimes h \cdot m,
\]
%
其中 $h$ 与 $P$ 如前，$m$ 是 $M$ 的元素。但现在第二、三项相消，因为我们是在
$U(\fk)$ 上作张量，于是只剩下所希望的第一项。

%=====================================================================
%  ch06 —— 原文 1794–2231：§3.4 Hidden Symmetries
%=====================================================================

\subsection{隐藏的对称性}

\ssszh 上一小节我们看到，局部化满足一条重要的相容性——它把等变 D-模与
Harish-Chandra 模互换。把这一断言展开之后，它归结为关于函子形式的一个相当具体而
初等的观察。

然而必须认识到，这样的陈述，尤其是在它刚被提出的时候，并不容易嵌入表示论中已知的
一般形式体系。

现在我们要说明，这条相容性其实是局部化两侧都存在的一个更丰富结构的影子。在几何
一侧，进一步的对称性即使不算显然，也并不完全出人意料。然而在表示论一侧，它就远不那么
显然了。可以恰当地把它看作李代数表示范畴的一个隐藏对称性。不过正如我们将要说明的，
它的各个片段曾在个案基础上被看到并被卓有成效地使用，但完整的结构在几十年间都没有浮现。

\ssszh 我们的出发点将是下面的类比：
%
\[
   \Dmod(\sB) \quad \text{与} \quad \Fun(\sB(\Fq)).
\]
%
回忆后者带有 $G(\Fq)$ 的作用，这是 $G(\Fq)$ 在 $\sB(\Fq)$ 本身上作用的线性化。

这提示了如下的想法——我们应当把 $\Dmod(\sB)$ 看成群 $G$ 的某种表示，同样是把 $G$ 在
$\sB$ 上的作用线性化。关键在于，局部化等价
%
\[
    \Dmod(\sB) \simeq \fg\mod_0
\]
%
应当是 $G$-模的等价，特别地，李代数表示的范畴本身应当是一个 $G$-模。要把这一点说
精确，需要对代数群在范畴（而非向量空间）上的作用有一个好的概念。

\ssszh 我们先在第 \ref{s:kandx} 小节那样的较一般设定（$K$ 与 $X$）中，对这个想法作
一个合理性检查。

回忆局部化与整体截面函子
%
\[
     \on{Loc}: \fk\mod \rightleftarrows \Dmod(X): \Gamma(X, -).
\]
%
无论 $K$ 在范畴 $\sC$ 上的作用该是什么，至少对 $K$ 的每个元素 $g$，我们应当有范畴的
自同构
%
$$
    g: \sC \xrightarrow{\sim}\sC,
$$
%
而且它们在复合下应当适当相容。现在我们来为 $\fk\mod$ 与 $\Dmod(X)$ 找出这样的
自同构，使局部化把两者交织起来。

我们从几何一侧开始。由于 $X$ 是 $K$-簇，用 $g$ 作用给出簇本身的自同构
$$g: X \xrightarrow{\sim} X$$，从而给出其 D-模范畴的自同构
%
\[
      g_*: \Dmod(X) \xrightarrow{\sim} \Dmod(X).
\]
%
例如，对 $X$ 的任意点 $x$，它把 $\delta_x$ 送到 $\delta_{gx}$。注意对 $K$ 中任意
$g$ 与 $g'$，有典范同构
%
\[
     g_* \circ g'_* \simeq (gg')_*,
\]
%
即群的乘法与自函子的复合之间具有相容性。

为看出 $\fk\mod$ 上相应的算子应当是什么，我们来看它对整体截面的作用。注意对 $X$ 上
任一 D-模 $M$，
%
\[
   \Gamma(X, M) \quad \text{与} \quad \Gamma(X, g_* M)
\]
%
的底向量空间有典范的等同，但整体微分算子的作用被 $g$ 扭了一下。

有鉴于此，记
%
\[
       \on{Ad}_{g, *}: \fk\mod \xrightarrow{\sim} \fk\mod
\]
%
为下述自同构：它把作用为 $\fk \otimes N \rightarrow N$、$X \otimes n \mapsto X
\cdot n$ 的 $\fk$-模 $N$，送到底向量空间相同、作用为 $X \otimes n \mapsto
\on{Ad}_{g^{-1}}(X) \cdot n$ 的 $\fk$-模。这里同样有相容性
%
\[
   \on{Ad}_{g, *} \circ \on{Ad}_{g', *} \simeq \on{Ad}_{gg', *}.
\]
%

\ssszh 由上述讨论，对 $K$ 中每个 $g$，我们得到典范地交换的图
%
\[
\xymatrix{\Dmod(X) \ar[d]^{g_*} \ar[rr]^{\Gamma(X, -)} && \fk\mod
\ar[d]^{\on{Ad}_{g, *}} \\ \Dmod(X) \ar[rr]^{\Gamma(X, -)} & &
\fk\mod. }
\]
%
也就是说，$\Gamma(X, -)$ 把目前为止在两个范畴上构造的 $K$ 作用交织起来。读者不妨
验证 $\on{Loc}$ 也有同样的性质。

\ssszh 现在到了重要的一点。我们至今为止期望一个代数群的范畴表示所具有的结构量，
显然还不够。例如，我们只用到它的闭点，并把它们看作抽象群。

为继续推进，回忆 $G(\Fq)$ 在 $\Fun(\sB(\Fq))$ 上的作用可延拓为群代数
$\Fun(G(\Fq))$ 的作用。按照函数-层对应，我们发现导出范畴 $\Dmod(K)$ 在卷积下是
幺半的。具体地说，给定 $\Dmod(K)$ 中两个层 $\sF$ 与 $\sG$，其卷积为
%
\begin{equation} \label{e:convalgdmod}
  \sF \star \sG := m_*\big( p_1^! \sF \overset{!}{\otimes} p_2^! \sG \big),
\end{equation}
%
其中用 $m, p_1, p_2: K \times K \rightarrow K$ 分别记乘法映射与到第一、第二因子的
投影，并对态射使用显然的函子性。

我们强调，在函数-层对应下，公式 \eqref{e:convalgdmod} 与群代数中通常的乘法相符。
与此相关，幺半单位是支撑在群单位元处的 delta 函数 $\delta_e$。有点启发式地说，正如
通常的群代数那样，可以把一般的 D-模想成 delta D-模的「连续线性组合」。

类似地，$\Dmod(X)$ 自然地成为 $\Dmod(K)$ 的模范畴。具体地说，给定
$\Dmod(K)$ 中的层 $\sF$ 与 $\Dmod(X)$ 中的 $\sM$，其卷积为
%
\begin{equation} \label{e:conv}
     \sF \star \sM := a_*\big( p_1^! \sF \overset{!}{\otimes} p_2^! \sM\big),
\end{equation}
%
其中 $a, p_1, p_2: K \times X \rightarrow X$ 记作用映射与投影映射。

按构造，与支撑在 $K$ 的点 $g$ 处的 delta 函数作卷积，重现已讨论过的对称性，即有
函子的典范同构
%
\begin{equation}
\label{e:action}
 \delta_g \star (-) \simeq g_*: \Dmod(X) \rightarrow \Dmod(X).
\end{equation}
%
非正式地说，我们因此把作用从单个 delta 函数延拓到了它们的全部线性组合，即
$\Dmod(K)$。

下一个任务是为李代数表示找出 $\Dmod(K)$ 的相应作用。在此之前，我们想对 D-模一侧
再作几点评论。

\begin{re} 注意，关于自同构 $g_*$ 与 $\on{Ad}_{g, *}$ 的讨论在阿贝尔范畴或导出范畴
中都同样可行；但要使这个幺半结构按所希望的方式运作，即甚至要有一个像样的 D-模前推，
就必须在导出范畴中工作。目前三角范畴已经够用，但下一节我们将被迫要求更多。

还要注意，要在卷积下封闭，我们必须立刻离开有限生成 D-模的世界。原型例子是 $\sD_K$
与自身的卷积，可以证明
%
\[
    \sD_K \star \sD_K \simeq \sD_K \underset{\underline{k}}{\otimes} \OO_K.
\]
%
\end{re}

\begin{re} 把卷积与作用的公式改写为
%
\begin{align*}
&\Dmod(K) \otimes \Dmod(K) \simeq \Dmod(K \times K) \xrightarrow{m_*}
\Dmod(K) \quad \text{与} \\ &\Dmod(K) \otimes \Dmod(X) \simeq
\Dmod(K \times X) \xrightarrow{a_*} \Dmod(X),
\end{align*}
%
会更有启发，正如函数情形那样，参看公式 \eqref{e:convfun} 与 \eqref{e:actionfun}。
不过为此需要导出范畴的张量积，而这只能在 dg 范畴的层次上进行，见第
\ref{s:catrepgrp} 节。
\end{re}

\ssszh\label{ss:weakinv} 最后，转向 $K$ 在 $\fk\mod$ 上的范畴作用。首先注意，
目前这个作用的构造不像 $\Dmod(X)$ 的情形那样直接。确实，那里我们只需要 D-模作用在
D-模上，这可以用几何函子性；这里我们却需要 D-模作用在李代数表示上，两者表面上相距
更远。

不过，我们可以把范畴 $\fk\mod$ 改写为 D-模的范畴来推进。回忆若 $K$ 作用在簇 $X$
上，则微分算子层 $\sD_X$ 不是等变 D-模。然而它是\emph{弱等变} D-模的原型例子——
这种对象的阿贝尔范畴记作
%
\[
     \Dmod(X)^{K, w, \heartsuit}
\]
%
其定义与第 \ref{s:rigimon} 小节完全相同，只是不加关键的条件 (ii)。

若取 $X = K$，看作右乘下的 $K$-簇，我们断言有典范等价
%
\begin{equation} \label{e:weakinvdmod}
    \Dmod(K)^{K, w, \heartsuit} \simeq \fk\mod^\heartsuit.
\end{equation}
%
为方便读者，我们勾勒证明。首先验证左端有投射生成元 $\sD_K$。这来自它余表示
$K$-不变的整体截面，即 $\Gamma(K, -)^K$，再结合 $K$ 上 $K$-等变拟凝聚层与
$\on{Vect}$ 之间的重言同构。其次，回忆右不变微分算子与包络代数的等同，即
%
\[
     \Gamma(K, \sD_K)^K \simeq U(\fk).
\]
%
总之，我们找到了一个自同态为 $U(\fk)$ 的投射生成元，由 Morita 理论即得结论。

公式 \eqref{e:weakinvdmod} 的等价可延拓为导出范畴的等价
%
\[
    \Dmod(K)^{K, w} \simeq \fk\mod,
\]
%
而前者通过左卷积带有 $\Dmod(K)$ 的作用，所需作用的构造至此完成。

\ssszh 有了这些，局部化与整体截面函子所希望的升级如下。

\begin{theo}[Beilinson--Drinfeld, \cite{bdh}]\label{t:loc+equiv} 设 $X$ 与 $K$ 如
第 \ref{s:kandx} 小节。有 $\Dmod(K)$-等变函子的典范伴随对
%
\[
    \on{Loc}: \fk\mod \rightleftarrows \Dmod(X): \Gamma(X, -).
\]
%
\end{theo}

特别地，定理 \ref{t:bbloc} 中所述的 Beilinson--Bernstein 局部化，是
$\Dmod(G)$ 表示的等价。\footnote{细心的读者会注意到，该局部化定理涉及的不是整个
$\fg\mod$，而是 $\fg\mod_0$。$\Dmod(G)$ 在 $\fg\mod_0$ 上的作用，形式地可由
$Z(\fg)$ 与 $\Dmod(G)$ 在 $\fg\mod$ 上作用的适当交换性推出，也可像第
\ref{ss:weakinv} 小节那样手工构造。}定理 \ref{t:loc+equiv} 的解释，包括对
$\Dmod(K)$-等变性的讨论，我们推迟到第 \ref{s:catrepgrp} 节。

\ssszh 下面指出这一理论的动机性应用。为定向，考虑有限群的子群，不妨暗示性地记作
$H(\Fq) \subset K(\Fq)$。回忆对 $K(\Fq)$ 的任一表示 $\pi$，其
$H(\Fq)$-不变量 $\pi^{H(\Fq)}$ 带有卷积代数
$$\Fun(H(\Fq) \backslash K(\Fq) /H(\Fq))$$ 的作用。

类似地，给定 $K$ 的代数子群 $H$，对 $K$ 的任一范畴表示 $\sC$ 可以取其 $H$-不变
对象 $\sC^H$，它们将带有卷积范畴
%
\[
   \Dmod(H \backslash K  / H)
\]
%
的作用。在目前的例子中，我们得到定理 \ref{t:heqv} 及其导出变体的如下升级。

\begin{theo}[Beilinson--Drinfeld \cite{bdh}]\label{t:lochchmod+eq}存在
$\Dmod(H \backslash K / H)$-等变函子的典范伴随对
%
\[
\on{Loc}: \fk\mod^H \rightleftarrows \Dmod(X)^H: \Gamma(X, -).
\]
%
\end{theo}

一旦形式体系就位，这可以由定理 \ref{t:loc+equiv} 施加 $H$-不变量得到。特别地，
范畴不变量的一般概念（将在第 \ref{s:catrepgrp} 节讨论）在李理论一侧重现
Harish-Chandra 模的范畴
%
\[
   (\fk, H)\mod \simeq \fk\mod^H
\]
%
而在几何一侧重现等变 D-模的范畴。

\ssszh\label{ss:hitchinhecke} 定理 \ref{t:lochchmod+eq} 是 Beilinson--Drinfeld
在他们关于 Hitchin 系统量子化的奠基性工作 \cite{bdh} 中证明的。\footnote{严格说来，
他们需要并证明的是某个无穷维的类似物。}在这个动机问题中，与
$\Dmod(H \backslash K / H)$ 作用的相容性至关重要。即，所希望的结果是 Hecke 特征层的
存在，那是一类关于这个作用具有等变性质的 D-模。它们经由局部化从李代数表示构造
出来，而那些表示被直接证明具有所需性质。

\ssszh 下面是一个不那么精巧的例子。考虑最高权表示的范畴
%
\[
    \fg\mod^B.
\]
%
由上述形式体系，它重言地带有 Hecke 范畴
%
\[
   \Dmod(B \backslash G / B).
\]
%
的作用。我们要强调，在其历史发展中，最高权模与旗簇上 D-模之间的联系是通过局部化
定理才被看到的，特别地，它高度不显然。然而范畴作用的形式体系使它们的关系一目了然。
而且，这迅速而概念性地重现了该范畴的若干重要自函子。

例如，对有限 Weyl 群的每个元素 $w$，可附上某些自函子
%
\[
  j_{w, !} \quad \text{与} \quad j_{w, *}: \fg\mod^B \rightarrow \fg\mod^B,
\]
%
分别称为 Enright 完备化函子与 Arkhipov 扭化函子 \cite{enright79}、\cite{bb82}、
\cite{ark}。几何上，它们由与
%
\[
     BwB \subset G
\]
%
上常值 D-模的 $!$-和 $*$-扩张作卷积给出，其许多基本性质都有简单的几何证明。

作为相关的例子，在 Schubert 簇相交上同调的若干非平凡事实之下，
Kazhdan--Lusztig 猜想可由下述断言推出：
%
\begin{equation} \label{e:brungus}
  \fg\mod_0^B
\end{equation}
%
与 $\Dmod(B \backslash G / B)$ 典范等价。这一点最初是通过建立局部化定理证明的。不过，
利用 Hecke 范畴在 \eqref{e:brungus} 上的作用，也可以直接验证：与反支配 Verma 模作
卷积给出等价
%
\[
  \Dmod(B \backslash G / B) \simeq \fg\mod_0^B,
\]
%
由此形式地推出 $\Gamma: \Dmod(\Fl_G) \rightarrow \fg\mod_0$ 的完全忠实性，参看
\cite{ahc}。

\section{群的范畴表示}
\label{s:catrepgrp}

\subsection{概览}

\ssszh 在第 \ref{s:bbloc} 节中，我们追随 Beilinson--Bernstein 局部化的对称性，得到了群的范畴表示这一概念。本节要给这套形式体系一个更细致的介绍。

这里的基本想法与结果归功于 Beilinson--Drinfeld 与 Frenkel--Gaitsgory。讨论全体范畴表示所构成的范畴、以及它们之间缠绕函子的形式体系，则主要归功于 Gaitsgory。

这套形式体系需要现代同调代数的一些较精细的要点才能起步，我们就从这里讲起。

\subsection{从三角范畴到 dg 范畴}

\ssszh 在第 \ref{s:bbloc} 节中，我们把 $\Dmod(K)$ 当作幺半三角范畴，它的模则是三角范畴。然而，要对它的模做许多基本操作，三角结构并不够用。

下面是这类问题的一个基本例子。给定 $\Dmod(K)$ 的两个模 $\sC_1$ 与 $\sC_2$，人们想谈论 $\Dmod(K)$-等变三角函子所构成的范畴
%
\[
   \Hom_{\Dmod(K)\mod}(\sC_1, \sC_2).
\]
%
不仅如此，我们还希望它们的行为像函数论中的对应物。例如，给定有限群 $K(\Fq)$ 的子群 $H_1(\Fq)$ 与 $H_2(\Fq)$，有典范同构
%
\begin{align*}
  \Hom_{K(\Fq)\mod}\big(\Fun(K(\Fq)/H_1(\Fq)),&\Fun(K(\Fq)/H_2(\Fq))
  \big)\\
  &\simeq
 \Fun\big( H_1(\Fq) \backslash K(\Fq) / H_2(\Fq)\big).
\end{align*}
%
人们希望对 $K$ 的代数子群 $H_1$ 与 $H_2$ 有类似的同构
%
\begin{align*}
\Hom_{\Dmod(K)\mod}\big(\Dmod(K/H_1),\Dmod(K/H_2)
\big)
\simeq
\Dmod( H_1 \backslash K / H_2).
\end{align*}
%

遗憾的是，即便 $K$ 是平凡群，构造合适的 Hom 空间也已经有麻烦。此时我们想要的是从 $\sC_1$ 到 $\sC_2$ 的三角函子的一个合理范畴。而这是三角范畴众所周知的老问题。例如，给定带态射 $F_1 \rightarrow F_2$ 的三角函子，如何把它补成一个正合三角，并不清楚。人们想逐点取锥，即对 $\sC_1$ 的每个对象 $c$ 令 $\on{Cone}(F_1 \rightarrow F_2)(c) \simeq \on{Cone}(F_1(c) \rightarrow F_2(c))$，但要把这变成一个函子，就会因 $\sC_2$ 中锥的非函子性而碰壁。

再看一个更精细的例子。人们希望能取 $\Dmod(K)$ 表示的极限、余极限与张量积，而在三角范畴中这些同样不易定义。

幸运的是，从三角范畴过渡到 dg 范畴，可以把这些问题一并解决。

\ssszh 我们先简述 dg 范畴的基本想法。回忆：表示论中出现的多数三角范畴，例如 $\fk\mod$，两个对象 $t_1$ 与 $t_2$ 之间的同态都算作某个向量空间复形的零次上同调。类似地，$t_1$ 与移位 $t_2[i]$（$i$ 为任意整数）之间的同态，就是这个复形的第 $i$ 次上同调。

具体地说，对阿贝尔范畴的导出范畴，可以取 $t_1$ 的投射消解 $p_1$ 或 $t_2$ 的单射消解 $i_2$\footnote{严格说来，由于讨论的可能是同调维数无穷的阿贝尔范畴上的无界复形，这里指的是 Bernstein 与 Spaltenstein 意义下的 K-投射消解与 K-单射消解 \cite{spaltenstein}。}，并使用相应的同态复形
%
\[
  \underline{\Hom}(p_1, t_2) \xrightarrow{\sim} \underline{\Hom}(p_1, i_2) \xleftarrow{\sim}
  \underline{\Hom}(t_1, i_2).
\]
%
我们强调，上述两个同构都是拟同构，也就是在上同调上诱导同构，但通常并不等同整条链复形本身。

dg 范畴的基本想法是：不仅记住这些复形的上同调，还要记住复形本身，并且只差一个同时的拟同构。熟悉导出范畴的读者不必对此感到意外。毕竟，导出范畴之所以有用，正是因为人们不只记住单个导出函子（例如 $\Ext^1$），还要记住计算它的复形。从三角范畴过渡到 dg 范畴，无非是把这一想法不仅用于对象，也用于态射。

\ssszh 直白地说，一个 $k$-线性 dg 范畴 $\sC$ 有一族对象。对每对对象 $c_1$ 与 $c_2$，有一个 $k$-向量空间复形 $\Hom(c_1, c_2)$；对每个三元组对象，有一个复形的复合态射
%
\[
\Hom(c_1, c_2) \otimes \Hom(c_2, c_3) \rightarrow \Hom(c_1, c_3),
  \]
%
它在朴素的意义下结合。

下面是 dg 范畴中的一些基本操作。首先，给定对象 $c$ 与整数 $i$，移位 $c[i]$（若存在）余表示函子
%
\[
\Hom(c[i], -) \simeq \Hom(c,-)[-i].
\]
%
类似地，给定零链 $f: c_1 \rightarrow c_2$，$f$ 的锥（若存在）余表示函子
%
\[
   \Hom( \on{Cone}(f), -) \simeq \on{Cone}( \Hom(c_2, -) \rightarrow
   \Hom(c_1, -))[-1],
\]
%
其中右边的锥是向量空间复形通常的锥。最后，直和按通常方式定义：当直和存在时，它余表示函子
%
\[
   \Hom( \underset{i}{\oplus} \hspace{.2mm} c_i, -) \simeq \underset{i}\Pi
   \hspace{.2mm} \Hom( c_i, -).
\]
%
若一个 dg 范畴对上述操作封闭，即对移位、锥与任意直和封闭，就称它余完备。\footnote{下文还默认任何余完备的 dg 范畴都是可表现的。这是关于其大小的技术性假设，所涉及的一切例子都满足。}

\ssszh 下面是余完备 dg 范畴的两个基本例子。

\begin{ex} $k$-向量空间的范畴 $\Vect$ 是余完备的 dg 范畴。具体地说，$\Vect$ 的对象是 $k$-向量空间的链复形 $(V_i, d_i)$，$i \in \mathbb{Z}$。

给定两个复形
$(V_i, d_V)$ 与 $(W_i, d_W)$，令
%
\[
\Hom(V_1, V_2) := \underline{\Hom}(V_1, V_2),
\]
%
其中 $\underline{\Hom}(V_1, V_2)$ 是这样的链复形：其 $i$-链为
%
\[
    \underline{\Hom}(V, W)^i := \underset{j}{\Pi} \hspace{.2mm} \Hom(V_j,
    W_{i+j}),
\]
%
微分则取通常的那个
%
\[
     d \phi := d_W \circ \phi -(-1)^i \phi \circ d_V, \quad \text{对 } \phi
     \in \underline{\Hom}(V,W)^i.
\]
%
特别地，$i$-闭链就是链复形的映射 $V \rightarrow W[i]$，$i$-边缘则是零伦的映射。\label{e:dgvect}
\end{ex}

\begin{ex}更一般地，设 $\sA$ 是 $k$-线性的 Grothendieck 阿贝尔范畴，例如 $k$ 上代数簇的 D-模，或某个 $k$-代数的左模。此时无界导出范畴
%
$
    \on{D}(\sA)
$
%
如下自然地提升为一个 dg 范畴：对象取 $\sA$ 中对象的 K-单射复形，同态则按例 \ref{e:dgvect} 中的方式定义。
\end{ex}

\ssszh\label{s:homdgcats} 给定两个余完备的 dg 范畴 $\sC_1$ 与 $\sC_2$，有一个余完备的 dg 范畴
%
\begin{equation} \label{e:bigbrainhoms}
\udHom(\sC_1, \sC_2),
\end{equation}
%
人们把它想成两者之间的函子范畴；其中的下标 $\on{cont}$ 将在下文第 \ref{s:dgcat} 小节进一步讨论。\eqref{e:bigbrainhoms} 的构造并不完全直接，因为要求它在 dg 范畴的拟等价下不变。\footnote{dg 范畴之间的朴素 dg 函子 $F: \sC \rightarrow \sC'$，由对象集之间的相容映射与态射复形之间的相容映射给出。若 (i) 它在同态复形上诱导拟同构，且 (ii) 对 $\sC'$ 的任意对象 $c'$，存在 $\sC$ 的对象 $c$ 与零链 $F(c) \rightarrow  c'$，使下式中的映射为拟同构
%
\begin{equation} \label{e:bigbrainboms}
  \Hom_{\sC'}(-, F(c)) \simeq \Hom_{\sC'}(-, c'),
\end{equation}
%
就称 $F$ 是拟等价。这可以改述如下。取所有同态复形的零次上同调之后，朴素 dg 函子 $F$ 诱导出相应「同伦」范畴之间的函子
%
\[
\on{Ho}(F): \on{Ho}(\sC) \simeq \on{Ho}(\sC').
\]
%
在 $\sC$ 与 $\sC'$ 都对移位封闭的假设下，$F$ 是拟等价当且仅当 $\on{Ho}(F)$ 是通常范畴的等价。

直白地说，实践中人们并不在乎严格的同态复形本身，而在乎复形到拟同构为止。这既引出了拟等价的概念，也引出这样的愿望：函子范畴应当在拟等价下不变。}

不过，可以像导出范畴中的映射那样，用模型范畴论的手段来构造它：在 $\sC_1$ 与 $\sC_2$ 的适当替换之间取朴素的 dg 函子范畴；参见 \cite{toenlectures} 及其中的文献。

我们来描述 \eqref{e:bigbrainhoms} 的一些基本性质。第一，\eqref{e:bigbrainhoms} 的对象集是 Keller 意义下的拟函子 \cite{kellerderiving}，并且是连续的，即与任意直和交换。第二，$F_1$ 与 $F_2$ 之间态射的 $i$-闭链集，在（余）纤维替换之后，就是自然变换 $F_1 \rightarrow F_2[i]$。

\begin{re} 需要说明，这里的「连续」指与余极限交换；而在文献的其他地方，它指与极限交换。\end{re}

下面是这类函子范畴的一些重要例子。

\begin{ex} 对任意余完备的 dg 范畴 $\sC$，有典范等价
%
\[
     \udHom(\Vect, \sC) \simeq \sC,
\]
%
它把 $\sC$ 的对象 $c$ 与张量以 $c$ 的 dg 函子对应起来，即 $V \mapsto V \otimes c$。

具体地说，这个 dg 函子把向量空间 $k^{\oplus I}$（$I$ 为集合）对应到对象 $c^{\oplus I}$；给定映射 $k^{\oplus I} \rightarrow k^{\oplus J}$，所得的映射 $c^{\oplus I} \rightarrow c^{\oplus J}$ 来自映射 $k \rightarrow \on{Hom}_\sC(c, c)$，即 $\sC$ 的 $k$-线性性。有界向量空间复形的像通过取锥得到，无界复形的像则通过取余极限得到。
\end{ex}

\begin{ex} 下面这个例子是第 \ref{sss:endos} 小节中函数论讨论的一个版本，后文会反复用到；对代数簇它归功于 \cite{toenpaper}，对一大类导出叠则由 \cite{benzvifrancisnadler} 推广。

对代数簇 $X$ 与 $Y$，有典范等价
%
\[
   \Hom_{\DGCat}(\QCoh(X), \QCoh(Y)) \simeq \QCoh(X \times Y),
\]
%
它把乘积上的层 $\mathscr{K}$ 对应到相应的积分变换，即
%
\[
   \QCoh(X) \xrightarrow{\pi_X^*} \QCoh(X \times Y) \xrightarrow{\mathscr{K} \otimes - } \QCoh(X \times Y) \xrightarrow{\pi_{Y, *}} \QCoh(Y),
\]
%
其中 $\pi_X: X \times Y \rightarrow X$ 与 $\pi_Y: X \times Y \rightarrow Y$ 表示投影。

类似地，有典范等价
%
\[
   \Hom_{\DGCat}(\Dmod(X), \Dmod(Y)) \simeq \Dmod(X \times Y),
\]
%
它同样把乘积上的 D-模 $\mathscr{K}$ 对应到积分变换
%
\[
   \Dmod(X) \xrightarrow{\pi_X^!} \Dmod(X \times Y) \xrightarrow{ \mathscr{K} \overset ! \otimes - } \Dmod(X \times Y) \xrightarrow{\pi_{Y, *}} \Dmod(Y).
\]
%
\end{ex}

\ssszh\label{s:dgcat} 要谈论 dg 范畴的（余）极限，需要用到如下结构。存在一个 $(\infty, 1)$-范畴，记作 $\DGCat$，其对象是余完备的 dg 范畴，1-态射是连续的拟函子；参见 \cite{gaitsroz} 第 1 章及其中的文献。以下是它的一些基本性质。

第一，它在 $\infty$-范畴的意义下具有一切（余）极限。粗略地说，dg 范畴的极限 $\varprojlim \sC_\alpha$ 的对象，是各 $\sC_\alpha$ 中同伦相容的一族对象；态射则是相应的映射复形的逆向极限。\footnote{这里再次强调，这正是三角范畴不够用的一个基本之处。也就是说，由于逆向极限不正合，人们确实需要复形本身，而不能只看它们的上同调。}

dg 范畴图的余极限则由下述等式刻画
%
\[
\udHom(\varinjlim \sC_\alpha, -) \simeq \varprojlim
\udHom(\sC_\alpha, -)
\]
%
与此相关，dg 范畴的极限满足等式
%
\[
 \udHom(-, \varprojlim \sC_\beta) \simeq \varprojlim
 \udHom(-, \sC_\beta).
\]
%

一般 dg 范畴的余极限中的对象或态射并不容易描述（其实一般只对极限才指望有直白的描述，因为它们由映射性质刻画）。不过为了帮助读者定位，我们描述一个可以显式描述的重要例子。

\begin{ex}\label{e:colimitcats}回忆：若函子 $$\Hom(c, -)$$ 与余极限交换，就说 dg 范畴 $\sC$ 的对象 $c$ 是\emph{紧}的。应当把它们想成小对象，例如在 $\Vect$ 中，一个对象紧当且仅当它等价于有限维向量空间的有界复形。与之相关，若每个对象都能写成紧对象的余极限，就说 $\sC$ 是紧生成的。

给定紧生成范畴的滤过图 $\sC_\alpha$，其中每个 1-态射 $\sC_\alpha \rightarrow \sC_{\alpha'}$ 都把紧对象送到紧对象，则余极限
%
\[
  \varinjlim \sC_{\alpha}
\]
%
有如下描述：它仍由各 $\sC_\alpha$ 的紧对象嵌入而紧生成，这些对象之间的同态则是余极限各步上同态的滤过余极限。
\end{ex}

\ssszh $\DGCat$ 的一个很有用的性质是，可以把余完备的 dg 范畴张量起来
%
\[
   \sC \text{ 与 } \sD \rightsquigarrow \sC \otimes \sD.
\]
%
更确切地说，$\DGCat$ 关于 Lurie 张量积是对称幺半 $\infty$-范畴。

在实际的许多例子中，下面这种可以动手算的描述很有用。

\begin{ex}设 $\sC$ 与 $\sD$ 是紧生成的，比如分别由对象集 $c_\alpha$ 与 $d_\beta$ 紧生成。则张量积由简单张量 $c_\alpha \otimes d_\beta$ 构成的集合紧生成，其同态由 K\"unneth 公式给出，即
%
\[
\Hom( c_\alpha \otimes d_\beta, c_{\alpha'} \otimes d_{\beta'}) \simeq
\Hom(c_\alpha, c_{\alpha'}) \otimes \Hom(d_\beta, d_{\beta'}).
\]
%
\end{ex}

上述例子有一些特殊情形，希望它们能让读者感受到，这个张量积与读者从函数-层对应中得来的直觉是一致的，参见第 \ref{ss:opsfinitesets} 小节。

\begin{ex} 给定代数簇 $X$ 与 $Y$，拟凝聚层范畴之间有典范等价
%
\[
   \on{QCoh}(X) \otimes \on{QCoh}(Y) \simeq \on{QCoh}(X \times Y).
\]
%
具体地说，给定 $\on{QCoh}(X)$ 的对象 $\sM$ 与 $\on{QCoh}(Y)$ 的对象 $\sN$，左边的简单张量 $\sM \otimes \sN$ 被送到外积对象 $$\sM \boxtimes \sN
\simeq \pi_X^* \sM
\overset{}\otimes \pi_Y^* \sN.$$
\end{ex}

\begin{ex} 仍取上例中的 $X$ 与 $Y$。D-模范畴之间也有典范等价
%
\[
     \Dmod(X) \otimes \Dmod(Y) \simeq \Dmod(X \times Y).
\]
%
同样地，给定 $\Dmod(X)$ 的对象 $\sM$ 与 $\Dmod(Y)$ 的对象 $\sN$，左边的简单张量 $\sM \otimes \sN$ 被送到外积对象
%
\[
      \sM \boxtimes \sN \simeq \pi_X^! \sN \overset{!}{\otimes} \pi_Y^! \sM.
\]
%
\zy{[译者按（上下文推断，无外部依据）：原文此处把 $\sM$ 与 $\sN$ 的位置写颠倒了。按上例的样式，应为 $\pi_X^! \sM
\overset{!}{\otimes} \pi_Y^! \sN$；原文写作 $\pi_X^! \sN \overset{!}{\otimes} \pi_Y^! \sM$。本译本照原文排印，读者按此理解即可。]}
\end{ex}

\begin{ex}给定 $k$-代数 $A$ 与 $B$，有典范等价
%
\[
    A\mod \otimes \hspace{.2mm} B\mod \simeq (A \underset{k} \otimes
    B)\mod.
\]
%
注意这推广了前两个例子的仿射情形。
\end{ex}

需要一提的是，这个幺半结构是幺半的，其单位是向量空间范畴 $\Vect$。特别地，对任意 dg 范畴 $\sC$，有典范等价
%
\[
   \Vect \otimes \hspace{.7mm}  \sC \simeq \sC.
\]
%
读者不妨在 $\sC$ 紧生成的情形直接验证这一点。

\ssszh 给定 dg 范畴 $\sC,\sD$，第 \ref{s:homdgcats} 小节中的范畴 $\udHom(\sC,
\sD)$ 关于 $\DGCat$ 上的幺半结构是内 Hom 对象。

这里有一个技术要点：它并不等同于把 $\sC$ 与 $\sD$ 看成 $\DGCat$ 的对象之后二者之间的态射范畴。准确地说，态射范畴是把 $\udHom(\sC, \sD)$ 截断到 $\infty$-群胚，也就是丢掉不可逆的自然变换。人们预期 $\DGCat$ 有典范的提升，成为对称幺半 $(\infty,2)$-范畴，其中同态确实就是 $\udHom(-, -)$；但据我们所知，文献中目前还没有这一构造。

\ssszh 借助幺半结构，可以谈论 $\DGCat$ 的代数对象，即幺半 dg 范畴。粗略地说，它由一个 dg 范畴 $\sA$ 连同乘法
%
\[
    \sA \otimes \sA \rightarrow \sA,
\]
%
组成，该乘法在相容同伦的意义下结合。我们强调，相容同伦是额外的结构，是幺半 dg 范畴完整数据的一部分。

下面是幺半 dg 范畴的一些基本例子。

\begin{ex}\label{ex:qcohalg} 对代数簇 $X$，$\QCoh(X)$ 在 $*$-张量积下是幺半 dg 范畴，$\Dmod(X)$ 在 $!$-张量积下是幺半 dg 范畴。

给定映射 $f: X_1 \rightarrow X_2$，相应的拉回函子
%
\[
f^*: \QCoh(X_2) \rightarrow \QCoh(X_1) \quad \text{与} \quad f^!: \Dmod(X_2)
\rightarrow \Dmod(X_1)
\]
%
带有典范的幺半结构数据。
\end{ex}

\begin{ex} 设簇 $Y$ 是簇范畴中的代数对象。也就是说，$Y$ 带有映射
%
\[
    \mu: Y \times Y \rightarrow Y,
\]
%
它在显然的意义下结合。此时 $\QCoh(Y)$ 与 $\Dmod(Y)$ 都获得「卷积」幺半结构，由沿 $\mu$ 的 $*$-推前给出。

给定代数的映射 $f:Y_1 \rightarrow Y_2$，相应的推前映射
%
\[
 f_*: \QCoh(Y_1) \rightarrow \QCoh(Y_2) \quad \text{与} \quad f_*: \Dmod(Y_1)
 \rightarrow \Dmod(Y_2)
\]
%
带有典范的幺半结构数据。

我们主要关心的是 $Y = K$ 为代数群、$\mu$ 为其乘法这一情形，由此得到先前见过的群代数 $\Dmod(K)$（更确切地说，是它的典范 dg 提升）。

类似地，给定群的同态 $\phi: K_1 \rightarrow K_2$，得到群代数的同态
%
\[
   \phi_*: \Dmod(K_1) \rightarrow \Dmod(K_2).
\]
%
\end{ex}

\begin{re} 把前两段中的 $\Dmod(K)$ 换成 $\QCoh(K)$，可以得到另一个卷积下的幺半 dg 范畴。$\Dmod(K)$ 的模有时称作强范畴表示，$\QCoh(K)$ 的模则称作弱范畴表示。

有一个或许能澄清问题的注：给定从簇到幺半 $\infty$-范畴 $(\sS, \otimes)$ 的任意幺半函子，都可得到类似的、把 $K$ 线性化的卷积代数。特别地，可以谈论 $\on{Fun}(K(\Fq))$ 的更高范畴版本，例如把层换成 $K$ 上层的范畴。\end{re}

\ssszh 给定幺半 dg 范畴 $\sA$，有与之相伴的左模与右模 $\infty$-范畴，分别记作
%
\[
  \sA\mod \quad \text{与} \quad \rmod\sA.
\]
%
说得不那么正式，一个左模由一个余完备 dg 范畴 $\sM$ 连同作用映射
%
\[
  \sA \otimes \sM \rightarrow \sM
\]
%
组成，该作用在相容同伦的意义下「结合」；右模类似。同样，相容同伦是指定模的数据的一部分。

给定两个左模 $\sM_1$ 与 $\sM_2$，有 $\sA$-等变函子构成的 dg 范畴
%
\[
   \Hom_{\sA\mod}(\sM_1, \sM_2)
\]
%
（$\sA$ 上右模的情形类似）。具体地说，它可以借助 bar 复形、作为 dg 范畴半余单纯图的极限来计算
%
\begin{equation} \label{e:hombar}
   \scriptsize
   \xymatrix{  \udHom(\sM_1, \sM_2)
   \ar@<-.4ex>[r] \ar@<.4ex>[r] &  \hspace{.5mm} \udHom(\sA \otimes \sM_1,
   \sM_2)
   \ar@<-.8ex>[r] \ar@<.8ex>[r] \ar[r] & \hspace{.5mm} \udHom(\sA^{\otimes 2}
   \otimes \sM_1, \sM_2)   \hspace{.5mm} \cdots \hspace{.5mm},
   }
   \end{equation}
%
其中各箭头按 bar 消解的标准模式，由 $\sA$ 在 $\sM_1$ 与 $\sM_2$ 上的作用映射、以及乘法映射 $\sA^{\otimes 2} \rightarrow \sA$ 诱导。

此外，给定右模 $\sM$ 与左模 $\sN$，可以作它们的张量积
%
\[
     \sM  \underset{\sA}{\otimes} \sN.
\]
%
与同态的情形一样，它可以借助 bar 消解、作为余完备 dg 范畴半单纯图的余极限来计算
\zy{[译者按（上下文推断，无外部依据）：下式第二项中的 $N$，按全式其余各项的写法应作花体的 $\sN$。原文此处漏写花体，本译本照原文排印。]}
%
\begin{equation} \label{e:tensorbar}
\xymatrix{  \cdots \hspace{.5mm} \sM \otimes \sA^{\otimes 2} \otimes N
\ar@<-.8ex>[r] \ar@<.8ex>[r] \ar[r]   &  \hspace{.5mm} \sM \otimes \sA \otimes
\sN \ar@<-.4ex>[r]
\ar@<.4ex>[r]  &  \hspace{.5mm} \sM \otimes  \sN}.
\end{equation}
%


现在给出几个基本例子。

\begin{ex} 对任意左 $\sA$-模 $\sM$，有显然的等价
%
\[
   \on{Hom}_{\sA\mod}(\sA, \sM) \simeq \sM \simeq \sA \underset \sA \otimes
   \sM
\]
%
右模的情形类似。
\end{ex}

\begin{ex} \label{e:fibprodvar}给定簇的映射 $X \rightarrow Z
\leftarrow Y$，如例 \ref{ex:qcohalg} 那样，得到 $\QCoh(Z)$ 在 $\QCoh(X)$ 与 $\QCoh(Y)$ 上的典范作用。于是它们的张量积由下式给出
%
\[
    \QCoh(X)
    \underset{\QCoh(Z)}{\otimes} \QCoh(Y) \simeq \QCoh( X \underset{Z}{\times}
    Y),
\]
%
其中上述纤维积是导出的，即得到导出概形。\footnote{直白地说，这是逐仿射地用导出环 $\OO_X
\underset{\OO_Z}{\overset{L}{\otimes}} \OO_Y$ 粘起来的，与通常纤维积的情形完全一样。人们把结构层的负上同调群想成给簇添上「幂零绒毛」，就像一般概形的非既约性那样，只不过后者或许更为读者所熟悉。本书会遇到的一切导出概形都以这种方式产生。}
具体地说，这个等价把简单张量 $\sM \otimes \sN$ 送到它们的外积在 $X \underset{Z}{\times} Y$ 上的限制。

类似地，它们的同态也由下式给出
%
\[
 \Hom_{\QCoh(Z)}(\QCoh(X), \QCoh(Y)) \simeq \QCoh( X \underset{Z}{\times}
 Y).
\]
%
具体地说，把纤维积上的拟凝聚层 $\sS$ 对应到积分变换
%
\[
  \QCoh(X) \xrightarrow{ \pi_X^*} \QCoh(X \underset Z \times Y) \xrightarrow{ -
  \overset * \otimes \sS} \QCoh(X \underset Z \times Y) \xrightarrow{\pi_{Z,
  *}} \QCoh(Y).
\]
%

对 D-模，同样的断言作必要的调整后成立。也就是说，有等价
%
\begin{align*}
  \Dmod(X) \underset{\Dmod(Z)}{\otimes} \Dmod(Y) &\simeq \Dmod(X \overset{\on{cl}}{{\underset Z {\times}}} Y) \text{，以及}  \\
  \Hom_{\Dmod(Z)}(\Dmod(X), \Dmod(Y)) &\simeq \Dmod(X \overset{\on{cl}}{{\underset Z {\times}}} Y),
\end{align*}
%
它们同样由外张量积与积分变换给出。这里上标 $\on{cl}$ 表示经典的，即非导出的纤维积。但要注意，当导出纤维积的结构层只有有限多个非零上同调时（例如 $Z$ 光滑时），由于 D-模范畴在幂零加厚下不变，取簇的导出纤维积还是经典纤维积并不重要。

与我们已经遇到的其它层论陈述一样，这些等式是下面这个函数等式在层论中的类似物。给定有限集 $S$，把 $\Fun(S)$ 看作逐点乘法下的代数。给定有限集 $X \rightarrow Z
\leftarrow Y$，便有向量空间的典范同构
%
\[
    \Fun(X) \underset{\Fun(Z)}{\otimes} \Fun(Y) \simeq
    \Hom_{\Fun(Z)}(\Fun(X), \Fun(Y)) \simeq \Fun(X \underset Z \times Y).
\]
%
\end{ex}

\subsection{范畴表示论的基础}

\ssszh 有了这些一般性工具，我们来描述范畴表示论中的一些基本操作，并借此重新得到第 \ref{s:bbloc} 节中遇到的定理 \ref{t:loc+equiv} 与定理 \ref{t:lochchmod+eq}。

回忆：对代数群 $K$，有与之相伴的范畴表示的 $\infty$-范畴
%
\[
   \Dmod(K)\mod.
\]
%
表示论中别处也一样，$K$ 上的取逆诱导了 $\Dmod(K)$ 与其反向幺半 dg 范畴之间的同构；特别地，可以在左模与右模之间自由切换。

\ssszh 有一族基本的表示可以如下几何地构造出来。

给定带 $K$ 作用的簇 $X$，得到 $\Dmod(K)$ 在 $\Dmod(X)$ 上的典范作用。具体地说，底层的二元乘法由
%
\[
   \Dmod(K) \otimes \Dmod(X) \simeq \Dmod(K \times X) \xrightarrow{a_*}
   \Dmod(X)
\]
%
给出，与第 \ref{ex:actspace} 小节中函数论的情形完全一样。给定 $K$-簇的映射 $f: X \rightarrow X'$，沿 $f$ 的 $*$-推前与 $!$-拉回函子都带有典范的 $K$-等变数据。

\ssszh 特别地，取 $X = \on{pt}$，就得到 $K$ 在
%
\[
    \on{Vect} \simeq \Dmod(\on{pt})
\]
%
上的平凡表示。具体地说，若用 $\pi: K \rightarrow \on{pt}$ 记投影，则作用映射形如
%
\[
    \Dmod(K) \otimes \Vect \rightarrow \Vect, \quad \quad \sM \otimes V \mapsto
    \pi_*(\sM) \underset k \otimes  V.
\]
%

于是对 $K$ 的任意表示 $\sC$，可以作它的不变量与余不变量
%
\[
 \sC^K := \Hom_{\Dmod(K)\mod}(\on{Vect}, \sC) \quad \text{与} \quad \sC_K :=
 \on{Vect} \underset {\Dmod(K)} \otimes \sC.
\]
%
正如有限群通常表示的情形，二者典范地等同，论证如下。

由 $K$ 的光滑性，有 $\Dmod(K)$-等变的伴随对
%
\[
     \pi^![-2 \dim K]: \Dmod(K) \rightleftarrows \Vect: \pi_*,
\]
%
对它作用 $\Hom_{\Dmod(K)\mod}( -, \sC)$，便诱导出伴随对
%
\[
    \on{Oblv}: \sC^K \rightleftarrows \sC: \on{Av}_*.
\]
%
正如记号所暗示的，人们把 $\on{Oblv}$ 想成遗忘函子，把 $\on{Av}_*$ 想成平均函子。出于形式上的理由，平均函子经过余不变量分解：
%
\[
 \on{Av}_*: \sC_K \rightarrow \sC^K.
\]
%
根据 Beraldo 与 Gaitsgory 的一个很有用的定理 \cite{Bearaldo}、\cite{whitlocglob}，它是等价。

下面给出（余）不变量的一些基本例子。

\begin{ex}给定 $\Dmod(K)$ 的表示 $\sC$ 与 $\sD$，有典范等价
%
\[
       \Hom_{\Dmod(K)\mod}(\sC, \sD) \simeq \udHom(\sC, \sD)^K,
\]
%
其中的不变量关于 $\Dmod(K)$ 在映射空间上的对角作用来取。
\end{ex}

\begin{ex} \label{e:dmodquotstack} 对 $K$-簇 $X$，$\Dmod(X)$ 中 $K$-等变的对象典范地等同于商叠 $X/K$ 上 D-模的范畴，即
%
\[
     \Dmod(X)^K \simeq \Dmod(X/K).
\]
%
直白地说，D-模满足光滑下降，而图 \eqref{e:hombar} 正是沿覆盖 $X
\rightarrow X/K$ 计算 $X/K$ 上 D-模的「$\check{\on{C}}$ech 复形」。
\end{ex}

\begin{ex} 在 $\Dmod(K)$ 对其李代数表示范畴 $\fk\mod$ 的自然作用下，不变量等同于群本身的代数表示，即
%
\[
      \fk\mod^K \simeq K\mod.
\]
%
更一般地，给定同态 $K \rightarrow K'$，范畴
%
\[
 \fk'\mod^K
\]
%
典范地等同于对 $(\fk', K)$ 的 Harish-Chandra 模范畴。
\end{ex}

\begin{ex} 还有另一类例子，它们同时给出 $\fk\mod$ 上作用的另一种视角。设 $A$ 是带 $K$ 的自同构作用的 $k$-代数。此时得到 $\QCoh(K)$ 在 $A\mod$ 上的作用。

在单位元处求导，$K$ 在 $A$ 上的作用给出 $\fk$ 通过导子的作用 $\alpha$。再设这个作用有一个平凡化，即李代数的同态 $$\tau: \fk \rightarrow A,$$使得 $\alpha$ 与经 $\tau$ 给出的伴随作用一致。此时得到 $\Dmod(K)$ 在 $A\mod$ 上的作用。 \end{ex}

\begin{re}\label{r:invariantstwice}这里记录一个与通常表示情形基本不同之处。取不变量这一操作并非幂等，即平凡映射
%
\[
   \on{Oblv}_{\sC^K}: (\sC^K)^K \rightarrow \sC^K
\]
%
通常不是同构。取而代之的是典范等价
%
\[
   (\sC^K)^K \simeq (\sC^K) \otimes \Dmod(\on{pt/K}).
\]
%

之所以如此，是因为除非 $K$ 可缩，对任意（非零）表示 $\sD$，遗忘映射 $\on{Oblv}_{\sD}: \sD^K \rightarrow \sD$ 都不会是全忠实的，即 $\sD^K$ 与 $\sD$ 中的 Ext 并不一致。确实，在最普遍的 $\sD = \Dmod(K)$ 情形，有
%
\[
\Dmod(K)^K \simeq \Dmod(K / K) \simeq \Vect,
\]
%
其生成元是常值层。特别地，它在 $\Dmod(K)$ 中的自 Ext 是 $\pi_*(\underline{k})$，即 $K$ 的上同调；而它在 $\Vect$ 中的自 Ext 则是平凡的。

一般地，复合
%
\[
    \sD^K \xrightarrow{\on{Oblv}_{\sD}} \sD \xrightarrow{\on{Av}_*} \sD^K
    \xrightarrow{\on{Oblv}_{\sD}} \sD
\]
%
典范地等同于 $\on{Oblv}_{\sD} \otimes \pi_*(\underline{k})$，也就是把 $K$-等变对象底层的对象与 $K$ 的上同调作张量。读者不妨借助函数-层对应，把它与通常表示的平均过程中出现的有限群的基数因子作比较。

但要注意，若 $K$ 幂幺，则 $\on{Oblv}$ 是全忠实的，因为此时 $K$ 可缩，从而 $\pi_*(\underline{k}) \simeq k$。 \end{re}

\ssszh 讨论完不变量函子，我们来说明如何得到定理 \ref{t:loc+equiv} 与定理 \ref{t:lochchmod+eq}；先从后者开始。

给定代数群的映射 $H \rightarrow K$，有诱导函子与限制函子
%
\[
  \on{ind}_H^K: \Dmod(H)\mod \rightarrow \Dmod(K)\mod,
\]
%
%
\[
  \on{res}_K^H: \Dmod(K)\mod \rightarrow \Dmod(H)\mod.
\]
%
直白地说，$\on{res}_K^H$ 沿幺半函子
%
\[
         \Dmod(H) \rightarrow \Dmod(K)
\]
%
限制 $\Dmod(K)$ 在 dg 范畴 $\sC$ 上的作用，而 $\on{ind}_H^K$ 则把带 $\Dmod(H)$ 作用的 dg 范畴 $\sD$ 张量上去，即
%
\[
      \sD \mapsto \Dmod(K) \underset{ \Dmod(H) } \otimes \sD.
\]
%
正如有限群的 Frobenius 互反性，诱导与限制典范地互为左伴随，也互为右伴随。

\label{ss:invariantscosets} 特别地，得到 $\Dmod(K)$ 表示上的 $H$-不变量函子，即
%
\begin{align*}
     \big(\on{Res}_K^H(-)\big)^H &\simeq   \Hom_{\Dmod(H)\mod}( \on{Vect},
     \on{res}_K^H-).
\intertext{由伴随性，可以把它改写成 }
&\simeq  \Hom_{\Dmod(K)\mod}( \on{ind}_H^K \on{Vect},
-) \\ & \simeq \Hom_{\Dmod(K)\mod}( \Dmod(K) \underset{\Dmod(H)}{\otimes} \Vect
,
-)
\intertext{由 $H$ 的不变量与余不变量的等同、以及例 \ref{e:dmodquotstack}，可以认出这是}
& \simeq \Hom_{\Dmod(K)\mod}( \Dmod(K/H)
,
-).
\end{align*}

结论是：任意 $\Dmod(K)$ 表示的 $H$-不变量，经前复合而带有
%
\[
\Hom_{\Dmod(K)\mod}( \Dmod(K/H), \Dmod(K/H)) \simeq \Dmod(H \backslash K / H)
\]
%
的典范作用，也就是我们在第 \ref{s:bbloc} 节末尾遇到的 Beilinson--Drinfeld 的 Hecke 代数作用。

\ssszh 最后，对光滑的 $K$-簇 $X$，我们简述如何得到 $\Dmod(K)$-等变的伴随对
%
\[
    \on{Loc}: \fk\mod \rightleftarrows \Dmod(X): \Gamma(X, -).
\]
%
首先，从 $\OO_K$-模诱导到 $\sD_K$-模给出幺半函子
%
\[
      \QCoh(K) \rightarrow \Dmod(K).
\]
%
类似地，有 $\QCoh(K)$ 在 $\Dmod(X)$ 上的典范作用，使得遗忘函子 $\Dmod(X) \rightarrow \QCoh(X)$ 带有典范的 $\QCoh(K)$-等变数据。特别地，整体截面函子，即复合
%
\[
    \Dmod(X) \rightarrow \QCoh(X) \xrightarrow{\pi_*} \QCoh(\pt) \simeq \Vect,
\]
%
带有典范的 $\QCoh(K)$-等变数据。

然而，对任意 $\Dmod(K)$ 的表示 $\sC$，由伴随性有
%
\begin{align*}
\Hom_{\QCoh(K)\mod}( \on{Res} \sC, \Vect)
&\simeq \Hom_{\Dmod(K)\mod} \big(\sC,\\
&\qquad \Hom_{\QCoh(K)\mod}(\Dmod(K), \Vect)\big) \\
&\simeq \Hom_{\Dmod(K)\mod}(\sC, \fk\mod),
\end{align*}
%
其中最后的等价用了 $\fk\mod$ 与 $K$ 上弱 $K$-等变 D-模的等同，参见第 \ref{ss:weakinv} 小节。这就给出整体截面函子所需的 $\Dmod(K)$-等变性。至于局部化，可以利用 Gaitsgory \cite{whitlocglob} 的一个漂亮事实：$\Dmod(K)\mod$-等变函子的任意伴随都继承相容的等变数据。

\subsection{补遗一：特征标层}

\label{ss:charsheaves}
\ssszh 本节开头略述了 dg 范畴的一些预备知识，随后讨论了范畴表示论中若干基本定义与构造，足以重新解释第 \ref{s:bbloc} 节预告过的一些结果。

本节余下部分，我们想再综述几条基本结果与构造。不过读者也可以直接跳到第 \ref{s:catloop} 节。

\ssszh 特征标层是群上（某些）伴随等变的 D-模。Lusztig 着眼于它们与有限李型群的不可约特征标之间的密切关联而开始了这项研究，尤其是在约化情形 \cite{luschar1, luschar23, luschar4, luschar5}。这是函数-层对应中一个格外丰富而不平凡的例子。

当然，有限群的特征标并不只是群上伴随等变的函数，而是从表示取迹得来的。人们可以问：特征标层是否有这样的解释？答案是有的，下面就来讨论。

\ssszh[对偶性] \label{ss:dualizable}回忆：要定义有限群表示的特征标，表示必须是有限维的，才能安全地取迹。在范畴情形中，「小到可以取迹」的类比物如下。

借助 $\DGCat$ 上的幺半结构，可以谈论可对偶化的 dg 范畴。在 $\DGCat$ 与 $\Vect$ 的类比下，它们对应于（有限维向量空间的有界复形）；确实，dg 范畴 $\sC$ 可对偶化当且仅当自然映射
%
\[
     \sC \rightarrow \udHom( \udHom(\sC, \Vect), \Vect)
\]
%
是等价。

对可对偶化的 dg 范畴，有典范等价 $\sC^\vee \simeq \udHom(\sC,
\Vect)$，即它在幺半意义下的对偶就是它的连续 dg-模。对任意 $\sS$，自然映射
%
\[
    \sS \otimes \sC^\vee \rightarrow \udHom(\sC, \sS)
\]
%
是等价。\footnote{把 $\sC^\vee$ 换成 $\Hom(\sC, \Vect)$，这又是可对偶化 dg 范畴的一条刻画。}更一般地，对任意 dg 范畴 $\sS$ 与 $\sS'$，有等价
%
\[
\Hom( \sS \otimes \sC, \sS') \simeq \Hom( \sS, \sC^\vee \otimes \sS').
\]
%

设 $\sA$ 是幺半 dg 范畴。由函子性，若 $\sC$ 是左 $\sA$-模，则 $\sC^\vee$ 是右 $\sA$-模。特别地，对可对偶化的 $\Dmod(K)$ 表示 $\sC$，它的对偶获得 $\Dmod(K)$ 的典范作用，即可以作反变表示。

在讨论如何定义可对偶化 $\Dmod(K)$ 表示的特征标层之前，我们先谈可对偶化范畴与模的一些基本例子。首先，下面这条一般事实保证了实践中遇到的多数 dg 范畴都可对偶化。

\begin{ex}若范畴 $\sC$ 紧生成，则它可对偶化。此外，它的对偶也紧生成，并且 $\sC$ 的紧对象与非余完备 dg 范畴 $\sC^\vee$ 的紧对象的反范畴之间有典范等价。
\end{ex}

上例有一些有用的特殊情形。

\begin{ex} 给定 $k$-代数 $A$，其左模范畴 $A\mod$ 可对偶化，对偶是右模范畴 $\on{mod-}A$。具体地说，求值映射
%
\[
   \rmod A \otimes A\mod \rightarrow \Vect
\]
%
把右模 $M$ 与左模 $N$ 送到它们的张量积 $M
\otimes_A N$。
\end{ex}

\begin{ex} \label{ex:dmodselfdual}给定簇 $X$，其 D-模范畴通过紧 D-模上的 Verdier 对偶而典范地自对偶。具体地说，所得的求值映射由复合
%
\[
 \Dmod(X) \otimes \Dmod(X) \simeq \Dmod(X \times X) \xrightarrow{\Delta^!}
 \Dmod(X) \xrightarrow{\pi_*} \Vect
\]
%
给出，其中与第 \ref{s:funselfdual} 小节一样，用 $\Delta: X
\rightarrow X \times X$ 记对角映射，用 $\pi: X \rightarrow \on{pt}$ 记到点的投影。

若 $X$ 是 $K$-簇，这个配对典范地 $\Dmod(K)$-等变，即把 $\Dmod(X)$ 与其反变表示等同起来。把 $\QCoh(X)$ 看作 $\QCoh(K)$ 的表示，作必要的调整后也有同样的陈述。
\end{ex}

\ssszh[特征标层] 现在可以说明如何取一个表示的特征标层。$\Dmod(K)$ 在范畴 $\sC$ 上的作用，等同于一个幺半函子
%
\[
\Dmod(K) \xrightarrow{\rho} \udHom(\sC, \sC).
\]
%
若 $\sC$ 可对偶化，则由等同 $\udHom(\sC, \sC) \simeq
\sC^\vee \otimes \sC$，它的自同态也可对偶化。于是可以把作用映射对偶化，得到矩阵系数的 $\Dmod(K \times K)$-等变函子
%
\[
   \udHom(\sC, \sC) \simeq \udHom(\sC, \sC)^\vee \xrightarrow{\rho^\vee}
   \Dmod(K)^\vee \simeq \Dmod(K).
\]
%

特别地，过渡到对角的 $K$-不变量，得到函子
%
\begin{equation} \label{e:mcequivfunc}
   \Hom_{\Dmod(K)}(\sC, \sC) \rightarrow \Dmod(K / K ).
\end{equation}
%
把它用于恒等自同态 $\id_{\sC}$，就得到 $K$ 上的伴随等变 D-模 $\chi_\sC$，即 $\sC$ 的\emph{特征标层}。

\ssszh 下面给出 $\chi_\sC$ 的更可动手的描述。给定可对偶化的范畴 $\sC$ 与自函子 $f: \sC \rightarrow \sC$，可以取它的迹，那是一个向量空间复形。与线性代数中完全一样，它由复合
%
\[
\Vect \xrightarrow{\id_{\sC}} \sC \otimes \sC^\vee \xrightarrow{f \otimes
\on{id}_{\sC^\vee}} \sC \otimes \sC^\vee \xrightarrow{\on{ev}} \Vect
\]
%
计算。也就是说，$\Vect$ 的这个复合自函子由与某个向量空间作张量给出，而这个向量空间按定义就是 $\on{tr}(f, \sC)$。要说明的是，当 $f$ 就是 $\sC$ 的恒等自同态时，所得的复形称作 $\sC$ 的 Hochschild 同调。

理解了在当下范畴情形中如何取迹之后，我们就能相当显式地描述特征标层了。也就是说，它的茎承载 $\sC$ 相应自同构的迹，与通常表示论中完全一样。即：对任意闭点 $g: \pt \rightarrow K$，有典范同构
%
\[
     g^!(\chi_\sC) \simeq \on{tr}( \delta_g \star -, \sC).
\]
%

\ssszh 下面提几个特点与例子。

\begin{ex} 可对偶化模之间的 $\Dmod(K)$-等变函子
$\sC \rightarrow \sD$
不一定诱导特征标层之间的映射。但是，若它带有连续的右伴随，则确实得到适当函子性的映射 $\chi_\sC
\rightarrow
\chi_{\sD}$。

要说明的是，在实际情形中 $\sC$ 是紧生成的，而「有右伴随」这一条件等价于把紧对象送到紧对象，也就是没有「小」的东西变成「大」的。
\end{ex}

\begin{ex} 范畴表示上的诱导函子与限制函子都保持可对偶化性质。在特征标上，它们重现特征标层通常的诱导函子与限制函子。
\end{ex}

\begin{ex} 若 $X$ 是 $K$-簇，它的特征标层可以如下描述。考虑 $K \times X$ 的闭子概形 $\on{Stab}$，它在 $X$ 的一点上的纤维就是该点的稳定化子。也就是说，$\on{Stab}$ 是作用映射与投影映射到 $X$ 的等化子。于是在平凡 $G$-等变映射
%
\[
\on{Stab} \rightarrow G
\]
%
下，对偶化层的推前就是 $\chi_{\Dmod(X)}$。特别地，由真基变换，$K$ 的任一元素在 $\Dmod(X)$ 上的迹就是其不动点的 Borel--Moore 同调。读者不妨把它与有限群、有限集上的类似断言作比较：在那里，群元素的迹给出不动点的个数。
\end{ex}

\begin{ex} 若 $K=G$ 是约化群，$\sC = \Dmod(\sB)$，则由上例，其特征标层是 $G/G$ 上的 Grothendieck--Springer 层。更一般地，等变自同态的矩阵系数函子（参见公式 \eqref{e:mcequivfunc}）形如
%
\[  \Dmod(B \backslash G / B) \rightarrow \Dmod(G/G).\]
%
这就重现了 Harish-Chandra 变换，或称极限圆变换，即沿对应 $$G/G \leftarrow G/B \rightarrow B \backslash G / B,$$ 的推-拉，其中左边与中间两项是关于伴随作用的商。后者正是 Lusztig 研究特征标层的出发点。
\end{ex}

下面有两条相关的注。

\begin{re}在当前的讨论中，我们把特征标层当作一个直接的构造：输入可对偶化的表示，输出伴随等变的 D-模。

另一种视角是把 $\Dmod(K/K)$ 解释为 $\Dmod(K)$ 的水平迹。确实，对任意幺半 dg 范畴 $\sA$，后者定义为 $\sA \otimes_{\sA \otimes \sA^{\on{op}}} \sA$；或者等价地，定义为 $\sA\mod$ 作为可对偶化 $\DGCat$-线性范畴的 2-范畴迹。若把范畴层级降低一层，即用一个代数 $A$ 代替幺半 dg 范畴，则这段讨论的类比物是把 $A\mod$ 的 Hochschild 同调实现为
%
\[
  A \otimes_{A \otimes A^{\on{op}}} A.
\]
%

在当下情形，利用范畴余不变量与商叠上 D-模之间的对应，有平凡的等同
%
\[
   \Dmod(K) \otimes_{\Dmod(K \times K)} \Dmod(K) \simeq \Dmod(K
   \overset{K \times K} \times K) \simeq \Dmod(K / K).
\]
%
一般的形式体系把可对偶化的 $\sA$-模对应到其水平迹的一个对象，而这就把上面关于特征标层的讨论作为特殊情形包含进来；进一步的讨论参见 \cite{gkrvtoyshtuka} 与 \cite{bznadlerchartheory}。
\end{re}

\begin{re}本节所讨论的理论还有平行的版本：把 $K$ 上的 D-模换成解析的或 \'etale 构造层，参见下文的第 \ref{ss:variants} 小节。当取定义在有限域 $\Fq$ 上的 $K$ 的 \'etale 层时，函数-层对应的一个版本把 $K(\Fq)$ 的模范畴等同于 $\on{Shv}(K)\mod$ 的 Frobenius 迹。这把范畴表示送到通常的表示，并且相应地，把特征标层送到通常的特征标；参见 \cite{Gaitsgoryinventshtukas}。
\end{re}

\subsection{补遗二：生成}

\ssszh 最后一个话题是几条关于生成的好定理。回忆一条基本而有用的常识：Morita 理论告诉我们，给定环 $A$，其模范畴 $A\mod$ 虽由 $A$ 本身紧生成，却也容许许多别的紧生成元 $c$。任取其一，取它的自同态就给出等价
%
\[
  \rmod\End(c)  \simeq   A\mod, \quad \quad M \mapsto M \otimes_{\End(c)} c.
\]
%
换一种说法，给定的 dg 范畴 $\sC$ 可以有许多紧生成元，相应地也有许多把它表示成模范畴的方式。这里想综述两条有用的此类陈述；它们必然发生在高一个范畴层级上，即对 $\Dmod(K)\mod$ 而言。

\ssszh \label{ss:hchbimod}首先，Beraldo \cite{beraldo} 在 Gaitsgory \cite{1affine} 的工作基础上证明了一条非平凡的定理：$\Dmod(K)$ 由 $\fk\mod$ 生成。注意后者的自同态为
%
\begin{align*}
   \Hom_{\Dmod(K)\mod}(\fk\mod, \fk\mod) & \simeq \udHom(\fk\mod, \fk\mod)^K \\
   & \simeq (\fk\mod^\vee \otimes \fk\mod)^K \\ & \simeq (\fk\mod \otimes
   \fk\mod)^K \\ & \simeq (\fk \oplus \fk\mod)^K.
\end{align*}
%
用文字表述：$\fk\mod$ 的等变自同态构成 Harish-Chandra 双模范畴，也就是对角作用可积为群 $K$ 的作用的 $\fk$-双模。\zy{[译者按（上下文推断，无外部依据）：按上面最后的计算，此行应读作 $(\fk \oplus \fk)\mod$，即 $\fk \oplus \fk$ 的模范畴；原文把它印作 $\fk \oplus \fk\mod$，本译本照原文排印。]}

\begin{re} 我们要强调，这条等同是该形式体系一个漂亮而有用的特点。也就是说，Harish-Chandra 双模在表示论中已被研究和使用数十年，人们早已认识到它们具有许多良好的性质，参见 \cite{bege1980} 及其中的参考文献。上面的陈述在精确的意义上说明了这一点：它们恰好就是与范畴 $K$-作用交换的 $\fk\mod$ 的自函子，反之亦然。
\end{re}

于是 Beraldo 的定理给出等价
%
\[
    \Dmod(K)\mod \simeq (\fk \oplus \fk\mod^K)\mod.
\]
%
特别地，对任意 $\Dmod(K)$ 表示 $\sC$，有典范等价
%
\[
    \fk\mod \otimes_{\fk\oplus \fk\mod^K} \sC^{K, w} \simeq \sC.
\]
%
也就是说，回忆等同 $\Dmod(K)^{K,w} \simeq \fk\mod$：一个 $\Dmod(K)$ 表示可以由它的弱不变量重构出来。

\ssszh 对约化群 $G$，我们还想提一条由 Ben-Zvi--Gunningham--Orem \cite{BZO} 给出的最高权理论在当下框架中的漂亮类比。他们证明了：对任意极大幂幺子群 $N$，任意的 $\Dmod(G)$ 表示 $\sC$ 都可以由它的 $N$-不变量重构，甚至只需后者的弱 Cartan 不变量：
%
\[
   \sC^N \quad \text{和} \quad (\sC^N)^{T,w}.
\]
%
过渡到它们的自同态（这些自同态大致是沿 $T$-轨道没有指定单值性的 Hecke 范畴），便得到等价
%
\[
  \Dmod(G)\mod \simeq \Dmod(N \backslash G / N)\mod \simeq \Dmod(N \backslash G
  / N)^{T \times T, w}\mod
\]
%
特别地，还得到等同
%
\[
   \Dmod(G/N) \otimes_{\Dmod(N \backslash G / N)}  \sC^N \simeq
   \Dmod(G/N)^{T,w} \otimes_{\Dmod(N \backslash G / N)^{T \times T,w}}
   (\sC^{N})^{T,w} \simeq \sC.
\]
%


\begin{re} \cite{BZO} 的论证还给出下面这条有用的变体。给定群 $K$ 的子群 $H$，使得 $K$ 对 $H$ 的正规化子的商是真紧的，则典范映射
	%
	\[
	     \Dmod(K/H) \otimes_{\Dmod(H \backslash K / H)} \sC^K
	     \rightarrow \sC,
	\]
	%
虽未必本质满，却总是全忠实的嵌入。由此可以检验「最高权向量」上的范畴等价，而后者往往容易直接分析得多；进一步的讨论与例子参见 \cite{ahc}、\cite{yangmoyprasad}。
\end{re}

\section{环路群的范畴表示}

\label{s:catloop}

\subsection{概览}

\ssszh 到此为止，我们讨论了约化群的 Beilinson--Bernstein 局部化与范畴表示论。

在本综述的余下部分，我们想转向它们的仿射类比，即从约化群过渡到环路群。这方面的故事要丰富得多，正是局部几何朗兰兹对应的主题。

本节将对环路群及其范畴表示作若干介绍性的评论，为第 \ref{s:affbbloc} 节讨论朗兰兹对应作准备。

\subsection{环路群}

\ssszh 与之前一样，仍用 $G$ 记代数闭域 $k$（特征为零）上的连通约化群。

本节的主角不是 $G$，而是它的环路群 $LG$。初步地说，正如 $G$ 由元素在 $k$ 中的某些矩阵构成（取一个忠实表示即可），$LG$ 由同样形状的矩阵构成，只是元素换成了 Laurent 级数域 $k(\!(t)\!)$。更仔细地说，由于我们需要把 $LG$ 理解成代数几何的对象，它是某个群归纳概形，其闭点正如上面所述。

\begin{ex}\label{ex:ga} 当 $G$ 是（非约化的）加法群 $\mathbb{G}_a$ 时，有一个显式的构造，它涵盖了大部分有代表性的想法。此时 $\mathbb{G}_a(k) \simeq k$，即一维仿射空间。如上所述，我们希望 $L\mathbb{G}_a(k) \simeq k(\!(t)\!)$。基本想法是：$k(\!(t)\!)$ 仍是无限维仿射空间；事实上，把 Laurent 级数极点的阶数作界，它就成为一个射有限维仿射空间的递增并。显式地，作为向量空间
%
\[
  k(\!(t)\!) = \varinjlim_n \varprojlim_m t^{-n} \cdot k[[t]] / t^{-n+m}
  \cdot k[[t]],
\]
%
而描述其归纳概形结构用的是同一个公式。也就是说，固定 $n$，得到一个射有限维仿射空间；显式地说，它是可数多个变量的多项式环的谱，而这些变量正是级数 Taylor 系数的函数。让 $n$ 变动，就得到无限维仿射空间沿闭嵌入的递增并，即一个归纳概形。
\end{ex}

\begin{re} 一般地，$LG$ 可以用它的点函子简洁地描述，定义如下。对任意簇 $X$，考虑从仿射簇到集合的函子 $LX$，它在任意 $k$-代数 $R$ 上的 $R$-点取为
	 %
	\[
	     LX(R) = X(R(\!(t)\!) ).
	\]
	%
	这个函子对一般的簇未必表现良好；但若 $X$ 是仿射的，它就由某个归纳无限型的归纳概形表示。取 $X = G$，就得到环路群 $LG$。
\end{re}

给出了形式定义之后，我们仿照例 \ref{ex:ga} 的精神，试着给出更可动手的描述。回忆前几节讨论过的 $G(\Fq)$ 与 $G(k)$ 之间的类比，更一般地，有限 $G(\Fq)$-集与有限维 $G(k)$-簇之间的类比。当下情形的一条基本指导原则，是把 $p$-adic 群 $G(\Fq(\!(t)\!))$ 与环路群 $LG$ 作类比。

\begin{re} 有些作者用 $p$-adic 群一词专指 $G$ 在 $\mathbb{Q}_p$ 或其有限扩张中的点；这里我们略微宽松地用这个词指任意非 Archimedes 局部域（例如 $\Fq(\!(t)\!)$）中的点。 \end{re}


回忆：$p$-adic 群带上其自然拓扑是归纳射有限集。显式地说，取定任一紧开子群 $K$，即满足 $K \cap G(\Fq[[t]])$ 在 $K$ 与 $G(\Fq[[t]])$ 中都是有限指数的子群。那么 $K$ 自然地是有限群的逆向极限，而整个群分解为它的 $K$-陪集的不交并。具体到 $G(\Fq[[t]])$，相应的逆向极限来自幂级数的拓扑，即
%
\begin{equation}
\label{e:invlim}
      G(\Fq[[t]]) = \varprojlim_n G(\Fq[t]/(t^n)).
\end{equation}
%


类似地，可以理解 $LG$ 的紧开子群。具体地说，有弧群 $L^+G$，它是无限型的群概形，$k$-点为 $G(k[[t]])$。它的定义与 \eqref{e:invlim} 类似，是有限维代数群的逆向极限。$LG$ 的子群概形 $K$ 称作紧开的，如果 $K \cap L^+G$ 在 $K$ 与 $L^+G$ 中都是有限余维的。此时 $K$ 自然是代数群的逆向极限，而陪集空间 $LG/K$ 是归纳有限型的归纳概形。我们要强调：在有限集与有限维簇的类比之下，这对应于 $p$-adic 群情形中陪集的可数离散并。



\subsection{环路群上的 D-模与范畴作用}

\ssszh 接下来想讨论 $LG$ 上 D-模的卷积代数。由于这是无限维的对象，需要多加小心。

它的定义可以这样得到。回忆我们把 $\Dmod(G)$ 的模看成类似于 $G(\Fq)$ 的表示；特别地，D-模的卷积代数本身就类比于群代数 $\Fun(G(\Fq))$。

我们同样希望 $\Dmod(LG)$ 的模表现得像 $p$-adic 群 $G(\Fqt)$ 的表示。后者是拓扑群，相应地就要用光滑表示。\footnote{显式地说，复向量空间 $V$ 上的一个表示是光滑的，如果作用映射 $G(\Fqt) \times V \rightarrow V$ 关于 $V$ 上的离散拓扑连续，即每个向量的稳定化子都是开子群。对不熟悉 $p$-adic 表示论的读者，这里的要点是：这个条件对应李群表示中的类似条件，那里要求作用在流形的意义下适当光滑。不过在 $p$-adic 情形中，群上的拓扑与向量空间上的拓扑（不严格地说）是「互素的」。实际上，研究自守形式时自然出现的正是光滑表示。}因此，$G(\Fqt)$ 的模并不等同于其群代数的模，而应当是其 Hecke 代数
$$\on{Dist}(G(\Fqt)) $$
即紧支分布代数的模。具体地说，它可以写成余极限
%
\begin{equation} \label{e:padicgroups}
\on{Dist}(G(\Fqt))\simeq \varinjlim_K  \on{Dist}( K \backslash G(\Fqt) / K),
\end{equation}
%
其中 $K$ 取遍紧开子群，而每个 $\on{Dist}( K \backslash G(\Fqt) / K)$ 不过是双陪集上 delta 分布的有限线性组合，即把函数沿某个双陪集 $KgK$（关于某个 Haar 测度）积分的分布。


有了这些准备，环路群上的 D-模范畴定义为
%
\[
\Dmod(LG) \simeq \varinjlim \Dmod(K \backslash LG / K),
\]
%
其中 $K$ 取遍紧开子群，余极限在 $\DGCat$ 中沿 $*$-拉回函子取。剩下来要描述的是，对固定的 $K$，范畴 $\Dmod(K\backslash LG/K)$ 是什么。回忆 $LG/K$ 是归纳有限型的归纳概形，故还可以把它写成递增并
%
\[
   LG/K = \bigcup_i X_i,
\]
%
其中每个 $X_i$ 是有限型的闭子概形，并且是 $K$-不变的。由于 $K$ 通过一个有限维商 $Q_i$（其核为射幂幺群）作用在每个 $X_i$ 上，有
%
\[
   \Dmod(X_i)^K \simeq \Dmod(X_i)^{Q_i} \simeq \Dmod(Q_i \backslash X_i),
\]
%
也就是说，我们面对的是通常的有限维叠上的 D-模。于是令
%
\[
  \Dmod(K \backslash LG/K) \simeq \varinjlim \Dmod( Q_i \backslash X_i).
\]
%

这里的要点是：粗略地说，$\Dmod(LG)$ 中的对象类似于 $G(\Fqt)$ 上的紧支分布。更精确地说，由例 \ref{e:colimitcats}，$\Dmod(LG)$ 是紧生成的，其紧对象由某个 $K$ 的 $LG/K$ 的子簇上凝聚 D-模有界复形的 $*$-扩张给出。


这个构造的确切细节对我们并不那么重要，不过下面给出两条帮助定位的注。

\begin{re}第一，可以把上述余极限 \eqref{e:padicgroups} 中所有可能的 $K$ 换成任一共尾子集。一个具体的取法是 $L^+G$ 的同余子群 $K_n$（$n \geqslant 1$），我们只需在 $k$-点的层面用正合列来描述它
	%
	\[
	      1 \rightarrow K_n \rightarrow L^+G = G(k[[t]]) \rightarrow
	      G(k[[t]]/(t^n)) \rightarrow 1.
	\]
	%

注意 $K_1$ 是 $L^+G$ 的射幂幺根，由此可知任意充分小的紧开子群 $K$ 都是射幂幺的。特别地，对两个这样的射幂幺子群 $K' \subset K$，拉回
%
\[
    \Dmod(K \backslash LG / K) \rightarrow \Dmod(K' \backslash LG / K')
\]
%
是全忠实的嵌入，参见注 \ref{r:invariantstwice}。
\end{re}



\begin{re}上述构造看起来或许有些凑合，但应当指出：它是为更大一类无限维簇（称作平缓归纳概形）给出 D-模定义的做法中的一个特例 \cite{rdm}。

事实上，这一无限维设定中有一处基本的微妙之处：D-模范畴有两个对偶的版本。一个带有 $*$-推前的函子性，另一个带有 $!$-拉回的函子性。上面的表述对应于带有 $*$-推前的那个。这正是 $G(\Fqt)$ 上函数与分布之区别的类比。与此相关，把维数主丛（dimension torsor）平凡化（即固定一个紧开子群）之后，$LG$ 的这两个对偶范畴可以等同起来。
\end{re}



\ssszh 与有限维情形一样，$\Dmod(LG)$ 在卷积下是幺半 dg 范畴。特别地，如前可以定义它的模范畴
%
\[
      \Dmod(LG)\mod.
\]
%

下面是一些表示的基本例子。

\begin{ex} 若 $X$ 是被 $LG$ 作用的平缓归纳概形，其 D-模范畴 $\Dmod(X)$ 通过卷积带有 $\Dmod(LG)$ 的典范作用。
\end{ex}


\begin{ex}\label{ex:kmmods} 给定环路群的中心扩张
	\[
	1 \rightarrow \mathbb{G}_m \rightarrow \widetilde{LG} \rightarrow LG
	\rightarrow 1,
	\]
	%
	过渡到李代数，就得到环路代数的中心扩张，即仿射李代数
	\[
	 0 \rightarrow k \cdot 1 \rightarrow \gk \rightarrow L \fg \rightarrow
	 0.
	\]
	%
	光滑模的导出范畴
	$
	    \gk\mod
	$
	带有 $\Dmod(LG)$ 的典范作用；这来自 $\widetilde{LG}$ 在 $\widehat{\fg}$ 上的伴随作用通过 $LG$ 分解这一事实。按定义，相应阿贝尔范畴的对象是 $\gk$ 的表示，其中 $1$ 作用为恒等，且每个向量都被某个紧开子代数零化。依照 Frenkel--Gaitsgory \cite{fglocaff}，完整的 dg 范畴是对朴素无界导出范畴作小幅修正，使得 Verma 模、更一般地紧开子代数有限维表示的诱导，都是紧生成元。

	这些细节对我们并不重要，只需指出：这一修正的存在与用处，源于该阿贝尔范畴具有无限的上同调维数。与此相关，重整化后的导出范畴带有 $t$-结构，并在有下界部分与通常的导出范畴典范地等同。也就是说，重整化只在同调次数 $-\infty$ 处带来差别。类似地，范畴 $\Dmod(LG)$ 也带有 $t$-结构，它在整体平移的意义下是典范的，\footnote{这对应 $G(\Fqt)$ 上 Haar 测度在相差一个常数因子意义下的唯一性。}并具有这样的性质：除同调次数 $-\infty$ 处的问题外，$\Dmod(LG)$ 与其心的朴素无界导出范畴一致。
\end{ex}

\section{局部几何朗兰兹与仿射 Beilinson--Bernstein 局部化}

\label{s:affbbloc}

\subsection{概览}

\ssszh 前几节讨论了约化李代数的局部化理论，并说明了它与约化群的范畴表示论之间的关系。第 \ref{s:catloop} 节中，我们遇到了环路群及其范畴表示。

在这最后一节里，我们想概述局部几何朗兰兹中的一些基本猜想与结果；它们可以理解为对环路群的范畴表示作参数化。最后我们会回到原点，遇到环路群上 Beilinson--Bernstein 局部化的类比，它由 Beilinson--Drinfeld、Frenkel--Gaitsgory 与 Raskin--Yang 发展起来。

\subsection{对应的朴素表述}



\ssszh 对 $p$-adic 群 $G(\Fqt)$，局部朗兰兹对应（猜想地）用 Galois 理论的数据来参数化它在 $k$ 向量空间上的不可约表示；初步地说，这些数据是从 $\Fqt$ 的绝对 Galois 群到朗兰兹对偶 $\check{G}(k)$ 的同态。


本小节想综述：对 $\Dmod(LG)\mod$ 的类似参数化，人们期待什么、又已知什么。



\ssszh 回忆 $k(\!(t)\!)$ 由形式穿孔圆盘 $D^\times$ 上的函数构成。按照 Galois 表示与局部系统之间的类比（参见附录 \ref{s:frmfun}），当前设定中与 $\cG$-值 Galois 表示相对应的，是平坦的 $\check{G}$-联络，即 $D^\times$ 上的 $\cG$ de Rham 局部系统。

具体地说，它们由联络模去规范等价给出，即
%
\[
    \LS \simeq \{ d + \cfg(\!(t)\!)dt \} / \cG(k(\!(t)\!)).
\]
%
同一个公式还把 $\LS$ 定义为代数几何的对象：取 $L\cG$ 在联络形式 $d + \cfg(\!(t)\!)dt$ 的归纳概形上作用的商预叠，而后者与例 \ref{ex:ga} 类似地定义；进一步的讨论参见 \cite{raskinspectraldecomp}。



\ssszh \label{s:infdiscuss} 下面试着以略为启发式的方式给出局部几何朗兰兹对应的陈述。初步地说，可以期望不可约的 $\Dmod(LG)$ 表示与 $\LS$ 的点一一对应。也就是说，对每个 $\cG$-联络 $\sigma$，应当能典范地配上 $LG$ 的一个不可约范畴表示 $\sC_{\sigma}$。

由于 $\Dmod(LG)\mod$ 的一般对象 $\sC$ 应当分解为不可约表示的直积分，它应当定义 $\LS$ 上相应的一族范畴，而该族在每个点的纤维衡量 $\sC_{\sigma}$ 在直积分中相应的重数。

由于不可约性是阿贝尔范畴中对象的特征，这段讨论很难精确表述。不过，对于空间上的一族范畴，有一个好的概念，即拟凝聚的范畴层，下面就来综述。

作一些修正（下文公式 \eqref{e:loclang} 会适时遇到）之后，这就给出局部几何朗兰兹对应的陈述。也就是说，猜想是：环路群的全部范畴表示与局部系统模空间上的范畴层之间，应当有 $(\infty,2)$-范畴的等价，即
%
\begin{equation}
\label{e:loclangnaive}       \mathbb{L}: \Dmod(LG)\mod \simeq
\on{2-QCoh}(\LS).
\end{equation}
%



下面给出一条帮助定位的基本注。所述 2-范畴的等价并不是说：对给定的范畴表示 $\sC$，$\sC$ 与 $\mathbb{L}(\sC)$ 的底层 dg 范畴等价。而是说：对任意一对对象 $\sC$ 与 $\sD$，有等价
%
\[
         \Hom_{\Dmod(LG)\mod}(\sC, \sD) \simeq
         \Hom_{\TwoCoh(\LS)}(\mathbb{L}(\sC), \mathbb{L}(\sD)),
\]
%
并且当 $\sC$、$\sD$ 变动时，这些等价与复合相容。

\begin{re}  作为玩具模型，通常代数的 Morita 等价
    \[
        A\mod \simeq B\mod
    \]
	%
	匹配的并不是模的底层向量空间，而只是它们之间的同态。具体地说，对有限维 $k$ 向量空间 $V$，可以考察 Morita 等价
	%
	\[
	   V \otimes -  :\Vect \simeq  \End(V)\mod.
	\]
	%
	它把 $k$ 与 $V$ 互换，而这两者当然不同构；但在自同态上确实有
	%
	\[
	     \Hom_{\Vect}(k, k) \simeq k \simeq \Hom_{\End(V)}(V,V).
	\]
	%
		\end{re}


在转向修正之前，我们想先讨论范畴层，以及 $\mathbb{L}$ 应当满足的一些基本相容性。

\ssszh   要讨论范畴层，先从仿射簇 $X = \on{Spec }O_X$ 的情形开始。此时，$X$ 上的拟凝聚范畴层不过是一个 $O_X$-线性的 dg 范畴（而非仅仅是 $k$-线性的）。由于 $O_X$-线性等同于可以张量 $O_X$-模的复形，等价地有
%
\[
	  \TwoCoh(X) \simeq \QCoh(X)\mod.
\]
%
特别地，$\TwoCoh(\on{pt})$ 就是 $\DGCat$；并且对任意闭点 $\on{pt} \rightarrow X$，有拉回
%
\[
	  \Vect \otimes_{\QCoh(X)}  -: \TwoCoh(X) \rightarrow \DGCat,
\]
%
也就是说，可以谈论范畴层的纤维，它们就是通常的 dg 范畴。这样，范畴层就被看成 $X$ 上的一族范畴。



\begin{re}\label{r:pushpull2coh} 更一般地，给定任意映射 $f: X \rightarrow Y$，拉回定义了幺半函子 $\QCoh(Y) \rightarrow \QCoh(X)$，特别地，得到伴随的诱导函子与限制函子
	%
	\[
	    f^*: \TwoCoh(Y) \leftrightarrows \TwoCoh(X): f_*.
	\]
	%
\end{re}





下面是一个范畴层的基本例子。

\begin{ex}给定簇的映射 $\tilde{X} \rightarrow X$，$\QCoh(\tilde{X})$ 是 $\QCoh(X)$-代数，特别地是 $X$ 上的范畴层。对 $X$ 的任意点 $x$，由例 \ref{e:fibprodvar}，该范畴层的纤维不过是实际纤维上的拟凝聚层，即
%
\[
	  i_x^* \QCoh(\tilde{X}) \simeq \QCoh( \tilde{X} \times_X x).
\]
%
\end{ex}


\ssszh 对一般的 $X$（不必仿射），定义像拟凝聚层的定义那样通过粘接给出。细节对我们并不那么重要，但应当提到 Gaitsgory 在 \cite{1affine} 中建立的下述漂亮现象。即：对一般的 $X$，按定义是逐个仿射地拼贴范畴层；另一方面，又有拟凝聚层的整体幺半范畴，从而有其左模范畴 $\QCoh(X)\mod$。于是可以问这两者之间的关系。与通常的拟凝聚层一样，有整体截面函子
%
\begin{equation} \label{e:globsectshvcat}
 \Gamma: \TwoCoh(X) \rightarrow \QCoh(X)\mod.
\end{equation}
%
对通常的层，类似的映射
    %
    \[
    \Gamma: \QCoh(X) \rightarrow \Gamma(X, \OO_X)\mod
    \]
	%
	在 $X$ 非仿射时极少是等价。\footnote{对「极少」一词有疑问的读者：要点在于，若考虑整体函数的完整导出代数，则对非仿射簇（例如仿射簇的开子簇），$\Gamma$ 仍可以是同构。}
相反，对多数有限维的 $X$（例如任意有限型的代数叠），函子 \eqref{e:globsectshvcat} 都是等价；这一现象称作 1-仿射性。


本领域的一条基本猜想（属于 Gaitsgory）是：$\LS$ 也是 1-仿射的，参见 \cite{raskinspectraldecomp}。这是一条微妙的断言，因为对归纳概形或无限型群概形的分类空间，1-仿射性一般并不成立。所以从某种意义上说，需要这两类问题对 $\LS$ 恰好「抵消」。Raskin 的一项重要工作 \cite{raskinspectraldecomp} 部分地处理了这一点，他证明了：对 $X = \LS$，\eqref{e:globsectshvcat} 是全忠实的。

\subsection{对应的性质 I}

\ssszh 定义方面的讨论既已交代完毕，现在转入正题，看看 $\mathbb{L}$ 应当是什么样子。更精确地说，我们列出两侧各自的一些基本对象，它们预期会互相对应。再加上与抛物诱导的相容性，这预期能在很大程度上确定该等价。

讨论的结构如下。我们在图 \ref{f:table} 中列出了互相对应的范畴。对图中每一行，我们会花一点篇幅描述两侧出现的对象，并说明为什么应当期待它们对应。最后我们看看，在这些情形中匹配缠绕算子意味着什么；特别是，会遇到前面许诺过的对 \eqref{e:loclangnaive} 的修正。



\begin{table}[htbp]
\centering
\begin{tabular}{|c|c|}
\hline
$LG$ 的范畴表示 & $\LS$ 上的范畴层 \\
\hline
$\Dmod(\on{Gr_G})$ & $\QCoh( \on{pt}/\cG)$ \\
$\Dmod(\Fl^{\on{aff}}_G)$ & $\QCoh( \widetilde{{\sN}}/\cG)$ \\
$\Dmod(LG/LN, \psi)$ & $\QCoh(\LS)$ \\
$\gkc\mod$ & $\QCoh(\Op)$ \\
\hline
\end{tabular}
\caption{局部几何朗兰兹下（猜想地）互相交换的一些对象。}
\label{f:table}
\end{table}


\ssszh 对图 \ref{f:table} 中的前两个例子，回忆 $p$-adic 表示论的一些基本特征会有助于理解。由光滑性，$p$-adic 群的任意表示 $\pi$ 对所有充分小的紧开子群 $K$ 都含非平凡的不变量 $\pi^K$。对这样的 $K$，若 $\pi$ 不可约，则不动向量 $\pi^K$ 是 Hecke 代数 $\on{Dist}(K \backslash G(\Fqt) / K)$ 的不可约表示，后者通过卷积自然地作用于它。这样便得到 $G(\Fqt)$ 的含 $K$-不变向量的不可约表示与该 Hecke 代数的不可约表示之间的双射。

当子群 $K$ 缩小时，相应的 Hecke 代数变大，这给出表示复杂度的一个基本度量。即：对给定的不可约 $\pi$，可以问使它有不平凡不变量的最大 $K$ 是什么；从某种意义上说，$K$ 越大，$\pi$ 越不复杂。


在局部朗兰兹的另一侧，$\Fqt$ 的绝对 Galois 群（更准确地说，其 Weil--Deligne 变体）带有典范的滤过，由递降的正规子群给出，称作分歧滤过。对给定的 $\cG(k)$-值 Galois 表示，可以问它的作用经过该滤过的哪个商，以此度量其复杂度。一条基本而引人注目的组织原则是：这两个复杂度度量在局部朗兰兹下互换。\footnote{更精确地说，表示的深度（由 Moy--Prasad 定义的某个非负有理数 \cite{moyprasad}）应当等于相应 Galois 表示的斜率，后者由上编号的分歧滤过度量。}


\subsection{对应的性质 II：非分歧表示}

\ssszh 由上述讨论可知，$p$-adic 群最容易理解的表示，应当是那些关于极大紧开子群 $G(\Fq[[t]])$ 有不动向量的表示。\footnote{应当指出，对 $GL_n$ 以外的群，还存在其他不共轭的极大紧开子群。}类似地，最简单的 Galois 表示应当是那些作用经过分歧滤过最后一个商的表示，即可以延拓到非穿孔圆盘上的 $\check{G}$-局部系统。具体地说，这些正是经过剩余域的 Galois 群分解的表示
%
\begin{equation} \label{e:galunramified}
 1 \rightarrow I \rightarrow \on{Gal}( \Fqt^s / \Fqt) \rightarrow
 \on{Gal}(\overline{\mathbb{F}}_q /
 \Fq) \simeq \widehat{\mathbb{Z}} \rightarrow 1,
\end{equation}
%
其中上标「$s$」表示可分闭包。

由于实际使用的是 Weil 群（按定义它把 $\widehat{\mathbb{Z}}$ 换成 $\mathbb{Z}$），非分歧 Galois 表示在同构意义下由 $\check{G}$ 中共轭类中的元素给出。也就是说，我们追踪 $1 \in \mathbb{Z}$ 在映射 $\mathbb{Z} \rightarrow \cG$ 下的像。在这个最简单的例子中，局部朗兰兹归结为球 Hecke 代数与 $\cG$ 的特征标环之间的 Satake 同构，即
%
\[    \on{Dist}(G(\Fq[[t]]) \backslash G(\Fqt) / G(\Fq[[t]])) \simeq
\on{Fun}(\cG(k) / \cG(k)),
\]
%
其中右端表示 $\cG$ 上伴随不变的多项式函数，把 $\cG$ 看作 $\on{Spec} k$ 上的簇 \cite{satake63}。过渡到特征标，就得到：$G(\Fqt)$ 的不可约非分歧表示由 $\cG(k)$ 中的（半单）共轭类参数化，正如所期望的那样。

\begin{re} 对初次接触这些想法的读者，我们要强调：甚至上面出现的 Hecke 代数的交换性都不是先验显然的。这恰好说明，过渡到关于一个大紧子群的双不变分布，会多么戏剧性地简化问题。\end{re}

\ssszh 为几何版本作准备，我们应当指出：每个不可约非分歧表示都唯一地（相差一个标量）嵌入万有非分歧表示
%
\begin{equation} \label{e:fungr}
	    \Fun( G(\Fqt) /G(\Fq[[t]])).
\end{equation}
%
而且，要从万有情形提取出每一个，只需取 \eqref{e:fungr} 中在球 Hecke 代数的右卷积下按相应特征标变换的子模。



最后，我们讨论上述内容在局部几何朗兰兹中的类比。在自守侧，把仿射 Grassmann 簇上的 D-模范畴
%
\[
	   \Dmod(\Gr_G) = \Dmod(LG/L^+G)
\]
%
当作上述万有非分歧表示的替代物；读者现在应当已经熟悉这种类比了。在谱侧，穿孔圆盘上任何能延拓到整个圆盘的 de Rham 局部系统都是可平凡化的，于是得到叠点
%
\[
	    \on{pt}/\cG \rightarrow \LS
\]
%
相应地把 $\QCoh( \on{pt}/\cG)$ 看作 $\LS$ 上的范畴层。


还应当指出：不仅有万有情形的类比，也有单个不可约非分歧表示的类比。更详细的讨论留到第 \ref{ss:locaff} 节。

\begin{re}应当指出，非分歧情形虽是最基本的，但它在算术朗兰兹与几何朗兰兹中都完全根本。即在整体朗兰兹对应中（这里不详述），adelic 群 $G(\mathbb{A}_F)$ 的任意不可约光滑表示都分解为各位置处 $p$-adic 群 $G(F_x)$ 的不可约表示的限制张量积，而除有限多个位置外都是非分歧的。为了哪怕理解这样的无限张量积，都必须实质性地用到非分歧表示有典范的生成直线，即其球不变量。在整体几何理论中，迄今已知的几乎一切都只是关于处处非分歧（或许还有野分歧）自守表示的类比。
\end{re}

\subsection{对应的性质 III：非分歧主系列}

\label{s:unramps}
\ssszh 有了 $p$-adic 群的万有表示，可以如下造出更多表示。如同表示论的其他部分，有抛物诱导函子，可以从较小群（例如 Cartan 群 $T$）的表示出发构造 $G$ 的表示。把非分歧表示代入这一构造，就得到一族重要的模，即非分歧主系列。它们的单子商恰是那些关于紧开 Iwahori 子群 $I$ 有不平凡不变量的单 $G(\Fqt)$-模；这里把 $I$ 理解为 $G(\Fq[[t]])$ 在投影 $G(\Fq[[t]]) \rightarrow G(\Fq)$ 下对 $G(\Fq)$ 的 Borel 子群的原像。\footnote{这条定理常被归于 Borel 与 Matsumoto，我们也听到有人把它归于 Bernstein 或 Casselman；参见 \cite{cass80spherical} 的命题 2.4 及其中的参考文献。}也就是说，把抛物诱导作用于非分歧表示，得到的是下一类最简单的 $p$-adic 群表示。

用谱侧的语言说，$G(\Fqt)$ 的这类不可约表示对应于野分歧、幂幺单值（unipotent monodromy）的 Galois 表示。为解释这一点，从序列 \eqref{e:galunramified} 回忆惯性子群
%
\[
   I \simeq \on{Gal}( \overline{\mathbb{F}}_q(\!(t)\!)^s / \overline{\mathbb{F}}_q(\!(t)\!)),
\]
%
其中 $\overline{\mathbb{F}}_q$ 表示 $\Fq$ 的代数闭包。分歧滤过的下一个商把它分裂为
%
\[
	    1 \rightarrow I_{\on{wild}} \rightarrow I \rightarrow I_{\on{tame}} \rightarrow 1.
\]
%
对野惯性 $I_{\on{tame}}$，基本想法是：它的行为与特征零的对应物一样，即与 $k$ 上形式穿孔圆盘的覆叠一样。这些覆叠的分类方式与拓扑中圆盘的覆叠空间完全一样，即对正整数 $n$，添加 $t$ 的一个 $n$ 次根。在正特征情形，它们可分当且仅当 $n$ 与 $q$ 的特征互素；而这样的覆叠穷尽了所有次数与 $q$ 互素的 Galois 扩张。显式地，可以取同构
%
\[
   I_{\on{tame}} \simeq \underset{\ell} \Pi \hspace{.7mm} \mathbb{Z}_\ell,
\]
%
其中 $\ell$ 取遍不等于 $\Fq$ 的特征的素数。Galois 表示称作野分歧的，如果它经过对 $I_{\on{wild}}$ 的商分解；进而称作幂幺单值的，如果 $I_{\on{tame}}$ 的任意元素（或等价地，某个拓扑生成元）都被送到 $\cG$ 的幂幺元素。



\ssszh 为几何讨论作定位，我们稍为详细地回顾一下 $p$-adic 表示的相应抛物诱导如何进行。给定 Borel 子群 $B$，其 Levi 商是抽象 Cartan 群 $T$，于是在有理点上得到对应
%
\[
	    T(\Fqt) \leftarrow B(\Fqt) \rightarrow G(\Fqt),
\]
%
其中向左的箭头是满射，核为幂幺根 $N(\Fqt)$ 的有理点；向右的箭头是闭嵌入。由此得到抛物诱导函子
%
\begin{equation} \label{e:parindpadic}
	    T(\Fqt)\mod \xrightarrow{\on{Res}} B(\Fqt)\mod \xrightarrow{\on{Ind}}
	    G(\Fqt)\mod.
\end{equation}
%
具体地说，把 $T(\Fqt)$ 的光滑表示 $\pi$ 通过让幂幺根平凡作用而膨胀成 $B(\Fqt)$ 的表示，再诱导到 $G(\Fqt)$。这一诱导与通常一样，具体地由旗流形上相应向量丛的光滑截面给出，即 $G(\Fqt) \overset{B(\Fqt)} \times \pi \rightarrow G(\Fqt)/B(\Fqt)$。



特别地，由此可知：$T$ 的任意非分歧特征标的抛物诱导都嵌入万有非分歧表示的抛物诱导，即万有主系列模
%
\[
	    \Fun\big( G(\Fqt) / T(\Fq[[t]]) \cdot N(\Fqt)\big).
\]
%

\begin{re} 对尚未具体接触过抛物诱导的读者，指出下面这一点或许有助于定位：用最高权分类约化群的不可约代数表示，正是抛物诱导的一个例子。

确实，对上文那样的代数群 $G,B,T,N$，在任意域上的代数表示范畴之间有伴随对
%
\[
	    \on{Rep}(G) \overset{\on{Res}}{\underset{\on{Ind}}{\rightleftarrows}}
	    \on{Rep}(B)
	    \overset{\on{Coinv}_N}{\underset{\on{Res}}{\rightleftarrows}}
	    \on{Rep}(T),
\]
%
其中 $\on{Res}$ 表示限制，$\on{Coinv}_N$ 表示 $N$ 的余不变量函子，$\on{Ind}$ 表示余诱导函子。把它们接起来，就得到抛物限制函子：它把 $G$-模送到其 $N$-余不变量，看作 $T$-模；还得到它的右伴随，即抛物余诱导函子：它把 $T$-模以 $N$ 的平凡作用延拓为 $B$-模，再余诱导到 $G$。

有了这些，最高权定理（可以表述为每个不可约 $G$-模都有一条 $N$-余不变量直线）特别地说明：每个不可约 $G$-模都作为某个抛物余诱导模的子模出现。此外，在特征零时，由完全可约性可知每个不可约模都与某个抛物余诱导模重合。借助「特征标的余诱导即线丛截面」这一显式模型，这就重新得到关于 $G/B$ 上线丛截面的 Borel--Weil 定理。 \end{re}


\ssszh 转向几何情形。此时自守侧类似地是半无限旗流形上的 D-模，即
		%
\[
	  \Dmod( LG / L^+ T \cdot LN),
\]
%
如前所述，它可以由 $\Dmod(\Gr_T)$ 通过一个类似于 \eqref{e:parindpadic} 的适当抛物诱导函子构造出来。

\begin{re} 关于这个范畴有一条技术性的评论。目前还不知道如何把空间 $LG/L^+T \cdot LN$ 当作代数几何的对象来处理。不过，它的 D-模范畴可以通过强制光滑下降来定义，即取 $\Dmod(LG)$ 的 $\Dmod(L^+T \cdot LN)$-不变量。已知（由 Raskin 的一个论证 \cite{raskin2ndadjoint}）这典范地等同于余不变量，也等同于 Iwahori 子群的（余）不变量，即仿射旗簇上的 D-模。这与 $p$-adic 情形中 Borel--Matsumoto 的结果平行，而且实际上是用类似方法证明的。
\end{re}



\ssszh 在谱侧，再次施加抛物诱导函子，即沿对应
 %
 \[
	    \on{LocSys}_{\check{T}} \leftarrow \on{LocSys}_{\check{B}} \rightarrow
	    \LS
 \]
 %
 对 $\QCoh(\on{pt}/\check{T})$ 作推-拉。这里指的是范畴层的拉回与推前，如注 \ref{r:pushpull2coh} 所述。\footnote{更仔细地说，前面许诺的对公式 \eqref{e:loclangnaive} 的修正将涉及从拟凝聚范畴层过渡到某些归纳凝聚范畴层，但这一修正不影响本计算。}




	  下面逐步说明如何计算。第一步是拉回，注意有笛卡儿图
	  %
	  \[
	  \xymatrix{ \on{LocSys}_{\check{N}} \ar[r] \ar[d] &
	  \on{LocSys}_{\check{B}} \ar[d] \\ \on{pt} \ar[r] &
	  \on{LocSys}_{\check{T}}.}
	  \]
	  %
	  这不过是说：一个 $\check{B}$ 局部系统，若其相伴的 $\check{T}$ 局部系统已平凡化，就等同于一个 $\check{N}$ 局部系统。更通俗地说，若转移矩阵的对角线上全是 $1$，那它就不是 $\check{B}$-丛，而是 $\check{N}$-丛！由 $\check{N}$ 的幂幺性，穿孔圆盘上 $\check{N}$-局部系统的模空间就是 $\check{\fn} / \check{N}$，同构由映射
	  %
	  \begin{equation} \label{e:rsmap}
	    \check{\fn} / \check{N} \xrightarrow{\sim}
	    \on{LocSys}_{\check{N}},
	    \quad  X \mapsto d +
	    X \cdot
	    \frac{dt}{t}.
	  \end{equation}
	  %
	 给出。由这段讨论可知，拉回为
	 %
	 \begin{equation}
	 \begin{split}
	 \QCoh(\on{LocSys}_{\check{B}}) \otimes_{\QCoh(\on{LocSys}_{\check{T}})}
	 \QCoh(\on{pt}/\check{T}) &\simeq \QCoh( \on{LocSys}_{P}  \times_{
	 \on{LocSys}_{\check{T}}} \on{pt}/\check{T}  ) \\
	 &\simeq \QCoh(\check{\fn} / \check{B} ).
	 \end{split}
	 \end{equation}
	 %


最后，推前再次给出 $\QCoh(\check{\fn} / \check{B} )$，通过投影
%
\begin{equation} \label{e:projection}
	     \check{\fn} / \check{B} \rightarrow \on{LocSys}_{\check{B}}
	     \rightarrow \LS.
\end{equation}
%
把它看作 $\LS$ 上的范畴层。把它纯粹用 $\cG$ 重写一遍会更清楚。为此，可以平凡地用旗簇的余切丛重写我们的叠，即
%
\[
	      \check{\fn} / \check{B} \simeq (\cG \overset{ \check{B}}
	      \times \check{\fn}) / \cG =:
	      \widetilde{\sN} / \cG.
\]
%
具体地说，$\widetilde{\sN}$ 是数对 $({\check{B}}', X)$ 的模空间，其中 ${\check{B}}'$ 是 Borel 子群，$X$ 是其幂幺根的元素。于是投影 \eqref{e:projection} 可以用 Springer 消解来写，即复合
%
\[
	    \widetilde{\sN}/ \cG \rightarrow \cfg / \cG \rightarrow \LS.
\]
%
这里第一个映射把数对 $(\check{B}', X)$ 送到 $X$，第二个映射的定义与 \eqref{e:rsmap} 类似（但离同构差得很远）。




\begin{re} 下文第 \ref{s:2realz} 节将会综述：$G$ 的主系列表示与 $\cG$ 的 Springer 消解上凝聚层之间的关系，最早出现在 $p$-adic 理论中。
\end{re}



\ssszh 前面两个例子，即图 \ref{f:table} 的前两行，涉及的是低分歧的对象。局部几何朗兰兹对应的这些情形虽属最基本，却已是严肃的结果。

一方面，如前所述，它们为（猜想中的）整体几何朗兰兹对应中最简单的情形提供了足够的局部输入；那里研究的是整体曲线上 $G$-丛模空间上的 D-模，可能还在有限多个点处带有抛物归约。



另一方面，它们在通常的表示论中有重要应用。一个显著的例子是对临界水平仿射李代数表示论的含义，我们将在下文第 \ref{ss:locaff} 节更详细地讨论。另一项关键应用是在过渡到模系数层之后，\footnote{严格说来，这些更自然地属于局部几何朗兰兹的 \'etale 版本或 Betti 版本，参见下文第 \ref{ss:variants} 节，因为它们有系数可变的自然类比。}对约化群的模表示论的应用；参见 \cite{finkelbergmirkovicconjecture}、\cite{mirkovic-vilonen}、\cite{acharrider}、\cite{richewilliamsontilting}、\cite{bezgaitsmiriwahotiwhittaker}、\cite{bezricheriderregcentralizer}、\cite{bezricheprinblock} 及其中的参考文献，以及综述 \cite{willimasonsurvey1}、\cite{williamsonsurvey2}。


鉴于这些进展的重要性，相当引人注目的是：它们只触及整个猜想的表层，即只等同了两侧非常小的子范畴。确实，用谱侧的话说，前述情形只涉及 $\TwoCoh(\LS)$ 中支撑在具有正则奇点与幂幺单值性的局部系统的形式邻域上的全子范畴；在整个穷尽 $\LS$ 的无限递增的轨迹序列中，这是最小的一块。



相比之下，下面要讨论的两个对应例子，涉及任意分歧的范畴环路群表示与 $\LS$ 上的范畴层。因此，这些是相当壮观的猜想，属于当下表示论这一部分中最深刻的结果之列；要把它们解决，很可能需要超出几何表示论标准工具箱的新想法。

\subsection{对应的性质 IV：Whittaker 模型}

\ssszh 表示论的许多情形中，例如约化群的代数表示或 $(\fg,B)$-模，不可约表示都带有\nobreak 典范的生成直线，即最高权向量。

与之相对，对 $p$-adic 群 $G(\Fqt)$，大多数不可约模并没有典范的生成直线。粗略地说，原因在于 $\on{Dist}(K\backslash G/K)$ 对一切充分小的 $K$ 都非交换。

\begin{re} 对非分歧表示，存在典范的生成直线，即 $G(\Fqt)$-不变量。这条基本奇迹也就是球 Hecke 代数的交换性，在本领域至关重要。例如，它给出自守表示的张量积分解，并把非分歧自守表示的研究化归为 Hecke 特征函数与特征层的研究。\end{re}

不过，如果我们所要求的典范直线不在 $\pi$ 内部，而是 $\pi$ 的一个商，情形就好一些。确切地说，存在一个从 $G(\Fqt)\mod$ 到向量空间的典范函子，它把每个不可约表示送到一条直线或零。这就是 Whittaker 余不变量函子，可以这样理解。

为了逼近定义，不妨先考虑最高权理论一个看似直接的类比，即 $N(\Fqt)$-余不变量的 Jacquet 函子
\[
\on{Jac}: G(\Fqt)\mod \rightarrow \on{Vect}, \quad \pi \mapsto
\pi_{N(\Fqt)}.
\]
它虽然很有用，却把许多有趣的不可约表示送到零\footnote{这一现象大致就是尖点表示的存在性，即这样的单 $G(\Fqt)$-模：它不出现在抛物诱导模之内。我们要强调，这是 $p$-adic 情形以及较简单的有限李型群情形的基本事实，而约化群的代数表示论中没有这种现象。}，把另一些送到维数大于 $1$ 的向量空间。

为作修正，我们用一个通用特征标来扭转这个构造。确切地说，以 $i \in I$ 标记单根，考虑复合
\begin{equation} \label{e:whitchar}
\psi:  N(\Fqt) \rightarrow (N/[N,N])(\Fqt) \simeq \underset{i} \Pi
\hspace{.2mm}
\Fqt \xrightarrow{\on{sum}} \Fqt
\xrightarrow{\on{res}} \Fq,
\end{equation}
其中 $\on{sum}$ 是 Laurent 级数的加法，$\on{res}$ 把 Laurent 级数送到它的留数，即 $t^{-1}$ 的系数。再把它与一个适当通用的特征标 $\mathbb{F}_q \rightarrow k^\times$ 复合，就得到 $N(\Fqt)$ 的一个特征标，仍记作 $\psi$。Whittaker 余不变量由
\[
\on{Whit}: G(\Fqt)\mod \rightarrow \Vect, \quad \pi \mapsto \pi_{N(\Fqt),
\psi}
\]
给出。它把单模送到一条直线或零，并具有一个关键性质：许多尖点表示，例如 $G = GL_n$ 时的全部尖点表示，都被送到直线。

这里的一条基本性质是：$\on{Whit}$ 有右伴随，由把 $k$ 送到 Whittaker 函数空间
\[
\Fun( G(\Fqt) / N(\Fqt), \psi),
\]
给出，也就是把 $\psi$ 上诱导到 $G(\Fqt)$。

\ssszh 上述构造有如下几何类比。首先，拓扑群的同态 \eqref{e:whitchar} 提升为群归纳概形的映射
\[
\psi: LN \rightarrow L(N/[N, N]) \simeq \underset i \Pi \hspace{.2mm}
L\mathbb{G}_a \rightarrow L\mathbb{G}_a \rightarrow \mathbb{G}_a.
\]
由此，我们想从 $\mathbb{G}_a$ 拉回「通用特征标」的类比物。回忆 $G(\Fq)$ 的平凡表示与 $G$ 在 $\Vect$ 上恒等范畴作用之间的类比，我们因此需要 $\mathbb{G}_a$ 在 $\on{Vect}$ 上的一个通用作用。代数群 $H$ 在 $\Vect$ 上的作用，等同于 $H$ 上的一个特征标层，即一个 D-模 $\chi$，配上适当结合的等同
\[
\mu^! \chi \simeq \chi \boxtimes \chi,
\]
其中 $\mu: H \times H \rightarrow H$ 是乘法映射。我们要强调：这里的用法与第 \ref{ss:charsheaves} 节中较宽泛的「特征标层」不同。那里的用法是 $G(\Fq)$ 的有限维表示的特征标的类比；这里的用法是一维表示即特征标的类比。

\begin{re} 稍不正式地说，若对 $H$ 的任意闭点 $h$ 记 $\chi_h$ 为 $\chi$ 相应的 $!$-茎，则上式给出茎之间的同构
\[
\chi_{h_1 h_2} \simeq \chi_{h_1} \otimes \chi_{h_2}, \quad \text{对 } h_1,
h_2 \in H.
\]
也就是说，特征标层给出一个映入直线的同态，即映入 $\Vect^\heartsuit$ 的可逆元的同态；正如特征标给出一个映入非零标量的同态，即映入 $k$ 的可逆元的同态。
\end{re}

有了这些，$\Fq$ 的通用特征标的类比物，就是 $\mathbb{G}_a$ 上的指数 D-模 $\on{exp}(z)$。确切地说，若以 $z$ 记 $\mathbb{G}_a \simeq \mathbb{A}^1$ 上的标准坐标，这个 D-模承载指数函数 $e^z$ 所满足的微分方程，即
\[
\on{exp}(z) \simeq D_{\mathbb{G}_a} / ( D_{\mathbb{G}_a} \cdot (\partial_z -
1)),
\]
而它的乘性本质上来自恒等式 $e^{z_1 + z_2} = e^{z_1} \cdot e^{z_2}$。据此，我们以 $\psi$ 记 $LN$ 上相应的特征标层，即 $\psi^! \on{exp}(z)$，并同样可构造 Whittaker 余不变量函子
\[
\on{Whit}: \Dmod(LG)\mod \rightarrow \Vect, \quad \sC \mapsto \sC_{LN,
\psi}
\]
以及 Whittaker D-模构成的范畴
\begin{equation} \label{e:catwhitsheaves}
\Dmod(LG / LN, \psi).
\end{equation}
在转向它的谱侧对应物之前，我们想作几条帮助定位的评注。

\begin{re} 对于 $\mathbb{G}_a(k)$ 上的指数 D-模与 $\mathbb{G}_q(\Fq)$\zy{[译者按（上下文推断，无外部依据）：原文此处作 $\mathbb{G}_q(\Fq)$，按上下文应为 $\Fq$；本译本照原文排印。]} 的通用特征标之间的类比，我们来给出三种理解方式。

第一种相当朴素。为简单起见，设 $\mathbb{F}_q$ 与其素子域 $\mathbb{F}_p$ 重合。此时，若以 $e^{2\pi i / p}$ 记 $k^\times$ 中一个非平凡的 $p$ 次单位根，则通用特征标的一个具体取法是
\[
\mathbb{F}_p \rightarrow k^\times, \quad a \mapsto e^{ a\cdot 2\pi i / p },
\]
其余取法与它相差 $\mathbb{F}_p$ 中一个非零元素的倍数。也就是说，通用特征标看起来就像指数函数。对非平凡的扩域 $\Fq$，同样的说法在作必要调整后也成立，只需先用迹映射 $\Fq \rightarrow \Fp$。

第二种也相当实际。给定 $k$-向量空间 $V$，其特征标层与 $V^*$ 的闭点典范地一一对应，其中余向量 $\xi$ 对应特征标层 $\xi^! \on{exp}(z)$。特别地，当 $V = \mathbb{G}_a$ 时，在差 $\mathbb{G}_m$ 的伸缩的意义下只有一种非平凡的特征标层；正如对 $\Fq$，在差 $\Fq^\times$ 的伸缩的意义下只有一个非平凡特征标。

第三种更一般，走的是函数-层对应。首先，正特征的一个美妙之处在于：$\Fq$ 上任意交换群概形 $H$ 都有一个有趣的自同源，即 Lang 同源，其形状为
\[
1 \rightarrow H(\Fq) \rightarrow H \rightarrow H \rightarrow 1,
\]
也就是说，它以有理点构成的有限群为核。对 $H = \mathbb{G}_a$，具体就是短正合列
\[
1 \rightarrow \mathbb{G}_a(\Fq) \rightarrow \mathbb{G}_a \xrightarrow{x
\mapsto x^q - x} \mathbb{G}_a \rightarrow 1.
\]
沿此同源前推常值层并取其直和项，得到 $\mathbb{G}_a$ 上一族特征标层，以 $\mathbb{G}_a(\Fq)$ 的不可约表示为指标。取 Frobenius 的迹之后，这些特征标层恰好给出 $\mathbb{G}_a(\Fq)$ 的不可约特征标。也就是说，$\Fq$ 的特征标可以典范地提升为 $\mathbb{G}_a$ 上的特征标层。

对这些特征标层，若只把它们看作 $\mathbb{A}^1$ 上的 étale 局部系统，则可以问它们在 $\infty$ 处的行为，即紧化到 $\mathbb{P}^1$ 之后的行为。结论是：它们在这个剩余点处野分歧。于是，它们与特征零情形中 $\on{exp}(z)$ 的对应，就纳入野分歧 Galois 表示与带非正则奇点的 de Rham 局部系统之间更广泛的类比之中。
\end{re}

\begin{re} 我们来评注 Whittaker D-模范畴 $\Dmod(LG / LN, \psi)$ 的定义。它的定义是取 $\Dmod(LG)$ 的 $(LN, \psi)$-余不变量，与半无限旗流形上 D-模的范畴类似。根据 Raskin 的一个重要定理 \cite{whit}，它同时也等同于 $(LN, \psi)$-不变量。在 $p$-adic 情形中，Whittaker 不变量与余不变量并不字面地等同；但上述等同仍可类比于 Rodier \cite{rodier} 对 Whittaker 模型的紧逼近。
\end{re}

\ssszh 现在来描述 Whittaker 模型预期的谱侧对应物。它就是 $$\QCoh(\LS),$$ 即局部系统模空间上的全部拟凝聚层。

这里的基本直觉同样可以取自 $p$-adic 情形。$G(\Fqt)$ 的每个不可约表示 $\pi$，其 Whittaker 余不变量至多构成一条直线。改用
\[
\Hom( \on{ind}_{N(\Fqt)}^{G(\Fqt)} k_{-\psi}, \pi) =: \ell,
\]
会更方便，其中 $\on{ind}$ 是紧支诱导函子。它与 $\pi$ 的对偶表示的 Whittaker 余不变量等同，特别地仍是一条直线或零。

另一方面，在局部朗兰兹对应之下，人们给 $\pi$ 配上一个 Galois 表示 $\sigma$，并把 $\pi$ 想象成在 Galois 表示的空间 $\on{Gal}_{\cG}$ 中「坐落于 $\sigma$ 之上」。

在算术局部朗兰兹对应的现代表述中（即作为范畴的等价，参见下文第 \ref{ss:variants} 节），人们期望：凡有非零 Whittaker 余不变量的 $\pi$，都与一个支撑在 $\sigma$ 上的一维摩天楼层 $\mathbb{L}(\pi)$ 互换。而且，这条典范地联系于 $\pi$ 的直线，预期就正是 $\ell$。于是，综合等同
\[
\Hom( \on{ind}_{N(\Fqt)}^{G(\Fqt)} k_\psi, \pi) \simeq \ell \simeq
\Hom( \OO_{\on{Gal}_{\cG}},  \mathbb{L} \pi ),
\]
可以看出，$\mathbb{L}$ 应当把（紧支）诱导的 Whittaker 模与结构层 $\OO_{\on{Gal}_{\cG}}$ 互换。

在几何情形中，Whittaker 模已换成相应的 D-模空间 $\Dmod(LG/LN,\psi)$。\footnote{注意：由于 Raskin 对 Whittaker 不变量与余不变量的等同 \cite{whit}，诱导与紧支诱导之间的差别消失了。} 类似地，我们把结构层 $\OO_{\on{Gal}_{\cG}}$ 换成 $\QCoh(\LS)$；后者同样余表示整体截面，只不过现在表示的是范畴层的截面。

正如我们在讲到缠绕算子时很快就会讨论的，这个对应会引出几何与表示论中许多引人注目的具体猜想。不过目前，我们只用两条注来结束对图 \ref{f:table} 中这一行的讨论。

\begin{re} 在 $p$-adic 情形中，并非每个缓增不可约表示都有 Whittaker 模型；局部几何朗兰兹的预期则不同：每个缓增的 $\Dmod(LG)$ 表示都有 Whittaker 模型。基本想法如下。设 $x$ 是叠 $X$ 的一个点，$I_x$ 是其迷向群。则在 $X$ 上以概形论方式支撑于 $x$ 的拟凝聚层构成的阿贝尔范畴，等同于 $\on{Rep}(I_x)$；特别地，除非 $I_x$ 幂幺，摩天楼层不会生成它。在算术局部朗兰兹对应之下，这与 Shahidi 的一般分歧包猜想相吻合：该猜想断言，在一个缓增 $L$-包内部，\footnote{超出 $GL_n$ 的情形，从 $G(\Fqt)$ 的不可约表示到 Galois 表示的映射不是双射，而是有限到一；其纤维称作 $L$-包。} 恰有一个不可约表示有 Whittaker 模型 \cite{shahidi90}。

与之相对，$\QCoh(X)$ 的摩天楼模 $\Vect$ 会生成 $\TwoCoh(X)$ 中支撑于 $x$ 的子范畴；在这个意义上，每个缓增不可约模都是通用的。尽管如此，由于前面许诺的对朴素表述 \eqref{e:loclangnaive} 的修正所牵涉的问题，$\Dmod(LG/LN, \psi)$ 并不会成为 $\Dmod(LG)$ 的生成元。
\end{re}

\subsection{对应的性质 V：Kac--Moody 表示}

\ssszh 我们已来到图 \ref{f:table} 的最后一行，这里会发生某种相当引人注目的事情。确切地说，我们将遇到对应的一条猜想性质，由 Frenkel--Gaitsgory \cite{frenkgaitslocallangds} 提出，它涉及 Kac--Moody 表示。这样的性质在 $p$-adic 情形中没有对应物，是局部朗兰兹中几何情形特有的一个侧面。

\begin{re} 我们应当提到，Kac--Moody 表示在几何朗兰兹中的重要性，最早出现在整体的（global）故事中。如我们在第 \ref{ss:hitchinhecke} 节所暗示的，Beilinson--Drinfeld 给出的一种构造自守 D-模的方法，就是通过 Kac--Moody 表示的局部化。这可以看作 Frenkel--Gaitsgory 工作的整体类比，也是「与共形场论的联系提供了算术情形所不具备的工具」这回事第一次出现。\end{re}

\ssszh 我们想先让读者对 Kac--Moody 代数有一个基本的感觉。对一般的连通约化群，它们的讨论本质上化归到 $\fg$ 单的情形。因此，为避免多半只是记号上的困扰，我们假定 $G$ 是单连通、几乎单的群。

首先从环路代数开始
\[
\fg(\!(t)\!) \simeq \fg \underset  k \otimes k(\!(t)\!).
\]
这不过是把 $\fg$ 作标量扩张到 Laurent 级数域。确切地说，括号的形式为
\[
[ X \otimes f, Y \otimes g] = [X,Y] \otimes fg, \quad X,Y \in \fg, \quad
f, g \in k(\!(t)\!).
\]
把它想成从穿孔圆盘 $\sD^\times$ 到 $\fg$ 的映射空间、而李括号逐点进行，会很有帮助。也就是说，它是穿孔圆盘上平凡 $G$-丛的无穷小对称。

我们把环路代数看作 $k$ 上的无穷维李代数。一个基本的事实是：许多乍看像是 $\fg(\!(t)\!)$-模的 $k$-向量空间，实际上带有中心扩张的作用，即仿射李代数的作用
\[
0 \rightarrow k \cdot \mathbf{1} \rightarrow \gk \rightarrow
\fg(\!(t)\!) \rightarrow 0.
\]

\begin{re} 这可以看作一般现象的一例：场论的量子化会给其局部对称性带来反常。\end{re}

幸运的是，环路代数的一切可能中心扩张都有简单的分类——它们与有限维代数 $\fg$ 上不变双线性形式所在的直线典范地一一对应。确切地说，对这样的形式 $\kappa \in (\fg^* \otimes \fg^*)^G$，相应地有中心扩张
\[
0 \rightarrow k \cdot \mathbf{1} \rightarrow \fgk \rightarrow \fg(\!(t)\!)
\rightarrow 0,
\]
其中由中心性，$\mathbf{1}$ 与任何元素的括号为零，其余括号则修正为
\[
[X \otimes f, Y \otimes g] = [X,Y] \otimes fg - \kappa(X,Y) \on{Res} (f
\cdot dg).
\]

\ssszh 对水平 $\kappa$，以 $\fgk\mod$ 记光滑 $\fgk$-模构成的范畴，其中中心元 $\mathbf{1}$ 通过恒等作用。参见例 \ref{ex:kmmods}。若 $\kappa$ 是整的，即环路代数的扩张来自 $LG$ 的中心扩张，那么如例 \ref{ex:kmmods} 所述，这个范畴带有 $\Dmod(LG)$ 的作用；因此下面只关注这一情形。

回忆 $G$ 被假定为单且单连通，整水平构成一个格
\[ \mathbb{Z} \hookrightarrow (\fg^*
\otimes \fg^*)^G,\]
由所谓的基本形式生成；具体说来，它使短余根的长度的平方为二。

\ssszh 仿射李代数表示论中的一个基本观察是：随着水平 $\kappa$ 变化，相应的表示论有两个「相」。

确切地说，对一切充分负定的 $\kappa$，表示论的表现类似。例如，所有 Verma 模都有有限长度，而这些水平上其单商的特征标公式本质上具有相同的形式 \cite{kashiwaratanisakinegativeintegrallevel}。

同样，对一切正定的 $\kappa$，表示论随水平变化的表现也类似。此时 Verma 模的长度无限，其单商的特征公式不随水平发生实质变化 \cite{kashiwara90}。

因此可以预期存在一个特殊水平 $\kappa_c$，即临界水平，表示论在此发生相变。朴素地看，人们或许以为这个中心点就是零形式。其实，由于 $\fg$ 表示论中 $\rho$-移位的仿射类比，它是 Killing 形式的一半再取负，即
\[
\kappa_c = - \frac{1}{2} \cdot \kappa_{\on{Killing}}.
\]

\ssszh 正如所料，相变处即临界水平处，表示论的行为显示出其他水平所没有的许多特征。

Feigin 与 Frenkel \cite{feiginfrenkel91} 发现了这样一条基本性质。非临界水平处，包络代数的中心\footnote{技术上说，我们指适当的完备包络代数的中心，或等价地，光滑 $\fgk$-模阿贝尔范畴的 Bernstein 中心。}是平凡的，即只有标量；而他们证明临界水平处的中心非平凡。关键在于，他们进一步把它等同于穿孔圆盘上 $\cG$-oper 的模空间 $\Op$ 上的函数代数。虽然我们马上会回忆它的精确定义，但首先的要点是：临界水平处的表示论把局部朗兰兹对偶、尤其是 $\cG$-联络的几何，硬性地编码于自身之中。

\ssszh 粗略地说，一个 $\cG$-oper 由一个带联络的 $\cG$-丛和一个 $\cB$-约化构成，二者适当地横截。

我们把它写清楚。回忆一下，为了谈论朗兰兹对偶群，需要固定 $G$ 与 $\cG$ 的带钉结构，特别地固定 $\check{\fg}$ 的 Chevalley 生成元 $f_i, \check{\alpha}_i, e_i, i \in I$。若以 $f$ 记 $\cfg$ 中使
\[
(f, 2 \check{\rho}, \underset i \Sigma \hspace{.2mm} e_i )
\]
构成 $\fsl_2$ 三元组的元素，则 $\cG$-oper 的空间是联络的模空间
\begin{equation} \label{e:defopers}
\Op \simeq \{ d  + f dt + \check{\fb}(\!(t)\!)dt \} / L\check{N}.
\end{equation}
把平凡丛看作带有其恒等的 $\cB$-约化，则上式中的 $f dt$ 一项就实现了上述联络与约化的横截性。关于 oper 定义的其他表述（例如具有明显不依赖坐标 $t$ 的优点），可参见 \cite{BD}。

公式 \eqref{e:defopers} 可以读作：取平凡 $\cG$-丛上全体联络 $d + \cfg(\!(t)\!)dt$ 的某个仿射子空间，再按 $L\check{N}$ 的伴随作用取商。事实上，\eqref{e:defopers} 是仿射空间 $d + \cfg(\!(t)\!)dt$ 关于 $L\check{N}$ 的一个通用特征标的 Hamilton 约化，特别地带有典范的 Poisson 结构。Gelfand--Dikii 与 Drinfeld--Sokolov \cite{drinfeldsokolov} 在几何朗兰兹出现之前，就在可积层级的背景下研究了它的 Poisson 几何；对 $\fg = \fsl_2$，这恢复了著名的 KdV 层级。

对尚未遇到过这个模空间的读者，这里有两条帮助定位的注。

\begin{re} 虽然 $\Op$ 与 $\LS$ 类似，被表示成联络的空间模去环路群的作用，但它的几何要简单得多，两者不可同日而语。事实上，它本身就是个无穷维仿射空间；特别地，它的点没有自同构。
\end{re}

\begin{re} 上一条注通常这样证明：说明每个 $L\check{N}$-轨道都是单传递的，即没有非平凡稳定化子，并给出与每个轨道恰好相交一次的显式截面。

指出整个这套说法（包括上述论证）是 Kostant \cite{kostantslice} 一项工作的仿射类比，会很有帮助。他证明了有典范同构
\[
\cfg /\!/ \cG \simeq (f + \check{\fb}) / \check{N},
\]
也就是说，对 $\cfg$ 作类似的有限维 Hamilton 约化，会切出 Poisson 中心的一个拷贝。此外，他证明两边的自然量子化是匹配的，即中心 $Z(\cfg)$ 可以由 $U(\cfg)$ 关于 $\cN$ 的量子 Hamilton 约化切出。\end{re}

\ssszh 通过投影映射 $\Op \rightarrow \LS$（从 oper 过渡到其底层的联络），可以把 $\QCoh(\Op)$ 看作 $\LS$ 上的范畴层。

于是人们预期，在局部几何朗兰兹下 $\gkc\mod$ 与 $\QCoh(\Op)$ 互换；从而恒等的幺半函子
\[
\QCoh(\Op) \rightarrow \Hom_{\TwoCoh(\LS)}(\QCoh(\Op), \QCoh(\Op))
\]
经由 Feigin--Frenkel 同构 $\OO_{\Op} \simeq Z(\gkc)$，与中心作用
\[
Z(\fgk)\mod \rightarrow \Hom_{\Dmod(LG)\mod}(\gkc\mod, \gkc\mod)
\]
互换。

在下面几个小节中，我们会讨论这个对应给出的、关于临界水平表示论的若干有趣猜想。

不过我们想在此先提一个特别引人注目的推论。已知任何 $\cG$-联络 $\sigma$ 都带有 oper 结构 $\chi$，即投影 $\Op \rightarrow \LS$ 在 $k$-点上满射 \cite{frenkelzhuoper}。由上述，形式上可推出：环群对应于 $\sigma$ 上摩天楼 $\Vect$ 的「不可约」范畴表示 $\sC_{\sigma}$，由
\begin{equation} \label{e:kmcenc}
\gk\mod \underset{\QCoh(\Op)} \otimes \Vect,
\end{equation}
给出，即中心特征标为 $\chi$ 的 Kac--Moody 表示范畴。

一方面，这通过 Kac--Moody 表示给出了全部 $\sC_{\sigma}$ 的统一猜想构造，即局部几何朗兰兹对应的谱到自守方向。另一方面，这预言：对任何两个位于 $\sigma$ 之上的 oper $\chi, \chi'$，相应的 Kac--Moody 模范畴等价——这本身就是表示论中一个深刻的猜想。

\begin{re} 事实上，由于尚未解释的、对朴素表述 \eqref{e:loclangnaive} 的修正，范畴 \eqref{e:kmcenc} 其实只是 $\sC_{\sigma}$ 的全子范畴。实践上，这意味着通过对 \eqref{e:kmcenc} 作重正规化来得到 $\sC_{\sigma}$，进一步的讨论参见第 \ref{ss:locaff} 节。\end{re}

\ssszh 这就结束了我们对图 \ref{f:table} 中出现的对象的引导性讨论。不过，在继续之前，我们应当强调：来自算术故事以及与物理的联系，还有更多预期的相容性，参见下文第 \ref{ss:variants} 节。

我们忍不住要提一个特别壮观的例子。Braverman--Finkelberg 对 $G = GL_n$ 提出了卷积代数 $\Dmod(LGL_n)$ 本身的朗兰兹对偶的一个候选 \cite{bravermanfinkelbergconjecturecoloumb}。这基于物理上的考虑，最终可追溯到 Gaiotto--Witten \cite{gaiottowitten} 发现的一个引人注目的公式，用于计算四维 $\mathscr{N}=4$ Yang--Mills 理论的 S-对偶边界条件。确切地说，他们猜想：在局部朗兰兹对偶下，$\Dmod(LGL_n)$ 与参数化局部系统退化旗的模空间上的凝聚层互换，即秩为 $i$ 的局部系统 $\sL_i$（$1 \leqslant i \leqslant n$）连同未必单射的联络映射
\[
\sL_i \rightarrow \sL_{i+1}, \quad 1 \leqslant i < n.
\]
此后 Zhu \cite{zhufamilies} 提出了这一猜想的算术类比。

\subsection{缠绕算子}

接下来，我们想解释：考察图 \ref{f:table} 中出现的范畴之间的缠绕算子，如何重新得到若干基本定理，并产生许多引人注目的猜想。

在这一点上，前面许诺的对局部几何朗兰兹对应朴素表述 \eqref{e:loclangnaive} 的技术性修正，自然就以如下方式显现出来。

\ssszh 回忆 $p$-adic 情形中最简单的表示是非分歧表示。令人高兴的是，在非分歧情形中，我们已能遇到几何猜想朴素表述 \eqref{e:loclangnaive} 所需的修正。

为看清这一点，回忆局部朗兰兹应当把仿射 Grassmann 簇上的 D-模 $\Dmod(\Gr_G)$，与平凡丛上的拟凝聚层 $\QCoh(\on{pt}/\cG)$ 互换。特别地，二者的自同态应当典范地等同；我们来看看它们是什么。

在自守侧，回忆 $\Gr_G$ 是陪集空间 $LG/L^+G$，如第 \ref{s:catrepgrp} 节所述，有
\begin{equation} \label{e:dersat1}
 \Hom_{\Dmod(LG)\mod}( \Dmod(\Gr_G), \Dmod(\Gr_G)) \simeq \Dmod(L^+G \backslash
 LG / L^+G),
\end{equation}
即双球层的导出 Satake 范畴。作为对初次接触此事的读者的可能提醒：这个范畴不是相应阿贝尔范畴的导出范畴，而更为丰富。

在谱侧，利用映射 $\cfg / \cG \rightarrow \LS$ 在平凡丛处形式 étale 这一事实，可以如例 \ref{e:fibprodvar} 那样算出
\begin{equation} \label{e:dersat2}
\begin{split}
&\Hom_{\TwoCoh(\LS)}(\QCoh(\on{pt}/\cG), \QCoh(\on{pt}/\cG)) \\
&\qquad \simeq \QCoh( \on{pt}/\cG \underset{\LS} \times \on{pt}/\cG)
\simeq \QCoh( \on{pt} \underset{\cfg} \times \on{pt} / \cG).
\end{split}
\end{equation}
这里需要两点解释。第一，对 \eqref{e:dersat2} 中最后的等同，可以用下面这条一般观察。设有代数群 $K$ 作用在簇 $X, Y, Z$ 上。给定 $K$-等变的对应
$$X \rightarrow Z \leftarrow Y,$$
则有商（导出）叠的典范同构
\[
      X/H \underset{Z/H} \times Y/H \simeq (X \underset Z \times Y)/H.
\]
今后我们会反复使用这一事实，不再另行说明。

第二，提醒读者：\eqref{e:dersat2} 中出现的纤维积是导出概形。具体地说，这只是意味着我们同时记住 Tor；于是自交
$$\on{pt} \underset{ \fg} \times
\on{pt}$$
是用 Koszul 消解作标准计算得到的、外代数 $\Sym \cfg^*[1]$ 的谱。

因此，我们可以朴素地期望 \eqref{e:dersat1} 与 \eqref{e:dersat2} 之间有幺半范畴的等价。其实，作为 Bezrukavnikov--Finkelberg \cite{bezfinksat} 的一个重要定理，导出 Satake 等价把 \eqref{e:dersat1} 与 \eqref{e:dersat2} 的一个放大等同起来。更确切地说，他们证明了等价
\begin{equation} \label{e:dersat} \Dmod(L^+G \backslash LG / L^+G) \simeq
\ICn( \on{pt} / \cG  \underset{
\cfg / \cG} \times \on{pt} / \cG), \end{equation}
其中右端是带{\em 幂幺奇异支撑}的归纳凝聚层构成的幺半范畴。

\ssszh 我们简要说明什么是归纳凝聚层，以及施加奇异支撑条件是什么意思。它们算是对拟凝聚层这一想法的一些技术性修正，初次接触这一圈想法的读者可以放心跳过。我们先讲归纳凝聚层，再转向奇异支撑条件。

用最朴素的话说，归纳凝聚层的理论源自有限生成模与有限生成投射模（即向量丛）之间的差别。在阿贝尔范畴中，比如对仿射概形 $X$，向量丛当然比有限生成模少，除非 $X$ 是若干约化点的无交并。

然而在导出范畴中，这一区别更为微妙。事实上，若 $X$ 光滑，则任何有限生成模都容许由向量丛给出的有限消解。因此，人们无法区分有限生成模的有限复形与向量丛的有限复形。若 $X$ 不光滑，它们在导出范畴中又有差别：一般的有限生成模此时需要无限投射消解。具体地说，可以取 $X$ 的任一奇点处的摩天层。特别地，这两类复形之间的差别给出了 $X$ 奇点程度的一个度量。

\ssszh 我们按照 Krause \cite{krause} 与 Gaitsgory \cite{gaitsgoryindcoh} 的做法，把前面几段讲得更仔细些。首先回忆：对代数簇 $X$，其拟凝聚层范畴 $\QCoh(X)$ 由完美复形紧生成，即由局部等价于向量丛的有限复形的对象紧生成。若 $X$ 仿射，这就化归为同义反复：$O_X\mod$ 由正则函数环 $O_X$ 生成。对非仿射的 $X$，所断言的在拟射影簇情形容易，但一般而言是 Thomason \cite{thomasontrobaugh} 一个颇为精细的定理。

作为有限生成模的有限复形的整体类比，考虑凝聚层的有界导出范畴
$$\on{Coh}(X)^b \subset \QCoh(X).$$ 按定义，$\on{IndCoh}(X)$ 是它的归纳完备化，即 $\on{IndCoh}(X)$ 由 $\on{Coh}(X)^b$ 紧生成。显然，$\on{IndCoh}(X)$ 的对象是上同调凝聚的有界复形的形式余极限，而同态复形按复形的同伦极限计算
\[
\Hom_{\on{IndCoh}(X)}( \varinjlim_i \sE_i, \varinjlim_j \sF_j) \simeq
\varprojlim_i
\varinjlim_j \Hom_{\on{QCoh}(X)}( \sE_i, \sF_j).
\]
对任何凝聚层 $\sE$，此式迫使 $\Hom_{\on{IndCoh}(X)}(\sE, -)$ 与余极限交换；而在 $\QCoh(X)$ 中，除非 $\sE$ 完美，这一点不成立。

由于完美复形有界凝聚，把它们的包含函子作归纳完备化，就得到伴随对
\[
\on{QCoh}(X) \rightleftarrows \on{IndCoh}(X),
\]
其中左伴随全忠实，且当 $X$ 光滑时是等价。特别地，人们把 $\on{IndCoh}(X)$ 看作 $\on{QCoh}(X)$ 的放大，并且任一归纳凝聚层都有到拟凝聚层的截断。此外，$\on{IndCoh}(X)$ 带有自然的 $t$-结构，并在有下界的对象上与 $\on{QCoh}(X)$ 一致。不正式地说，$\on{IndCoh}(X)$ 是通过添加上同调次数为 $-\infty$ 的对象得到的。

\begin{re} 回忆一下，我们在讨论 Kac--Moody 表示与环路群上的 D-模时，已经遇到过这种对无限上同调维数的导出范畴的重正规化，参见例 \ref{ex:kmmods}；这或许有助于定位。\end{re}

\begin{re} 在奇点簇上系统地区分完美复形与有界凝聚复形这一想法，此前也出现在代数几何中。例如，商 $\on{IndCoh}(X)/\QCoh(X)$ 就是 Orlov 奇点范畴 \cite{orlovsingularity} 的归纳完备化。\end{re}

\ssszh 接下来解释奇异支撑的想法。它由 Gulliksen \cite{gulliksen}、Eisenbud \cite{eisenbud}、Avramov 与 Buchweitz \cite{avramovbuchweitz} 针对完全交发展起来，又由 Arinkin--Gaitsgory 推广到导出完全交 \cite{arinkingaitssingsupp}；另可参见重要的 \cite{bensoniyengarkrause}。按定义，归纳凝聚层与拟凝聚层之间的差别源自 $X$ 的奇点；特别地，商 $$\IC(X)/\QCoh(X)$$ 将沿 $X$ 的奇异轨迹支撑。给定凝聚层 $\sF$，人们可能想要一个不变量，粗略地记录它在奇点处的何处不再是完美的。奇异支撑对这一问题的回答更为精细。

下面是稍微精确一点的表述。设 $X$ 是拟光滑导出概形，即其余切复形 $T^*(X)$ 支撑在上同调次数 $-1$ 与 $0$ 处。则对上述对象 $\sF$，其奇异支撑是 $H^{-1} T^*(X)$ 的一个闭锥形子集 $\Lambda$。奇异支撑逐仿射地度量 $\OO_X$ 的元素的局部幂零性，从而恢复通常的支撑，以及 $\sF$ 的某些自然的二次自映射。此外，对任何这样的子集 $\Lambda$，都得到全子范畴
\[
   \IC_\Lambda(X) \subset \IC(X)
\]
由奇异支撑含于 $\Lambda$ 的归纳凝聚层构成。例如，若 $\Lambda$ 是零截面，这恢复 $\QCoh(X)$。

虽然我们不打算给出计算奇异支撑的一般定义，但马上会讨论几个特例。不过先给出几条帮助定位的注。

\begin{re}\label{r:singsupp}导出概形拟光滑，当且仅当它局部可写成纤维积 $X \times_Z Y$，其中 $X,Y,Z$ 都是光滑簇。在这个局部图景中，余切复形有如下显式形式。确切地说，对纤维积的一点 $(x,y)$，设其公共像为 $z$，则有切复形的 distinguished 三角形
\begin{equation} \label{e:dts}
    T_{(x,y)} (X \times_Z Y) \rightarrow T_x(X) \oplus T_y(Y) \rightarrow
    T_z(Z) \xrightarrow{+1}.
\end{equation}
对不熟悉这个 distinguished 三角形的读者，先理解 $X$ 与 $Y$ 横截地映入 $Z$ 的情形会有帮助。此时纤维积仍然光滑，上述 distinguished 三角形变成一个（希望是）直观的短正合列。在一般情形，要点在于：纤维积上的导出结构修正了它的切复形，使得横截情形下最朴素的修正依然成立。

把 \eqref{e:dts} 对偶化，可以看出余切丛的 $H^{-1}$ 具体就是那些与 $T_x(X)$、$T_y(Y)$ 的像正交的余方向 $\xi \in T^*_z(Z)$；更不正式地说，就是在 $Z$ 中还剩下可以走的方向。

上述情形的一个重要特例是形如
\begin{equation}
\label{e:vanlocus} \{0\}
\times_{\mathbb{A}^n} Y,\end{equation}
的纤维积，即 $Y$ 上 $n$ 个函数 $f_1, \ldots, f_n$ 的导出零点集。由此可见，拟光滑性是导出代数几何中「隐藏光滑性」哲学的一例。也就是说，纤维积底层的经典概形可以任意奇异，但它所承载的导出概形表现得像完全交，即只是轻度奇异。

\end{re}

下面用几个例子说明奇异支撑究竟是什么。

\begin{ex}\label{ex:singsupp1}对凝聚层 $\sF$，其奇异支撑与零截面的交就是通常的支撑。\end{ex}

\begin{ex}\label{ex:singsupp2}设 $X$ 可实现为光滑簇 $U$ 中的超曲面。此时 $$H^{-1} T^*(X)$$ 由 $X$ 本身（看作零截面）连同放在 $X$ 每个奇点 $x$ 处的一条直线 $\ell_x$ 构成。设凝聚层 $\sF$ 的通常支撑含有这样的 $x$。则其奇异支撑含有 $\ell_x$，当且仅当 $\sF$ 在 $x$ 的任一 Zariski 邻域内都不与完美复形同构。\end{ex}

\begin{ex}设 $X$ 如导出 Satake 的谱侧那样，由向量空间中原点的自交给出：$$X = 0 \underset V \times 0.$$ 此时有 $T^*_{(0,0)} X \simeq V^*[1]$。这里，奇异支撑可以理解为 Koszul 对偶另一侧的通常支撑。确切地说，摩天层 $k$ 生成 $\IC(X)$，取它的自同态就得到等价
\[
    \IC(X) \simeq \Sym V[-2] \mod.
\]
经由这个等同，对任何凝聚层 $\sF$，其奇异支撑与相应 $\Sym V[-2]$-模的通常支撑一致（忽略分次的平移）。这里可以把 $0$ 与 $V$ 换成任意簇的任一光滑点。\end{ex}

\ssszh 讨论了归纳凝聚层与奇异支撑的基础之后，我们回到导出 Satake 等价，即
\begin{equation} \label{e:der} \Dmod(L^+G \backslash LG / L^+G) \simeq
\ICn( \on{pt} / \cG  \underset{
	\cfg / \cG} \times \on{pt} / \cG). \end{equation}
在这个等价中，下标「nilp」指的是归纳凝聚层中奇异支撑落在幂零锥
$$\sN \subset \cfg^* \simeq
H^{-1}T^*( \on{pt}/\cG \underset {\cfg / \cG}   \times \on{pt}/\cG)$$
之中的子范畴。

\ssszh 我们马上会看到，非分歧情形中遇到的模式——必须放大谱侧，允许带幂幺奇异支撑的归纳凝聚层——在另一些匹配缠绕算子的情形中也会出现。这就迫使我们对局部几何朗兰兹的谱侧作调整，办法是从范畴的拟凝聚层过渡到带幂幺奇异支撑的范畴归纳凝聚层
\[
\TwoCoh(\LS) \rightsquigarrow \TwoIC(\LS).
\]
我们马上会谈这指的是什么；目前只需说，有了这一修正，就得到局部几何朗兰兹猜想的现有表述，即等价
\begin{equation}
\label{e:loclang}\Dmod(LG)\mod \simeq \TwoIC(\LS).
\end{equation}

\ssszh 我们简要勾勒带奇异支撑的范畴归纳凝聚层理论，它属于 Arinkin \cite{arinkinindcohcats}。

当范畴层级从归纳凝聚层升到范畴归纳凝聚层时，要建立奇异支撑理论，人们不再要求底簇 $X$ 拟光滑（即余切复形在次数 $-1$ 与 $0$），而要求它光滑（即余切复形在次数 $0$）。原因马上可见。

于是设 $X$ 是光滑代数簇。利用 1-仿射性可以证明：对任何从光滑簇出发的紧合映射 $f: Z \rightarrow X$，$\QCoh(Z)$ 紧生成 $\TwoCoh(X)$ 中支撑于 $\on{im}(f)$ 之上的对象构成的全子范畴。特别地，这个子范畴等同于卷积代数
\[
  \QCoh(Z \underset X \times Z)
\]
上的模。

$X$ 上范畴归纳凝聚层的 $(\infty,2)$-范畴有如下基本性质。第一，它带有 $\TwoCoh(X)$ 的幺半作用，因而可以谈论对象的支撑。此外，它由形如 $\QCoh(Z)$ 的对象（$Z$ 如上）生成。然而，此类对象的自同态现在提升为
\[
   \IC( Z \underset{X}{\times} Z),
\]
粗略地说，随着 $Z \rightarrow X$ 越来越远离光滑映射，它相对于拟凝聚的类似物会变得更大。更确切地说，如注 \ref{r:singsupp} 所述，这个卷积范畴中对象的奇异支撑，沿 $X$ 中与 $Z \times Z$ 的切丛正交的余方向分布。因此，$f$ 在给定点 $x$ 上越不光滑，$\on{2-IndCoh}(X)$ 中由 $\QCoh(Z)$ 生成、且支撑于 $x$ 之上的范畴归纳凝聚层就越多。

特别地，$\QCoh(Z)$ 不再生成 $\on{2-IndCoh}(X)$ 中对应于支撑于 $\on{im}(f)$ 之上的对象的全子范畴，而是生成如下微局部细化后的子范畴。\footnote{确切地说，要看出细化是必要的，注意：若没有细化，取 $Z = X$ 就会把 $\on{2-IndCoh}(X)$ 与 $\TwoCoh(X)$ 等同起来。} 由于假定 $f$ 紧合，可以构造余法簇 $N^\vee_{Z}$，即由沿 $$Z \underset X \times T^*X \rightarrow T^*X$$ 投影 $f^* T^*X \rightarrow T^*Z$ 的核所得到的、$T^*X$ 的闭各向同性子簇。可以给任何范畴归纳凝聚层配上它在 $T^*X$ 中的支撑，而 $\QCoh(Z)$ 生成微支撑落在 $N^\vee_Z$ 之中的对象构成的全子范畴。

更一般地，给定 $T^*X$ 的任一闭锥形各向同性子簇 $\Lambda$，可配上全子范畴
$$
\on{2-IndCoh}_\Lambda(X) \subset \on{2-IndCoh}(X).
$$
它同样由 $\QCoh(Z)$ 的（截断）生成，自同态为 $$\on{IndCoh}_\Lambda(Z \underset X \times Z),$$ 其中 $\Lambda$ 现在用来对通常的归纳凝聚层施加通常的奇异支撑限制，如注 \ref{r:singsupp}。特别地，取 $\Lambda$ 为零截面，就恢复 $\TwoCoh(X)$。

\begin{re} 归纳凝聚层与范畴归纳凝聚层的形式体系还有另一种进路，其中包含分别超出拟光滑与光滑情形的奇异支撑理论；该进路正由 di Fiore--Stefanich 发展中。\end{re}

\ssszh 前面的讨论并不字面地适用于 $\LS$。粗略地说，原因在于它无穷维，不是光滑 Artin 叠。不过人们预期，它带有类似的范畴归纳凝聚层理论。

承认这种形式体系存在，则由 \eqref{e:loclang}，只剩一件事要解释：奇异支撑条件「nilp」指的是什么。不过，鉴于我们在非分歧情形已经遇到的东西，这应当不足为奇。确切地说，$\LS$ 在联络 $(\check{\sP}, \nabla)$ 处的余切复形，由计算余伴随丛
\begin{equation} \label{e:datboi}
   \check{\sP} \overset{\cG} \times \cfg^*
\end{equation}
在穿孔圆盘上平坦截面的 de Rham 复形给出。特别地，余切空间即其零次上同调，典范地等同于 \eqref{e:datboi} 的平坦截面。于是，「nilp」在每个余切空间内切出那些幂幺的平坦截面。

\ssszh 处理完必要的修正，现在我们终于兑现承诺，来实际看看图 \ref{f:table} 中各对象之间缠绕算子构成的范畴。这样做，我们会遇到许多漂亮的定理与猜想，希望这能让读者对本学科有更多亲自动手的感觉。

两侧出现的范畴都是自对偶的。对 D-模与拟凝聚层的范畴，除开一些可驾驭的无穷维问题，这由例 \ref{ex:dmodselfdual} 得出。Kac--Moody 表示的情形我们会在下文第 \ref{ss:affhc} 节更详细地讨论。承认这种自对偶，经由典范等价
\[
\Hom_{\Dmod(LG)\mod}(\sC, \sD) \simeq \Hom_{\Dmod(LG)\mod}(\sD^\vee, \sC^\vee),
\]
以及对 $\LS$ 上的范畴类似的关系，我们只需写下「对角线上方」的缠绕算子范畴。

这些缠绕算子记录在图 \ref{f:tableaut} 与图 \ref{f:tablespec} 中。自守侧与谱侧这些范畴的计算，基本上遵循第 \ref{s:catrepgrp} 节、尤其是例 \ref{e:fibprodvar} 与第 \ref{ss:invariantscosets} 节讨论的模式。下面我们来比较两张表中的条目，并对出现的等价（无论已知还是未知）加以评注。

\begin{table}[H]
\centering
\scriptsize
\setlength{\tabcolsep}{1.6pt}
\begin{tabular}{|c|c|c|c|c|}
\hline
 & $\Dmod(\Gr_G)$ & $\Dmod(\Fl_G)$ & $\Dmod(LG/LN, \psi)$ & $\fgc\mod$ \\
\hline
$\Dmod(\Gr_G)$ & $\Dmod(L^+G \backslash LG / L^+G)$ & $\Dmod(L^+G
\backslash LG / I)$ & $\Dmod(L^+G \backslash LG / LN, \psi)$ &
$\fgc\mod^{L^+G}$ \\
\hline
$\Dmod(\Fl_G)$ & & $\Dmod(I \backslash LG / I)$ &
$\Dmod(I \backslash LG / LN, \psi)$ & $\fgc\mod^I$ \\
\hline
$\Dmod(LG/LN, \psi)$ & & & $\Dmod(LN, \psi \backslash
LG / LN, \psi)$ & $\fgc\mod^{LN, \psi}$ \\
\hline
$\fgc\mod$ & & & & $\fgc\oplus
\fgc\mod^{LG}$ \\
\hline
\end{tabular}
\caption{局部几何朗兰兹自守侧的一些缠绕算子。}
\label{f:tableaut}
\end{table}

\begin{table}[H]
\centering
\scriptsize
\setlength{\tabcolsep}{1.6pt}
\begin{tabular}{|c|c|c|c|c|}
\hline
 & $\QCoh(\on{pt}/\cG)$ & $\QCoh(\widetilde{\sN} / \cG)$
& $\QCoh(\LS)$ & $\QCoh(\Op)$ \\
\hline
$\QCoh(\on{pt} / \cG)$ & $\ICn( \on{pt} / \cG \underset {\cfg/\cG} \times
\on{pt} / \cG)$ & $\ICn(\on{pt} / \cG \underset { \cfg / \cG} \times
\widetilde{\sN} / \cG)$ & $\QCoh( \on{pt}/ \cG)$ &
$\QCoh( \Op^{\on{unramified}})$ \\
\hline
$\QCoh(\widetilde{\sN} / \cG)$ & & $\ICn( \widetilde{\sN}/ \cG \underset{
\cfg / \cG } \times \widetilde{\sN} / \cG)$ & $\QCoh( \widetilde{\sN}/
\cG)$ & $\QCoh(\widetilde{\sN}/\cG \underset \LS \times \Op)$ \\
\hline
$\QCoh(\LS)$ & & & $\QCoh(\LS)$ & $\QCoh(\Op)$ \\
\hline
$\QCoh(\Op)$ & & & & $\QCoh( \Op \underset \LS \times \Op)$ \\
\hline
\end{tabular}
\caption{局部几何朗兰兹谱侧的一些缠绕算子。}
\label{f:tablespec}
\end{table}

\ssszh 从非分歧情形开始，即表中的第一行。第一个等价
\begin{equation}
\Dmod(L^+G \backslash LG / L^+G) \simeq \ICn( \on{pt} / \cG \underset {\cfg/\cG}
\times \on{pt} / \cG)
\end{equation}
我们已经遇到过，即 Bezrukavnikov--Finkelberg \cite{bezfinksat} 的导出 Satake 等价。

导出 Satake 等价是 $t$-正合的，在阿贝尔范畴上恢复几何 Satake 等价
\[
 \Dmod(L^+G \backslash LG / L^+G)^\heartsuit \simeq \QCoh(\on{pt} / \cG) =
 \on{Rep}(\cG),
\]
后者是更早的一个基本定理，属于 Lusztig \cite{lusztigqchar}、Drinfeld、Ginzburg \cite{ginzburgsatake}、Mirković--Vilonen \cite{mirkovic-vilonen}。这两个等价（其中阿贝尔的那个稍更透明）都是 $p$-adic 情形中类似断言的范畴化：
\[
   \on{Dist}( G(\Fq[[t]]) \backslash G(\Fq(\!(t)\!)) / G(\Fq[[t]])) \simeq K_0(
   \on{Rep}(\cG)) \underset{\mathbb{Z}}{\otimes} k,
\]
该同构连同其混合特征的情形，属于 Satake \cite{satake63}。

对此作两点评注。

\begin{re} \label{r:paramcount} 第一，对初次接触的读者，为了对这个等价有所把握，配比参数（即（几乎）紧对象的 Grothendieck 群的基）是很有启发的。

基本要点是：$LG$ 上只有可数多个 $L^+G \times L^+G$-轨道，它们自然地以 $\cG$ 的不可约表示为指标。更仔细地说，任何余特征标
\[
   \check{\mu}: \mathbb{G}_m \rightarrow T
\]
限制到原点处的形式圆盘，就给出 $LT$ 的一点，记作 $t^{\check{\mu}}$。若以 $\check{\Lambda}^+$ 记支配余特征标，则 $LG$ 被轨道
\[
   L^+G \cdot t^{\check{\mu}} \cdot L^+G, \quad \text{对 } \check{\mu} \in
   \check{\Lambda}^+,
\]
分层。而这些正是 $\cG$-模的最高权；该等价把 $\Gr_G$ 上相应的相交上同调复形与单 $\cG$-模互换。\end{re}

\begin{re} 第二，这是一条帮助定位的注。导出 Satake 自守侧底层的空间，参数化三元组 $(\sP_1, \sP_2, \tau)$，其中 $\sP_1, \sP_2$ 是圆盘上的 $G$-丛，$\tau$ 是它们在穿孔圆盘上的等同。

类似地，谱侧底层的空间参数化三元组 $(\sE_1, \sE_2, \tau)$，其中 $\sE_1, \sE_2$ 是形式圆盘上的 $\cG$-联络，$\tau$ 是它们限制到穿孔圆盘之后的同构。注意，由于形式圆盘上每个 $\cG$-联络都可平凡化，底层的经典叠是 $\on{pt} / \cG$。然而这个模问题要求我们把 $\on{pt}/\cG$ 与它在 $\LS$ 中的自身相交，正是这提供了使该等价成立的导出结构。

因此，不正式地说，这个等价把 $G$ 的一个凝聚模空间上的构造性对象，与 $\cG$ 的类似构造性模空间上的凝聚对象互换。
\end{re}

\ssszh \label{s:abg}表中的下一个等价，即
\[
   \Dmod(L^+G \backslash LG / I) \simeq \ICn(\on{pt}/\cG \underset {\cfg / \cG}
   \times \widetilde{\sN} / \cG),
\]
也是 Arkhipov--Bezrukavnikov--Ginzburg \cite{abg} 一个非常好的定理。如注 \ref{r:paramcount} 那样，可以这样配比参数。在自守侧，$LG$ 的相关分层现在由层
\[
     L^+ G \cdot t^{\check{\mu}} \cdot I, \quad \text{对 } \check{\mu} \in
     \check{\Lambda},
\]
给出。也就是说，要把右侧的 $L^+G$ 缩小到 $I$，就必须乘以有限 Weyl 群这个因子，即从支配余特征标扩大到全部余特征标。

在谱侧，回忆 $\widetilde{\sN}$ 是（有限）旗簇的余切丛，我们考虑的是零截面（即旗簇本身）的导出加厚上的 $\cG$-等变凝聚层。这里，若忽略加厚，同一个格 $\check{\Lambda}$ 出现在
\[
   \QCoh( \Fl_{\cG}  / \cG  ) \simeq \QCoh(  \cB \backslash \cG / \cG) \simeq
   \QCoh( \cB \backslash \on{pt}) \simeq \on{Rep}(\cB),
\]
中，即作为 $\Fl_{\cG}$ 上的等变线丛，等价地作为 $\on{Rep}(\cB)$ 的特征标。

\ssszh 接下来把球与 Whittaker 两类范畴配对，得到
%
\[
   \Dmod(L^+G \backslash LG / LN, \psi) \simeq \QCoh( \pt / \cG).
\]
%
这同样是 Frenkel--Gaitsgory--Vilonen \cite{fgv} 的一条已知而且要重的定理。可以把它理解为 $p$-adic 表示论中
Casselman--Shalika 公式 \cite{cassshalika} 的几何精细化，该公式显式地
确定了一个非分歧表示的 Whittaker 余向量。

这里的参数配比类似。也就是说，$LG$ 被 double cosets
\[
   L^+G \cdot t^{\check{\mu}} \cdot LN, \quad \check{\mu} \in \check{\Lambda}
\]
分层，而一条轨道承载 Whittaker 层当且仅当其余权 $\check{\mu}$ 反支配。

\begin{re} 我们还想简短评注与双球情形（即导出 Satake）的一处基本差别。在那里，谱侧是
$\on{Rep}(\cG)$ 的一个导出增强。粗略地说，之所以如此，是因为
$L^+G \backslash LG / L^+G$ 上相交上同调层的茎沿着轨道闭包相当有趣，
对应着 $\cG$-表示的 $q$-权数重数 \cite{lusztigqchar}。

在当下情形，没有这样的导出增强。之所以如此，是因为
$L^+G \backslash LG / LN, \psi$ 上的相交上同调对象是 clean 的，
即在支撑中开轨道的边界处，其 $!$-茎与 $*$-茎都消没。

\end{re}

\ssszh \label{ss:sphericalkmmod}最后，把球向量与 Kac--Moody 表示配对，得到
\[
   \fgc\mod^{L^+G} \simeq \QCoh(\Op^{\on{unramified}}),
\]
其中后一个模空间按定义参数化平凡丛上的 oper 结构，即
\[
\Op^{\on{unramified}} \simeq \pt / \cG \underset \LS \times \Op.
\]

这同样是一条已知的定理，属于 Frenkel--Gaitsgory \cite{fg-sph}，它包含并
加强了先前关于临界水平 $L^+G$-等变 Kac--Moody 表示的许多结果，
这些结果属于 Feigin--Frenkel、Beilinson--Drinfeld 等人。

此时，Grothendieck 群不再是可数的。取而代之，单模构成由射有限
维仿射空间参数化的连续族。尽管如此，我们仍能给出如下粗略的
计数。

在自守侧，有一组紧生成元，即 Weyl 模；按定义，它们是有限维不可约
$G$-模的抛物诱导
\[
  \mathbb{V}_\lambda := \on{pind}_{\fg}^{\fgc}(L_\lambda), \quad \lambda \in \Lambda^+.
\]
在当下的李代数情形，抛物诱导由对应
\[
  \fg \leftarrow \fg[[t]] \rightarrow \fgc
\]
给出，也就是说，它是限制函子与诱导函子的复合
%
\[
\on{pind}_{\fg}^{\fgc}: \fg\mod \xrightarrow{\on{Res}} \fg[[t]]\mod \xrightarrow{\on{ind}} \fgc\mod.
\]
%

每个 $\mathbb{V}_\lambda$ 的自同态由中心 $Z(\fgc)$ 的一个商给出，并且它在
$\gkc\mod^{L^+G}_\lambda$ 中是投射的——这里 $\gkc\mod_\lambda$ 指 $Z(\fgc)$ 通过同一个商
作用的 $\gkc$-模范畴。这是一个非平凡的事实。

\begin{ex}\label{e:regop} 若 $\lambda = 0$，就得到临界水平处的真空
模\footnote{这里「真空」一词源于它在共形场论中的出处。也就是说，它是某个顶点算子代数的
底向量空间，而后者作为局部算子出现在具有 $G$-对称性的手征共形场论中。}
	\[
	  \mathbb{V}_0 \simeq \on{pind}_{\fg}^{\fgc}(k).
	\]
在 Feigin--Frenkel 同构之下，它的自同态典范地等同于非穿孔形式圆盘上
oper 代数上的函数代数
%
\[
    \Op^{\on{reg}} \simeq \{ d + f dt + \check{\fb}[[t]]dt \} / \cN[[t]].
\]
%
\end{ex}

在谱侧，随着 $\lambda$ 变化，这些自同态代数对应 $\Op$ 内部互不相交的闭
子概形。模空间 $\Op^{\on{unramified}}$ 是它们并集的形式加厚，而该定理进一步把这些
幂幺法方向与 Weyl 模的自扩张相匹配。

\ssszh\label{s:2realz} 可以进入图 \ref{f:tableaut} 与图 \ref{f:tablespec} 的下一行了，即
驯分歧与幂幺单值的情形。先遇到的是自同态，亦即一个幺半等价
%
\begin{equation} \label{e:2realz}
 \Dmod(I \backslash LG / I) \simeq \ICn( \widetilde{\sN} / \cG \underset {\cfg
 / \cG} \times \widetilde{\sN} / \cG).
\end{equation}
%
这一漂亮的陈述在多个意义上都值得注意。其一，有限 Hecke 范畴并没有这样的
凝聚实现，在这个意义上，该等价确实是仿射情形（具体地说，是其与算术的
联系）所特有的。\footnote{诚然，通过 Soergel 双模的构造可以把单值 Hecke
范畴写成某些凝聚层在
\[
\check{T} \underset {\check{T}/W_f}{\times} \check{T}
\]
上的同伦范畴，参见 \cite{bezrichetop}。不过这一等同在某种程度上更微妙，因为
它只在倾斜对象上、且仅在非导出意义下才全忠实；也就是说，相应的
Soergel 双模在双模范畴中有非平凡的 Ext。无论如何，\eqref{e:2realz} 还有另一个
令人愉快的特性：有限 Hecke 范畴的 Soergel 双模表述确实嵌入其中。}
如同我们迄今遇到的其它低分歧陈述，它是 Bezrukavnikov \cite{bez2realz} 的一条著名定理。

据我们所知，这一等价最初是作为猜想由 Ginzburg \cite{chrissginzburg} 提出的。一个主要
动机来自 $p$-adic 情形：在那里，$\on{Dist}(I \backslash G(\Fqt) / I)$（即仿射 Hecke 代数）
与 Steinberg 簇
\[
\widetilde{\sN} \underset \cfg \times \widetilde{\sN}
\]
的等变 $K$-理论之适当特化之间的类似同构，属于 Kazhdan--Lusztig \cite{kazhdanlusztig87}。他们以及
独立地 Ginzburg，用这一同构分类了仿射 Hecke 代数的不可约表示并确定了它们的特征标。等价地说，他们
确定了 $G(\Fqt)$ 的非分歧主系列的不可约子商的特征标。

来数参数。在自守侧，$LG$ 的相关分层其成分是
%
\[
    I\cdot w \cdot t^{\check{\mu}} \cdot I, \quad w \in W_f,\ \check{\mu} \in \check{\Lambda}.
\]
%
这是 $G$ 的 Bruhat 分解（双陪集 $B \cdot w \cdot B$，$w \in W_f$）的仿射类比，而这里的指标集
确实就是（扩张）仿射 Weyl 群 $W_f \ltimes \check{\Lambda}$。注意，如第 \ref{s:abg} 节那样，
我们再一次把 $L^+G$ 缩小到 $I$，因此分层多出了一个 $W_f$ 因子。

在谱侧，这些参数如下出现。首先，Steinberg 簇\footnote{虽然这是标准术语，我们要强调，
对于该幺半等价，我们实际用的是它的导出增强。}显式地参数化三元组
\[
(X, \cB', \cB''),
\]
其中 $X$ 是 $\cfg$ 的幂幺元，而 $\cB'$、$\cB''$ 是包含它的两个 Borel 子群。Steinberg 簇是
$\widetilde{\sN} \times \widetilde{\sN}$ 中 $W_f$ 个局部闭 Lagrange 子簇的并，其中每个
Lagrange 子簇由 $\cB_1$ 与 $\cB_2$ 的相对位置指定，即由 $\cB \backslash \cG / \cB$ 的一个点指定。

\begin{re} 回忆我们真正关心的是商
\[
\widetilde{\sN} / \cG \underset {\cfg/\cG} \times \widetilde{\sN} / \cG.
\]
把它认作 Hecke 叠
\begin{equation} \label{e:strestopesto} \cB \backslash \cG / \cB \simeq \Fl_{\cG} \times \Fl_{\cG} / \cG
\end{equation}
的余切（导出叠）会很有启发，这里的呈现方式是从 \eqref{e:strestopesto} 右端出发做 Hamilton 约化。
与此相关，所涉及的这些 $W_f$ 个成分不过是相应 Schubert 胞腔的余法丛。
\end{re}

如第 \ref{s:abg} 节那样，每个成分带有 $\check{\Lambda}$ 多个等变线丛，它们由沿任一投影
到 $\widetilde{\sN}/\cG$ 的拉回给出，\footnote{从 $\check{\Lambda}$ 到其 Picard 群的两个映射相差
$w$ 的相应元素的作用。}于是对成分求和，就得到所要的 $W_f \ltimes \check{\Lambda}$ 计数。

\ssszh 现在把非分歧主系列与 Whittaker 模型配对，得到等价
%
\begin{equation} \label{e:arkbez}
    \Dmod( I \backslash LG / LN, \psi) \simeq \QCoh( \widetilde{\sN} / \cG).
\end{equation}
%
这一断言同样是已知的，是 Arkhipov--Bezukravnikov \cite{ab} 的一条非常好的定理。他们的论证被广泛视为
本领域的一颗明珠。该结果在 Bezrukavnikov 对 \eqref{e:2realz} 的证明中起基本作用，
这与 $K$-理论中一条类似恒等式在 Kazhdan--Lusztig 把仿射 Hecke 代数等同于 Steinberg 簇
的等变 $K$-理论时所起的作用相平行。

这里的参数计数如下。在自守侧，过渡到半无限变体更方便，即
\begin{equation} \label{e:semiinfre}  \Dmod( LN \cdot L^+T \backslash LG / LN,
\psi).\end{equation}
在这一呈现下，按 $LG$ 的陪集
\[ (LN \cdot L^+T) \cdot w \cdot t^{\check{\mu}} \cdot LN,   \quad w
\in W_f,\ \check{\mu} \in \check{\Lambda} \]
分层，一条轨道承载 Whittaker 层当且仅当 $w$ 等于 $W_f$ 的最长元 $w_\circ$。特别地，相关轨道
的集合典范地由 $\check{\Lambda}$ 参数化。

在谱侧，有
%
\[
    \widetilde{\sN} / \cG = (\cG \overset \cB \times \check{\fn}) / \cG \simeq
    \check{\fn} / \cB,
\]
%
而 $\check{\Lambda}$ 再次对应于从 $\pt / \cB$ 拉回的等变线丛。

关于这一等价，有两点评注。

\begin{re} 在 $p$-adic 情形，相关的分布空间
\begin{equation} \label{e:antisph}\on{Dist}( I \backslash
G(\Fqt) / N(\Fqt), \psi)\end{equation}
是仿射 Hecke 代数的一个诱导表示。更准确地说，若以 $H$ 记（扩张）仿射 Hecke 代数，$H_f$ 记有限 Hecke
代数，$k_{\on{sgn}}$ 记其符号表示，则 \eqref{e:antisph} 就是所谓反球表示
\[ H \underset {H_f} \otimes k_{\on{sgn}}.\]
几何情形有类似的公式，即
\[ \Dmod(I \backslash LG / LN, \psi) \simeq \Dmod(I \backslash LG / I)
\underset {\Dmod(B\backslash G / B)} \otimes \Dmod(B \backslash G/ N, \psi),\]
参见 \cite{ahc}。
\end{re}

\begin{re} 我们在 \eqref{e:semiinfre} 中把 \eqref{e:arkbez} 的自守侧改写成了半无限形式。这是为了
稍稍简化相关轨道的组合学。不过，记住这一替代表述还有不那么琐碎的理由。也就是说，回忆范畴表示的
抛物诱导确实给出仿射旗簇的半无限版本，参见第 \ref{s:unramps} 节。与此相关，在
Arkhipov--Bezukravnikov 的证明中，起中心作用的正是半无限来源的对象，例如 Wakimoto 层。\end{re}

\ssszh \label{ss:iwahorikmmod}把非分歧主系列与 Kac--Moody 表示配对，得到等价
%
\[
 \fgc\mod^I \simeq \QCoh(\widetilde{\sN}/\cG \underset \LS \times \Op).
\]
%
这一断言是前面讨论的 Frenkel--Gaitsgory 球情形结果（见第 \ref{ss:sphericalkmmod} 节）的延伸。
他们先在 \cite{fg06loc} 中提出猜想，几年后又通过对仿射旗簇上临界水平局部化的工作得到它的一个版本，
其中还对中心特征作了进一步特化 \cite{fglocaff}。\footnote{严格地说，在 {\em loc. cit.} 中给出的是
Iwahori 单值变体。}

\ssszh 转到图 \ref{f:tableaut} 与图 \ref{f:tablespec} 的第三行，研究以 Whittaker 模型为源的缠绕算子。先遇到它的自同态，以及一个幺半范畴的等价
%
\begin{equation} \label{e:whitend} \Dmod(LN, \psi \backslash LG / LN, \psi)
\simeq \QCoh(\LS).\end{equation}
%
这是一个引人注目、且基本上仍完全开放的猜想。甚至连左端的对称幺半结构都远非显然。从精神上说，
在当前几何设定下，它是 Whittaker 余向量唯一性的一种体现——在通常的「重数自由」与「自同态代数交换」的通常对应之下。

我们强调，这一等价会确定猜想 \eqref{e:loclang} 中一大片内容，即谱侧所有奇异支撑沿零截面的
范畴层。这样的范畴表示称为缓增的。

\begin{re} 尽管 \eqref{e:whitend} 基本上仍是开放的，这里有一个已知的特例。也就是说，由 Raskin \cite{whit}
的工作，环路群的范畴表示上的深度滤过给出
\[
\Dmod(LN, \psi \backslash LG / LN, \psi)
\]
上的一个滤过，其中每个商都是幺半余局部化。第一个商对应「向量被 Iwahori 子群的
射幂幺根固定的表示」，即 Iwahori--Whittaker 范畴
%
\[
     \Dmod(\mathring{I}, \psi \backslash LG / \mathring{I}, \psi).
\]
%
它又容许一个进一步的商，后者对应从全部驯分歧表示过渡到幂幺单值的那些表示，即
Steinberg--Whittaker 范畴
%
\[
   \Dmod(\mathring{I}, \psi, \on{s} \backslash LG / \mathring{I}, \psi,
   \on{s}).
\]
%
显然，它就是由 $\Dmod(I^- \backslash LG / I^-)$ 中对象的平均所生成的满子范畴，其中 $I^-$ 是
$L^+G$ 中相对于与 $B$ 处于一般位置的 Borel 子群 $B^-$ 定义的 Iwahori 子群，参见 \cite{ahc}。
若以 $\sN^{\wedge}$ 记 $\cfg$ 沿幂幺锥 $\sN$ 的形式完备化，则有一个幺半等价
%
\[
    \Dmod(\mathring{I}, \psi, \on{s} \backslash LG / \mathring{I}, \psi,
    \on{s}) \simeq \QCoh( \sN^\wedge / \cG),
\]
%
这正是 \eqref{e:whitend} 所许诺的一块。

Bezrukavnikov 证明仿射 Hecke 范畴的等价 \eqref{e:2realz}，本质上就是把这一结果与前面提到的
Arkhipov--Bezrukavnikov 结果结合起来。更准确地说，以 $\widetilde{\sN}^\wedge$ 记
Grothendieck--Springer 簇
\[
\widetilde{\cG} := \cG \overset \cB \times \check{\fb}
\]
沿 Springer 簇的形式完备化。也就是说，Grothendieck--Springer 簇参数化数对 $(X, \check{\fb}')$，
其中 $X$ 是 $\cfg$ 的元素，$\check{\fb}'$ 是包含它的 Borel 子代数；而沿 $X$ 幂幺的轨迹作完备化。
于是所要的等价 \eqref{e:2realz}\footnote{更仔细地说，我们指的是它的单值变体；严格的陈述可以从它
形式地推出。}由重整化下述等价得到
%
{\small
\begin{align*} &\QCoh( \widetilde{\sN}^\wedge / \cG  \underset {\cfg / \cG}
\times
\widetilde{\sN}^\wedge / \cG) \\
&\qquad \simeq \QCoh(\widetilde{\sN}^\wedge / \cG)  \underset {\QCoh(\sN^\wedge / \cG)}
\otimes  \QCoh(\widetilde{\sN}^\wedge / \cG ) \\ &\qquad \simeq \Dmod(\mathring{I}
\backslash LG / \mathring{I}, \psi, \on{s}) \underset {\Dmod(\mathring{I},
\psi, \on{s} \backslash LG / \mathring{I}, \psi, \on{s})} \otimes
\Dmod(\mathring{I}, \psi, \on{s} \backslash LG / \mathring{I}).  \end{align*}}
%

\end{re}

\ssszh 接下来把 Whittaker 模型与 Kac--Moody 表示配对，得到等价
\begin{equation} \label{e:affskr} \fgc\mod^{LN, \psi} \simeq \QCoh(\Op).
\end{equation}
这同样是一条关于任意分歧对象的强断言。与我们正遇到的其他这类断言不同，它事实上是 Raskin
\cite{whit} 的一条定理。

也就是说，Raskin 的仿射 Skryabin 定理\footnote{附带一提，Skryabin 是「by Raskin」的变位词。}
对任意水平 $\kappa$ 给出一个典范等价
%
\[
  \fgk\mod^{LN, \psi} \simeq \sW\mod,
\]
%
其中 $\sW$ 是仿射 W-代数。另一方面，在临界水平，Feigin--Frenkel 已把 $\mathscr{W}_{\kappa_c}$ 等同于
包络代数的中心 $Z(\fgc)$ \cite{feiginfrenkel91}，因而也等同于 $\OO_{\Op}$。这与仿射 Skryabin 定理
结合，就给出 \eqref{e:affskr}。

\begin{re} 读者可参考本文集中 Arakawa 的综述，以了解 W-代数的更多细节；另见 \cite{araintro}、\cite{fbz}。
这里只附带一提：仿射 Skryabin 定理使得人们可以在 W-代数表示的研究中系统地使用范畴表示论的
技巧。例如可参见 \cite{locw}，其中讨论了 W-代数的最高权模的局部化理论。
\end{re}

\ssszh \label{ss:affhc}最后一行遇到的是 Kac--Moody 表示的自同态。它们如下等同于临界水平处的
仿射 Harish-Chandra 双模。

首先，范畴 $\fgc\mod$ 作为 $\Dmod(LG)$ 表示、特别地作为抽象 dg 范畴，典范地自对偶。

为看清这一点，从更简单的情形入手会有帮助。对带中心扩张的有限维李代数 $\mathfrak{a}$
%
\[
0 \rightarrow k \cdot \mathbf{1} \rightarrow \widetilde{\mathfrak{a}}_c
\rightarrow \mathfrak{a} \rightarrow 0,
\]
%
考虑 $\mathbf{1}$ 通过恒等作用的 $\widetilde{\mathfrak{a}}_c$-模范畴 $\widetilde{\mathfrak{a}}_c\mod$。
若以 $\widetilde{\mathfrak{a}}_{-c}$ 记相反的中心扩张，例如由任一表示 2-余循环的加法逆给出，则有一个
完美配对
%
\[
\widetilde{\mathfrak{a}}_c\mod \otimes \hspace{.7mm}
\widetilde{\mathfrak{a}}_{-c}\mod \rightarrow \mathfrak{a}\mod \rightarrow
\Vect,
\]
%
其中第一个箭头把表示张量起来以消去扩张，第二个箭头是李代数同调。

在仿射情形，Arkhipov--Gaitsgory \cite{ag} 在 Frenkel--Gaitsgory \cite{fg06loc} 先前工作的基础上证明了：
对任意 $\kappa$，复合
%
\[
\fgk\mod \otimes \hspace{.7mm} \widehat{\fg}_{-\kappa + 2\kappa_c}\mod
\rightarrow \widehat{\fg}_{2 \kappa_c}\mod \rightarrow \Vect
\]
%
是完美配对，其中第一个箭头把表示张量起来，第二个箭头是半无限上同调函子。简言之，后者是像 $\gk$
这样的无限维李代数（正式地说，Tate 李代数）的表示的自然同调理论，但它需要水平为 $2\kappa_c$ 的中心
扩张。对一般 Tate 李代数也有这样的对偶陈述，参见 \cite{mys}、\cite{lpw}。

特别地，有自对偶
%
\[
    \fgc\mod \otimes \hspace{.7mm} \fgc\mod \rightarrow \Vect,
\]
%
这是「临界水平是水平空间的中点」的另一种意义。此外，如第 \ref{ss:hchbimod} 节那样，这把
$\fgc\mod$ 的 $\Dmod(LG)$-等变自函子等同于临界水平仿射 Harish-Chandra 双模范畴
%
\[
      \fgc \oplus \fgc \mod^{LG}.
\]
%
这是一个漂亮却也多少难以把握的范畴。例如，它的全部对象相对于 $\fgc \oplus \fgc\mod$ 上自然的
$t$-结构都落在同调次数 $-\infty$，也就是说，它由通过重整化引入的同调幽灵构成。

\begin{re} 其他水平处的仿射 Harish-Chandra 双模同样微妙，其中一些基础工作与猜想属于 I. Frenkel--Malikov
\cite{frenkelmalikov1} 与 Gaitsgory \cite{quantum-langlands-summary}。
\end{re}

对这些双模，局部几何朗兰兹预测一个与谱卷积代数的幺半等价
%
\[
      \fgc \oplus \fgc \mod^{LG} \simeq \QCoh( \Op \underset \LS \times \Op).
\]
%
这是一个引人注目、且基本上仍完全开放的猜想。据我们所知，甚至连球情形与幂幺单值的驯分歧情形都还是
开放的；不过我们预期，利用前面提到的 Frenkel--Gaitsgory 的一些工作（参见第 \ref{ss:sphericalkmmod}
与 \ref{ss:iwahorikmmod} 节），再结合临界水平的局部化定理，这两个情形是可以达到的。下面转入该定理。

\subsection{局部化理论}
	\label{ss:locaff}
	\ssszh 我们想讨论局部几何朗兰兹的想法与技巧在仿射李代数表示论中的一个重要
	应用，而这一问题是纯代数的。先交代背景。

	\ssszh 回忆第 \ref{s:bbloc} 节讨论了
	约化李代数的局部化理论，并概述了如何用它确定
	单最高权模的特征标。

	对仿射李代数，Kashiwara--Tanisaki 在一系列工作中于非临界水平发展了类似的纲领
	\cite{kashiwara90}、\cite{kashiwaratanisakiposlevelii}、\cite{kashiwaratanisakinegativeintegrallevel}、
	\cite{kashiwaratanisakinegativerationallevel}、\cite{kashiwaratanisakipositiverational}、
	\cite{kashiwaratanisakialllevels}。扼要地说，在负水平，他们把仿射旗流形上的 D-模
	与某些 Kac--Moody 表示等同起来。在正水平，他们做了同样的事，但把仿射旗流形换成
	厚仿射旗流形，后者是无限型的概形，参数化 $\mathbb{P}^1$ 上在零点处有完整水平结构、在无穷远处有 Borel
	约化的 $G$-丛。作为推论，他们得到了非临界水平单最高权模的特征标。

	在他们工作之后，人们认识到，要推进剩下的情形，即临界水平的表示论，还需要新的想法与工具。下面我们
	想按 Beilinson--Drinfeld、Frenkel--Gaitsgory 与 Raskin--Yang 的工作，概述这里的局部化理论是什么样子。

	\ssszh 临界水平处的基本新困难在于：局部化定理的几何侧不再是某个空间上的 D-模，而是它的一种修正。
	令人愉快的是，只用我们迄今讨论过的局部几何朗兰兹内容，就能得到正确的陈述。下面来看怎么做到。

	首先，回忆对任一 $\cG$-oper $\chi$，设其底层的局部系统为 $\sigma$，局部几何朗兰兹预测：在映射
	\[
	  \sigma: \pt \rightarrow \LS
	\]
	之下对应于 $\Vect \simeq \QCoh(\pt)$ 的范畴环路群表示 $\sC_{\sigma}$，可以由重整化
	%
	\[
	  \fgk\mod \underset{\QCoh(\Op)} \otimes \Vect
	\]
	%
	得到，也就是取中心特征为 $\chi$ 的 Kac--Moody 表示。

   \begin{re} 既然我们已经讨论了范畴的归纳凝聚层，这一重整化可以理解为映射
   \[
   \Op \rightarrow \LS
   \]
   形式光滑的推论。特别地，这意味着由 $\QCoh(\Op)$ 生成的子 2-范畴只由拟凝聚范畴层构成。
   也就是说，没有空间容纳非平凡的奇异支撑。

   在朗兰兹对偶侧，这预测 $\fgc\mod$ 是缓增的，即典范映射
   %
   \[
        \Dmod(LG/ LN, \psi ) \underset {\Dmod(LN, \psi \backslash LG / LN,
        \psi)} \otimes \mathscr{W}_{\kappa_c}\mod \rightarrow \fgc\mod
   \]
   %
   是等价。这是关于临界水平 Kac--Moody 表示的一个值得注意的猜想。
   \end{re}

	\ssszh 把上述讨论（它对任意局部系统 $\sigma$ 都适用）特化到平凡联络。为此，可以平凡地把
	$\sigma$ 分解为
	%
	\[
	  \pt \rightarrow \pt/\cG \rightarrow \LS.
	\]
	%

	我们已经相当详细地讨论过 $\QCoh(\pt/\cG)$ 在局部几何朗兰兹下应当对应什么，即仿射
	Grassmann 簇上的 D-模。这给出 $\sC_{\sigma}$ 的另一种呈现，如下。由于 $\pt / \cG \rightarrow \LS$
	尽可能不光滑，$\Vect$ 落在由 $\QCoh(\pt / \cG)$ 生成的范畴之中。于是得到

	%
	{\footnotesize
	\begin{align*}
	     \Vect &\simeq  \QCoh( \pt / \cG) \\
	     &\quad \otimes_{\Hom_{\TwoIC(\LS)}(\QCoh(\pt / \cG), \QCoh(\pt / \cG))} \\
	     &\qquad \Hom_{\TwoIC(\LS)}( \QCoh(\pt / \cG), \QCoh(\pt)) \\
	     &\simeq
	     \QCoh( \pt / \cG) \otimes_{\ICn(\pt / \cG \underset \LS \times \pt /
	     \cG)} \ICn( \pt \underset \LS \times \pt ).
	\end{align*}}
	%
因此在自守侧，应当有
\begin{align}
   \sC_{\sigma} &\simeq \Dmod(\Gr_G)\underset {\ICn(\pt / \cG \underset \LS
   \times \pt /
   	\cG)} \otimes \ICn( \pt \underset \LS \times \pt ).
\intertext{而从 $\ICn( \pt / \cG \underset \LS \times \pt / \cG)$ 张量到
$\ICn(\pt
\underset  \LS \times \pt)$ 不过是忘掉 $\cG$ 作用的等变性，所以可以把它更直白地改写成 }
 \sC_{\sigma} & \simeq \Dmod(\Gr_G) \underset{\QCoh(\pt / \cG)} \otimes
 \QCoh(\pt).\label{e:heckeeig}
	\end{align}

显式地说，这就是仿射 Grassmann 簇上的 Hecke 特征层范畴。
也就是说，回忆 $\QCoh( \pt / \cG)$，即 $\on{Rep}(\cG)$，等同于阿贝尔 Satake 范畴
\[
\Dmod(L^+G \backslash LG /
L^+G)^\heartsuit.
\]
该范畴通过右卷积作用在仿射 Grassmann 簇的 D-模上，这一作用记为
%
\[
    - \overset {L^+G} \star -: \Dmod(LG/L^+G) \otimes \Dmod(L^+G \backslash LG
    / L^+G) \rightarrow \Dmod(LG / L^+G).
\]
%
所要的范畴由这一作用的特征对象构成，也就是说，由 $\Dmod(\Gr_G)$ 中的层 $\sF$ 构成，
并带有相容的同构
%
\begin{equation} \label{e:heckeeigen}
    \sF \overset {L^+G }\star \on{Sat}(V) \simeq \sF \otimes \on{Oblv}(V),
    \quad V \in \on{Rep}(\cG),
\end{equation}
%
其中 $\on{Oblv}(V)$ 记 $V$ 的底向量空间。关于这样的 $\sF$——与任一 Satake 层作卷积都返回
$\sF$ 的若干拷贝的直和——有一个基本观察：$\sF$ 必须有无限维的支撑，特别地，它不可能是
$\Dmod(\Gr_G)$ 的紧对象。

在比较 $\sC_\sigma$ 的两种描述之前，再作两条帮助定位的评注。

\begin{re}  把 Satake 层拖过张量积 \eqref{e:heckeeig}，也就是说，考虑
%
\[
     (\sF \overset {L^+G} \star V) \otimes k \simeq \sF \otimes (V \cdot k)
     \simeq
     \sF \otimes \on{Oblv}(V),
\]
%
其中 $\cdot$ 记 $\QCoh(\pt/\cG)$ 在 $\QCoh(\pt)$ 上的作用，就得到 Hecke 特征层更接地的表述
\eqref{e:heckeeigen}。
	\end{re}

	\begin{re} 初次接触这些想法的读者，可能觉得前面讨论的一些细节稍显复杂。因此值得强调，
	也可以通过类比来预测这一点。也就是说，在 $p$-adic 情形，要从
	%
	\[
	  \Fun( G(\Fqt)/G(\Fq[[t]]))
	\]
	%
	中提取 $G(\Fqt)$ 的一个不可约非分歧表示 $\pi_{\eta}$，人们取球 Hecke 代数右作用的特征向量，
	特征值为 $\eta$，即取相应的同态
	%
	\[   \eta: \Fun(G(\Fq[[t]]) \backslash G(\Fqt)/G(\Fq[[t]]))) \rightarrow k.\]
	%
	 在几何设定下，由 Tannaka 对偶，可能的特征值只有一个，即幺半函子
	 %
	 \[ \on{Oblv}: \on{Rep}(\cG) \rightarrow \Vect,\]
	 %
	  而我们再次是从万有表示出发、通过取特征向量来构造不可约非分歧表示。特别地，上面关于
	 Hecke 特征层支撑的观察，提升为关于 Hecke 特征向量支撑的类似断言。
	\end{re}

	\ssszh 比较非分歧情形下 $\sC_{\sigma}$ 的两种构造，就得到如下预测。也就是说，对平凡联络上的
	任一 oper 结构 $\chi$，应当有等价
	%
	\begin{equation} \label{e:locaff}
	      \Dmod(\Gr_G) \underset {\QCoh(\pt / \cG)} \otimes \QCoh(\pt) \simeq
	      \fgc\mod_\chi,
	\end{equation}
	%
	其中右端记中心特征为 $\chi$ 的 Kac--Moody 表示范畴的适当重整化。

	事实上，这是一条重要的定理，\footnote{更仔细地说，至少当 $\chi$ 对应形式穿孔圆盘上的 oper 时，
	它是定理，参见例 \ref{e:regop}。不过人们预期一般情形可以类似处理。}它的发展历时数十年，
	涉及多组作者的重要贡献。简要的历史如下。在对 Hitchin 系统的量子化的工作 \cite{bdh} 中，
	Beilinson--Drinfeld 证明了：在卷积作用
	%
	\[
	     \Dmod(L^+G \backslash LG / L^+G) \otimes \fgc\mod^{L^+G} \rightarrow
	     \fgc\mod^{L^+G}
	\]
	%
	之下，真空模的任一中心商都是阿贝尔 Satake 范畴作用的特征对象。这条定理是他们在整体几何朗兰兹中
	构造对应于整体 oper 的 Hecke 特征层时所用的主要局部输入。

	\ssszh Beilinson--Drinfeld 的工作对临界水平局部化理论的意涵，由 Frenkel--Gaitsgory 认识到，
	并在 \cite{frenkelgaitsgorylocgrassmannian04}、\cite{fg06loc} 与 \cite{fgloc09} 中大幅发展。首先，
	Frenkel--Gaitsgory 用 Beilinson--Drinfeld 建立的 Hecke 特征性质构造了一个函子
	%
	\begin{equation} \label{e:globgr}
	      \Dmod(\Gr_G)\underset {\QCoh(\pt / \cG)} \otimes \QCoh(\pt)
	      \rightarrow
	      \fgc\mod_\chi.
	\end{equation}
	%
	粗略地说，这相当于说明：$\Gr_G$ 上任一 D-模（以临界线丛扭之，从而得到临界水平的表示，并
	施加对应于 $\chi$ 的中心商）的整体截面自然地是 Hecke 特征对象。对 delta D-模，这正是
	Beilinson--Drinfeld 前述的定理，而一般情形可以化归到它。

	 除猜想 \eqref{e:globgr} 应是 $t$-正合的等价之外，Frenkel--Gaitsgory 还确立了许多基本性质。也就是说，
	 他们证明了它右 $t$-正合，且在下有界导出范畴上是全忠实的嵌入。\footnote{一般情形下的全忠实性，
	 特别是 Kac--Moody 表示范畴的正确重整化，后来由 Raskin \cite{RaskinLocCrit} 得到。}他们进一步证明了
	 它在 $I$-等变对象、即最高权模上是等价。

	\begin{re} 稍作加工，就可以从这些结果得到临界水平单最高权模特征标的 Kazhdan--Lusztig 猜想的类比。
	这一猜想由 Feigin--Frenkel 提出，见 \cite{arakawafiebig}，并在 \cite{DhillonYang} 中得到证明。
	\end{re}

\ssszh 在 Frenkel--Gaitsgory 的工作之后，其猜想剩下的部分是：\eqref{e:globgr} 是 $t$-正合的，且本质满。
后者可以有益地重新解释为环路群范畴作用下的一条生成性陈述，如下。\eqref{e:globgr} 的左端按定义由
其 $L^+G$-等变对象生成。因此，只需说明右端也如此。

一方面，这听起来是相当自然的陈述。确实，所施加的中心特征是非分歧性质的，对应形式非穿孔圆盘上的
oper，因此范畴本身应当是非分歧的，这颇为合理。然而，这一断言结果很难，很大程度上是因为人们对
驯分歧之外的 Kac--Moody 表示了解得很少，而且至今仍然如此。

\ssszh 十多年后，Raskin \cite{RaskinLocCrit} 取得了重要进展：他在第一个非平凡情形 $GL_2$
证明了该猜想，事实上对半单秩为一的任何群都证明了。为此，他证明了关于 $LGL_2$ 的范畴表示的一条
一般性断言，即它们由 Whittaker 不变量与 Iwahori 不变量生成；这与 $p$-adic 理论中一条更早、
同样基本的断言相平行。既然 \eqref{e:globgr} 在这类对象上是等价，完整的等价便随之而来。

猜想的一般情形随后由 Raskin--Yang \cite{RaskinYang} 解决。主要的新要素是 Yang \cite{yangmoyprasad}
把 Moy--Prasad 理论从 $p$-adic 情形移植过来；另见更早的 \cite{chenkamgarpour}。关键的是，除按深度
滤过环路群表示之外，Moy--Prasad 理论还为每个相继的商给出了一份紧生成元的显式清单，对应 $p$-adic
理论中的非精化极小 $K$-型。他们的证明随后通过一项非平凡的分析说明：从深度零以上的任一这样的生成元
出发若有非零映射，就与非分歧的中心特征不相容；于是得到所要的本质满性。

\begin{re} 事实上，前面讨论中我们略过了一个关键的技术要点：由于重整化，范畴 $\fgc\mod_\chi$ 并不
先验地带有 $LG$ 的范畴作用。这是 \cite{RaskinYang} 中克服的基本困难之一。
\end{re}

\begin{re} 我们想就 Moy--Prasad 理论在几何朗兰兹中的地位作一条评注，它与我们对局部化理论的讨论
本身无关。也就是说，由 Bremer--Sage \cite{sage1}、\cite{sage2}、\cite{sage3}，与 Galois 表示相对，
使用 $\cG$-联络有一个有趣的特点：对 $\cG$-联络也有极小 $K$-型理论。它与局部几何朗兰兹、特别是
与自守侧 Moy--Prasad 理论的相互作用，目前尚未被完全理解。	\end{re}

\ssszh 最后用两点评注结束关于临界水平局部化理论的讨论。第一，读者可以想见，上述内容可以从
非分歧情形推广到驯分歧情形，例如幂幺单值的情形。这由 Frenkel 与 Gaitsgory 在 \cite{fg06loc} 与
\cite{fglocaff} 中发展，做法大致是从仿射 Grassmann 簇过渡到仿射旗簇。

第二，颇为引人注目的是：尽管我们讨论的临界水平局部化定理已经很不平凡，它在某种意义上只是冰山
一角。也就是说，人们预期 $p$-adic 群更深层次的不可约表示的许多已知构造都应有几何版本，从而得到
等价，特别是带有野分歧中心特征的临界水平表示的局部化定理；相关讨论与猜想可参见 \cite{fg06loc} 与
\cite{kamgarpourschedler}。

特别地，综述作者预期：若一个 oper 的底层局部系统带有 $\cB$-约化，则相应的 Kac--Moody 模范畴
局部化为半无限旗流形上的扭 D-模范畴。对分歧特征标，这样的局部化不再能像 \cite{fglocaff} 那样在
仿射旗簇上实现，而要在先由 Roche \cite{roche} 在 $p$-adic 情形引入的某个更小的紧开子群的商上实现。

	\subsection{同类猜想} \label{ss:variants}
	本节讨论了局部几何朗兰兹猜想，希望读者对其陈述、以及它在表示论与
	几何中的一些具体预测（既有已知的，也有猜想的）有了基本的感受。

	在结束之前，我们应当提到附近领域中一些重要的变体，或者更确切地说，同类的猜想，以及
	一些其他正在浮现的联系。

	\ssszh[拟分裂群] 首先要作的评注是：本综述从头到尾使用的都是环路群，即 $k(\!(t)\!)$ 上的分裂
	约化群。然而，即便为了研究它们的范畴表示，正如数论中那样，考虑扭环路群的情形也很重要，
	即 $p$-adic 情形中拟分裂群的范畴类比。注意这也为扭仿射李代数的表示论提供了类似的图景，而
	下面讨论的局部几何朗兰兹对偶的各种形式，预期同样可以推广到拟分裂情形。

	\ssszh[量子局部几何朗兰兹] 回忆 $\fgk\mod$
	只是对整水平才是环路群的范畴表示。对其他水平，它反而带有扭 D-模范畴 $\Dmod_\kappa(LG)$ 的作用。
	可以问：对环路群的这类扭范畴表示，是否存在某种形式的局部朗兰兹对偶。确实存在，这正是量子几何
	朗兰兹的主题。

    为继续下去，为简单起见假设
    $\kappa - \kappa_c$ 是非退化双线性型。此时可以给 $\cG$ 配上一个对偶型 $\check{\kappa}$，如下。若以
    $\kappa_{\cfg, c}$ 记 $\cfg$ 的临界水平，则 $\check{\kappa}$ 由下述性质定义：
    \[
\kappa - \kappa_c \text{ 与 } \check{\kappa} - \kappa_{\cfg, c}
\]
    分别是 $\ft$ 与 $\check{\ft}$ 上的对偶双线性型。也就是说，在差一个我们反复遇到的临界平移之后，
    要求 $\kappa$ 与 $\check{\kappa}$ 互为对偶型。

    由 Gaitsgory \cite{quantum-langlands-summary} 提出的局部量子朗兰兹猜想是：存在 $(\infty,2)$-范畴的
    等价
    \[
\Dmod_\kappa(LG)\mod \simeq \Dmod_{-\check{\kappa}}(L\cG)\mod.
\]
    这建立在整体情形中 Beilinson--Drinfeld、Feigin--Frenkel 与 Stoyanovsky \cite{stoyanovsky} 的
    早期工作之上。

     这里只作几点评注；更多细节可参见 \cite{winter-langlands-summary}、\cite{paris-notes}。

    \begin{re} 关于量子情形，一个基本要点是：两端看起来更对称，因为它们只涉及扭的范畴环路群
    表示。此外，在确切的意义下，当 $\kappa \rightarrow \kappa_c$、从而 $\check{\kappa} \rightarrow \infty$ 时，
    就恢复通常的陈述 \eqref{e:loclang}；而当 $\check{\kappa} \rightarrow \kappa_{\cfg, c}$、从而
    $\kappa \rightarrow \infty$ 时，就恢复把 $G$ 与 $\cG$ 的角色互换后的类比。
    	\end{re}

    \begin{re} 「量子」这一形容词的来源如下。回忆 Feigin--Frenkel 同构
    %
    \[      \mathscr{W}_{\fg, \kappa_c} \simeq \OO_{\Op},
    \]
    %
    它在临界水平的故事中起基本作用。事实上，Feigin--Frenkel 在 \cite{feiginfrenkel91} 中证明了：
    它是仿射 W-代数的一族同构
    %
    \[   \mathscr{W}_{\fg, \kappa} \simeq \mathscr{W}_{\cfg, \check{\kappa}} \]
    %
    的准经典极限。这些同构（如今是非交换代数之间的同构）是前一个同构的形变量子化，在量子理论中
    起基本作用。

    与此相关，在临界水平我们看到对偶群 $\on{Rep}(\cG)$ 的表示在几何朗兰兹的非分歧情形中起基本作用；
    而在量子设定下，起类似作用的是量子群 $\on{Rep}_q(\cG)$ 的表示。
    \end{re}

   \ssszh[与物理的联系] 迄今为止，我们只触及了几何朗兰兹与物理之间的一些关系，而且大多是与二维
   共形场论的联系。然而，它与三、四维超对称规范理论之间有极其丰富的联系，这始于
   \cite{kapustinwitten}、\cite{wittenwild}、\cite{gukovwitten}、\cite{frenkelwitten} 等文。这是（量子）几何朗兰兹中
   许多有趣构造与预期相容性的来源，无论局部还是整体。部分近期工作可参见
   \cite{gaiotto16s-duality}、\cite{gaiottorapcak}、\cite{bfn}、\cite{gaiottotwistcompact}、\cite{costellogaiotto}、
   \cite{creutziggaiotto}、\cite{creutziggaiottolinshaw}、\cite{frenkelgaiotto}、\cite{bfgt}、\cite{hilburnraskin}，
   以及即将发表的 \cite{gaiottoyu}。

 \begin{re}
   上述部分工作的一个令人兴奋之处，是它与算术朗兰兹纲领的新进展密切相关。也就是说，由四维
   $\mathscr{N} = 4$ Yang--Mills 的边界条件的 S-对偶给出的朗兰兹对偶范畴的许多预测
   （Gaiotto--Witten \cite{gaiottowitten} 研究过），与算术中的类似预测相符，尤其是
   Sakellaridis--Venkatesh \cite{sakellaridisvenkatesh} 的相对朗兰兹纲领。这一联系在
   Ben-Zvi--Sakellaridis--Venkatesh 正在进行的工作中得到发展。
\end{re}

	\ssszh 我们来到最后一点，先略微非正式地陈述。在代数簇上使用构造层时，有几类不同的层可以
	选择，它们表现相似。显然，人们通常使用其中之一：（i）D-模，其相应的上同调理论（即「常层」的整体
	截面）是代数 de Rham 上同调；（ii）$\ell$-进层，其相应的上同调理论是 étale 上同调；（iii）在复数上
	工作时，使用解析拓扑下的构造层，其相应的上同调理论是奇异上同调，即 Betti 上同调。

	在我们讨论的（量子）几何朗兰兹版本中，对应的构造部分属于 de Rham 型。也就是说，在自守侧我们研究
	$G$-丛模空间上的 D-模，在谱侧研究带平坦联络的 $\cG$-丛模空间（即 de Rham 意义下的局部系统）上的
	凝聚层。与此相关，对 $k$ 上光滑射影曲线 $X$，Arinkin--Gaitsgory \cite{arinkingaitssingsupp} 用现代
	方式表述的整体猜想断言等价
	%
	\[
	\Dmod(\Bun_G(X)) \simeq \ICn(\LS(X)).
	\]
	%
	正如我们现在要描述的，本理论也有 étale 与 Betti 风味的版本。

	\ssszh[受限变化的局部几何朗兰兹] 局部算术朗兰兹的近期重要进展，是把它在表示的族中表述，
	即讨论整个范畴，而不只是一个一个不可约表示
	\cite{emertonhelm}、\cite{hellmann}、\cite{zhufamilies}、\cite{farguescholze}。也就是说，其现代表述
	大致是：$G(\Fqt)$ 的全部表示与某个适当的 Galois 表示模空间上的凝聚层之间的等价。

	\begin{re} 事实上，本着 Zelevinsky、Lusztig 与 Vogan \cite{voganloclang} 工作的精神，自守侧更
	微妙，它还要与其它群（例如 $G$ 的某些内形式）的表示范畴粘合起来。
	\end{re}

    在系列工作 \cite{restrictedvar1}、\cite{restrictedvar2}、\cite{restrictedvar3} 中，
    Arinkin--Gaitsgory--Kazhdan--Raskin--Rozenblyum--Varshavsky
    给出了一个在特征零与函数域情形下都说得通的整体范畴几何朗兰兹猜想。其形式是
    %
    \begin{equation} \label{e:etalelang}
           \on{Shv}^{\text{étale}}_{\on{nilp}}(\on{Bun}_G(X)) \simeq
           \on{IndCoh}_{\on{nilp}}(\on{LocSys}_{\cG}^{\on{restr}}(X)).
    \end{equation}
%
	在左端，考虑的是整体曲线 $X$ 上 $G$-丛模空间上的 ind-构造 $\ell$-进层，其奇异支撑落在
	$T^* \Bun_G(X)$ 的整体幂幺锥中。在右端，遇到的是 $X$ 上 $\ell$-进 $\cG$-局部系统的某个适当定义的
	模空间上、带幂幺奇异支撑的归纳凝聚层。

	\begin{re}
	非常粗略地说，基本图景是：在复数上，谱侧看起来像是 $\cG$-局部系统的 Betti 模叠的「撕开」版本，
	即参数化拓扑基本群表示的模空间。也就是说，按局部系统的半单化把 Betti 叠分解为许多块，
	再取它们形式邻域的不交并。例如，不可约局部系统的形式邻域在 $\on{LocSys}_{\cG}^{\on{restr}}$ 中是孤立的，
	并等同于 Betti 或 de Rham 模空间中相应的形式完备化。与此相关，上标「restr」代表受限变化。
	\end{re}

	在函数域设定下，取 Frobenius 的范畴迹，可以证明这恢复出（紧支撑）自守形式，特别地，
	恢复出函数域上尖点表示的非分歧整体朗兰兹的自守到 Galois 方向。

	在未发表的工作中，Gaitsgory 表述了 \eqref{e:etalelang} 的局部版本。它大致断言等价
	%
	\[
	     \on{Shv}^{\text{étale}}(G(\overline{\mathbb{F}}_q(\!(t)\!))\mod \simeq
	     \TwoIC(\LS^{\on{restr}}),
	\]
	%
	其中左端实际使用的是更微妙的范畴表示概念，以修正全忠实嵌入
	\[
\on{Shv}^{\text{étale}}(X) \otimes \on{Shv}^{\text{étale}}(Y) \hookrightarrow \on{Shv}^{\text{étale}}(X \times Y)
\]
	并非本质满的问题，并进一步过渡到「谱有限」对象的满子范畴。

	通过 Frobenius 的范畴迹，预期在函数域设定下这能恢复 Zhu \cite{zhufamilies} 所表述的算术局部
	朗兰兹猜想的某种形式。

	\ssszh[局部 Betti 几何朗兰兹] 最后，Ben-Zvi--Nadler \cite{benzvinadlerbetti} 表述了关于 Betti 层的
	几何朗兰兹版本。
	此时，对复数上的射影代数曲线 $X$，整体等价取形式
	%
	\[
	  \on{Shv}^{\on{Betti}}_{\on{nilp}}(\Bun_G(X)) \simeq
	  \ICn(\LS^{\on{Betti}}(X) ).
	\]
	%
	在左端，考虑的是 $\Bun_G(X)$ 的解析化上的 Betti 层（不必是 ind-构造的），其奇异支撑落在整体
	幂幺锥中。在右端，考虑的是局部系统的 Betti 模空间上的归纳凝聚层，即把 $X$ 的底拓扑空间的特征簇
	自然地增强为导出叠。

	局部情形没有很好的文献记载，但我们可以概述一些基本部分。首先，正则奇点的部分由等价
	%
	\begin{equation} \label{e:bezfam}
	      \on{Shv}^{\on{Betti}}_{\on{nilp}}( \mathring{I}\backslash LG /
	      \mathring{I}) \simeq \ICn( \widetilde{\cG} / \cG \underset {\cG /
	      \cG} \times \widetilde{\cG} / \cG).
	\end{equation}
	%
	确立。非正式地说，这是 Bezrukavnikov 等价 \eqref{e:2realz} 及其「非幂幺单值、且单值在族中变化」
	的变体的一个版本，其中单值在伴随商 $\check{G} / \check{G}$ 中变化。\footnote{这一等价可能尚未见于
	文献。不过综述作者与 H. Chen 有一个论证，希望不久能记录成文。}

	按构造，\eqref{e:bezfam} 右端的模等同于 $\LS^{\on{Betti}}$，即 $\cG / \cG$ 上带幂幺奇异支撑的
	归纳凝聚范畴层。完整的断言（即包括野分歧的对应物）应当把环路群的 Betti 范畴表示（一阶近似下
	即 $\on{Shv}^{\on{Betti}}(LG)\mod$）与 Stokes 数据模空间上带幂幺奇异支撑的归纳凝聚范畴层联系起来。

	\ssszh 至此，本综述已到结尾。希望读者对本领域的一些基本想法、结果与猜想有所感受，也能欣赏它们的
	美。有兴趣的读者不妨深入我们巡览过的文献；在一个有如此多迷人开放问题的领域中，边做边学尤其可行。

\zhappendix

\section{从 $X(\Fq)$ 上的函数到 $X(\CC)$ 上的 D-模}
\label{s:frmfun}

\subsection{概览}

\ssszh 本附录是第 \ref{s:bbloc} 节的补充，目的是解释：从旗簇的 $\Fq$-点上的函数过渡到旗簇的 $\CC$-点上的
D-模，何以是一件自然的事。

简短的回答是：人们把后者看作前者某种意义上的范畴化。任何解释都必然涉及几何表示论中三个重要的
主题——D-模、函数-层对应，以及 Riemann--Hilbert 对应。

我们希望下面关于这些主题的讲述，能为非专家提供足够的背景与直觉，以便欣赏正文中的讨论。
文献指引我们留到每一小节的末尾。

\subsection{代数 D-模}

\ssszh 本小节讨论 D-模的基础。先从仿射情形开始。

设 $X$ 是域 $k$ 上的光滑仿射代数簇，为简单起见假定 $k$ 的特征为零。粗略地说，$X$ 上的 D-模就是
这样一个对象：在其上既可以乘以 $X$ 上的函数，也可以沿 $X$ 上的向量场求导。

更形式地说，记 $O_X$ 为 $X$ 上正则函数构成的代数，$T_X$ 为 $X$ 上正则向量场构成的 $O_X$-模。
$X$ 上的 D-模 $M$ 就是带有平坦联络的 $O_X$-模。也就是说，$M$ 是向量空间，带有函数代数的作用
\[
O_X\otimes M\to M,\quad f\otimes m\mapsto f\cdot m,
\]
以及向量场李代数的作用
\[
T_X\otimes M\to M,\quad \xi\otimes m\mapsto \xi\cdot m.
\]
上面是李代数作用，这恰恰就是平坦性（即曲率消失）的假设。还要求这两种运算与「以函数缩放向量场」
相容
\[
f\cdot(\xi\cdot m)=(f\cdot\xi)\cdot m,
\]
并与「用向量场求导函数」相容，即满足「莱布尼茨法则」
\[
\xi\cdot(f\cdot m)=(\xi\cdot f)\cdot m+f\cdot(\xi\cdot m),
\]
其中 $\xi\cdot f$ 表示 $f$ 沿 $\xi$ 的导数。

等价地，可以构造 $X$ 上多项式线性微分算子构成的代数 $D_X$，即 $O_X$ 的全部自同态中由「乘以函数」
与「沿向量场求导」生成的子代数。有了它，D-模就不过是左 $D_X$-模。

\begin{ex} 对仿射空间 $X=\mathbb{A}^n_k$，正则函数与向量场形如
\[
O_X\simeq k[x_1,\ldots,x_n]\quad\text{和}\quad T_X\simeq
\underset{i}{\oplus}\ O_X\cdot\partial_i,
\]
其中 $\partial_i$（$1\leqslant i\leqslant n$）记与坐标 $x_i$ 相伴的标准坐标向量场。此时，$D_X$
作为 $k$-代数由 $x_i,\partial_j$（$1\leqslant i,j\leqslant n$）生成，满足标准的交换关系
\[
[x_i,x_j]=[\partial_i,\partial_j]=0,\qquad
[\partial_i,x_j]=\begin{cases}1 & \text{若 } i=j,\\ 0 & \text{否则}.\end{cases}
\]
具体地说，$D_X$ 以有序单项式
\[
x_1^{a_1}\cdots x_n^{a_n}\cdot\partial_1^{b_1}\cdots\partial_n^{b_n},\qquad
\text{其中 } a_i,b_j\in\mathbb{Z}^{\geqslant 0},
\]
为基，即由系数为多项式的线性微分算子构成。
\end{ex}

记 $X$ 上 D-模构成的阿贝尔范畴为 $\on{D-mod}(X)^\heartsuit$，其无界导出范畴记为
$\on{D-mod}(X)$。下面讨论这两个范畴中三个基本对象，也就是 D-模的三个例子。

\begin{ex} 微分算子代数 $D_X$ 自然地是自身的左模，即一个 D-模。
\end{ex}

下面转向一个不那么平凡、「整体」的例子。

\begin{ex} 对任意 $X$，正则函数环 $O_X$ 带有显然的 D-模结构：按标准方式乘函数、对函数求导。

该模有如下生成元与关系的表示。作用在常函数 $1$ 上给出 D-模的同构
\[
D_X/D_X\cdot T_X\simeq O_X.
\]
特别地，给出 D-模的映射 $O_X\to M$，等同于指定 $M$ 的一个平坦截面，即被 $T_X$ 零化的元素 $m$。

更一般地，任何带平坦联络 $\nabla$ 的向量丛 $E$ 都是 D-模；给出 D-模的映射 $E\to M$，等同于指定
张量积 $M\otimes_{O_X}E^\vee$（带其自然平坦联络）的一个平坦截面。
\end{ex}

最后，我们想描述一类「局部」的例子。

\begin{ex} 对 $X$ 的任意闭点 $x$，有一个相应的 delta 函数 D-模。记 $\fm_x\subset O_X$ 为在 $x$
处消失的函数构成的极大理想，令
\[
\delta'_x:=D_X/(D_X\cdot\fm_x).
\]
特别地，给出 D-模的映射 $\delta'_x\to M$，等同于指定 $M$ 的一个被 $\fm_x$ 零化的截面，
即在 $x$ 处概形论地支撑的截面。

D-模 $\delta'_x$ 是「支撑在 $x$ 处的分布构成的向量空间」的代数几何版本。\footnote{若改用相应的
右 D-模，这一类比会更贴切，但读者可以略过这一点。} 与此相关，若 $\xi_1,\ldots,\xi_n$ 是
$T_X$ 中的一族向量场，它们在 $x$ 处切空间 $T_xX$ 中的像构成一组基，则单项式
\begin{equation} \label{PBW}
\xi_1^{a_1}\cdots\xi_n^{a_n},\qquad\text{其中 } a_i\in\mathbb{Z}^{\geqslant 0}
\end{equation}
构成 $\delta'_x$ 的一组基。人们把 $1$ 在 $D_X$ 中的像想成通常的 delta 分布，整个模便由它的各阶
导数张成。

这里有一个稍微技术性的点，读者初读时可以放心略过。$\delta'_x$ 的底向量空间若用一条线扭一下会
更方便，即 $x$ 处切空间的行列式线：
\[
\delta_x:=\delta'_x\underset{k}{\otimes}\det(T_xX).
\]
理由如下。给定光滑仿射簇的映射 $f:X\to Y$，有相应的 D-模推前
$$f_*:\on{D-mod}(X)\to\on{D-mod}(Y),$$ 参见下文第 \ref{ss:pushpull} 节。把这一扭纳入之后，
便有典范同构
\[
f_*(\delta_x)\simeq\delta_{f(x)},\qquad\text{对 } x\in X.
\]
也就是说，delta 函数推前为 delta 函数，这与分布上的类比所预期的一致。
\end{ex}

\begin{re} \label{r:hwdelta} 这里给出 $\delta_x$ 的另一种构造，它有助于说明行列式线为何出现。

考虑在 $x$ 处消失的向量场构成的李代数，即 $\fm_x\cdot T_X\subset T_X$。它自然地作用在商
$T_X/\fm_x\cdot T_X\simeq T_xX$ 上。具体地说，这一作用通过
\[
\fm_x\cdot T_X/\fm_x^2\cdot T_X\simeq T^*_xX\otimes T_xX\simeq\on{End}(T^*_xX)
\]
分解，而它在 $T_xX$ 上的作用，对任意向量空间 $W$，就是自同态李代数 $\on{End}(W)$ 在对偶空间
$W^*$ 上的自然作用。

考虑 $\fm_x\cdot T_X$ 在行列式线 $\det(T_xX)$ 上的诱导作用。这条线还带有 $\OO_X$ 的自然作用，
这两个作用定义了由 $O_X$ 与 $\fm_x\cdot T_X$ 生成的 $D_X$ 的子代数 $D_{X,x}$ 的作用。要点在于：
若不先过渡到行列式，所需的种种相容性就不再成立。

最后，诱导出的 D-模带有典范同构
\[
D_X\underset{D_{X,x}}{\otimes}\det(T_xX)\simeq\delta_x,
\]
它在所生成的行列式线上作用为恒等，也就是我们已描述了 $D_{X,x}$ 在 delta 分布本身所张成的空间
上的作用。
\end{re}

正如读者可以想见，许多 D-模的行为介于上面最后两个例子之间：它们沿某个子簇像代数联络，而沿其
法向像 delta D-模。

\ssszh \label{ss:pushpull} 给定光滑仿射簇的映射 $X\to Y$，其 D-模范畴之间有自然的拉回与推前
函子，即
\[
f^!:\Dmod(Y)\to\Dmod(X)\qquad\text{和}\qquad f_*:\Dmod(X)\to\Dmod(Y).
\]
关于这两个函子的性质，就本书的目的而言，只需记住下面这一点。流形 $M$ 上有光滑函数代数及其对偶
的分布空间。给定映射 $M\to N$，可以拉回函数，并借助对偶推前分布。

直观地说，$f^!$ 记录的是：$Y$ 上满足某个线性偏微分方程的函数，其拉回满足哪些线性偏微分方程。
类似地，$f_*$ 记录的是：$X$ 上满足某个线性偏微分方程的分布，其推前满足哪些线性偏微分方程。
特别地，$f^!$ 与 $f_*$ 并不总是伴随函子，而是互为对偶。这可以用第 \ref{ss:dualizable} 节的
dg 范畴对偶性精确表述。

为方便读者，下面更详细地讨论 $f^!$ 与 $f_*$，不过读者可以放心略过。拉回 $f^!$ 就是把拟凝聚
拉回再作一个上同调平移，即
\[
f^!(M)\simeq O_X\underset{O_Y}{\otimes}M[d_X-d_Y],
\]
其中 $d_X,d_Y$ 分别记 $X,Y$ 的维数。这里 $O_X$ 以显然的方式作用，$T_X$ 则通过链法则作用，
即经由映射 $$T_X\to O_X\underset{O_Y}{\otimes}T_Y$$ 以及后者在 $f^!(M)$ 上的自然作用。要描述
$f_*$，注意上面拉回的公式可以改写为用双模
\[
O_X\underset{O_Y}{\otimes}D_Y
\]
在 $D_Y$ 上作张量。同一个双模也定义推前函子 $f_*$。更仔细地说，它给出右 D-模范畴之间的函子
$$f_{\on{r},*}:\on{mod-D}(X)\to\on{mod-D}(Y).$$ 这一差别由「左 D-模与右 D-模典范等同」来弥合。
即在光滑仿射簇 $Z$ 上，典范丛 $K_Z$ 带有 $D_Z$ 的自然右作用。粗略地说，人们把典范丛看作分布的
代数替代，而该作用就是分部积分。用它与左 D-模作张量，便定义了所需的等价
\[
-\underset{O_Z}{\otimes}K_Z:\Dmod(Z)\simeq\on{mod-D}(Z),
\]
在两边施用，便得到推前
\[
f_*:\Dmod(X)\simeq\on{mod-D}(X)\to\on{mod-D}(Y)\simeq\Dmod(Y).
\]

\ssszh 仿射情形下 D-模的定义可以直接整体化，如下。设 $X$ 是 $k$ 上的光滑代数簇，不必仿射。

$X$ 上的 D-模是带平坦联络的拟凝聚层 $\sM$。也就是说，记 $\OO_X$ 为正则函数层、$\sT_X$
为正则向量场层、$\underline{k}$ 为 $k$ 对应的常层，则 $\sM$ 是 $k$ 向量空间的层，带有映射
\[
\OO_X\underset{\ck}{\otimes}\sM\to\sM\qquad\text{与}\qquad
\sT_X\underset{\ck}{\otimes}\sM\to\sM,
\]
满足与上一小节相同的相容性。所谓拟凝聚，是进一步要求：对任意仿射开集 $U$ 及 $U$ 上正则函数
$f$（其非零轨迹记 $U_f$），自然映射
\[
\Gamma(U,\sM)\underset{O_U}{\otimes}O_{U_f}\xrightarrow{\sim}\Gamma(U_f,\sM)
\]
为同构。由这一假设，若 $X$ 仿射，则当下这个层论版本的 D-模范畴与上一小节的范畴经由「取整体
截面」函子典范等价。

于是，粗略地说，一般光滑簇上的 D-模范畴由其上各仿射开子簇的 D-模范畴粘合而成。类似地，
D-模 $\sM$ 大致就是这样一个层：其局部截面既可以乘以正则函数，也可以沿正则向量场求导。

上一小节的例子如下搬到这一设定中。

\begin{ex} 有 $X$ 上微分算子的整体层 $\sD_X$，D-模就是 $\sD_X$-模构成的拟凝聚层。
\end{ex}

\begin{ex} 对任意 $X$，正则函数层 $\OO_X$ 自然地是 D-模；任何带平坦联络的向量丛
$(\sE,\nabla)$ 的截面层亦然。
\end{ex}

\begin{ex} 对 $X$ 的闭点 $x$，同样有支撑在 $x$ 处的 delta 函数 D-模，记作 $\delta_x$。
\end{ex}

\label{ss:dmodglob}
\ssszh 最后指出：如第 \ref{ss:dmodglob} 节所讨论，$X$ 上的 D-模自然地构成一个阿贝尔范畴，
使得忘掉 D-模结构、只看阿贝尔群层的函子是 $t$-正合的。也就是说，核、余核与直和按层的通常方式
计算，它们自然地继承 D-模结构，并满足正确的泛性质。

\begin{defn} 记 $\on{D-mod}(X)^\heartsuit$ 为 $X$ 上 D-模构成的阿贝尔范畴，其无界导出范畴记为
$\on{D-mod}(X)$。
\end{defn}

给定映射 $f:X\to Y$，同样有推前与拉回函子
\[ f^!:\Dmod(Y)\to\Dmod(X)\qquad\text{和}\qquad f_*:\Dmod(X)\to\Dmod(Y), \]
定义同第 \ref{ss:pushpull} 节。应当强调：在非仿射情形，$f_*$ 是「用双模
$\sO_X\otimes_{\sO_Y}\sD_Y$ 作张量」的左导出函子与「沿 $f$ 推前层」的右导出函子的复合。因此
它不是两个 D-模阿贝尔范畴之间某个函子的导出函子，必须从一开始就在导出范畴中工作。

\ssszh 关于 D-模的进一步阅读，读者可以参阅下面这些（在众多可能的文献中，按对读者的要求大致
递增排列）：\cite{coutinho}、\cite{arkhipovnotes}、\cite{bernsteinnotes}、
\cite{kashiwarabookd-modules}、\cite{gaitsgorynotes}、\cite{hottatakeuchitanisaki}。

\subsection{函数-层对应}

\ssszh 如前所约，我们已经理解了复旗簇上的 D-模范畴
\[
\on{D-mod}(\sB(\CC)).
\]
接下来要说明，它与函数构成的向量空间
$$\Fun(\sB(\Fq))$$
在何种意义上类似。

\ssszh 回忆：对任意 $\Fq$-簇 $X$ 与特征零的域 $k$，我们给它配上
\[
\Fun(X(\Fq),k),
\]
即有理点构成的有限集 $X(\Fq)$ 上 $k$-值函数构成的向量空间。

生活中一个美妙的事实是：$\Fun(X(\Fq),k)$ 中的特殊函数，往往是 $X$ 上特殊层的「影子」。
下面把这一点说得更精确。

\ssszh\label{s:ladic} 为此，取 $k=\Qlb$，即 $\ell$-进数 $\mathbb{Q}_\ell$ 的代数闭包，其中
$\ell$ 是与 $q$ 互素的任意素数。

有了这样的 $k$，就有性质良好的「$X$ 上 $k$ 向量空间的 étale 局部系统」范畴。该范畴等价于
$X$ 的 étale 基本群在 $k$ 向量空间上的连续表示范畴。

\begin{re} 对不熟悉这些概念的读者，这里给一条帮助定位的注。设 $X$ 是光滑复代数簇，带有
局部系统 $\sL$，即解析拓扑意义下局部常值的有限维 $k$ 向量空间层。

要作的基本观察是：几乎不可能在 $X$ 的某个 Zariski 开子集上把 $\sL$ 平凡化。具体地说，
读者不妨想一想 $X=\mathbb{G}_m$、$\sL$ 是绕原点单值性为单位根的秩一局部系统的情形。

然而，若基本群的相应表示通过某个有限商分解，则拉回到 $X$ 的有限覆叠空间 $\tilde{X}$ 之后
可以把 $\sL$ 平凡化。关键在于，$\tilde{X}$ 又典范地是代数簇——我们强调这在一般情形下并非
显然。在上面的例子中，若单值性是 $n$ 次单位根，可以取覆叠
\[
\mathbb{G}_m\to\mathbb{G}_m,\quad z\mapsto z^n.
\]
覆叠映射 $\tilde{X}\to X$ 是 étale（即平坦且非分歧）映射的原型。要点是：若容许不仅在
Zariski 开集上、而且在 étale 覆叠上检验层的局部性质，就得到一套稳健的局部系统理论，它对
一般概形（不必在复数上）都有意义。
\end{re}

更一般地，若还想考虑 $X$ 的局部闭子簇上的局部系统，它们都自然地属于构造 étale 层的有界导出
范畴
$
\on{Sh}(X)^b,
$
并且事实上生成该范畴。

\ssszh 下面说明这一范畴在最简单的情形 $X$ 为一个点（即域 $\kappa$ 的谱）时是什么样子。

若 $\kappa$ 可分闭，其 $k$ 向量空间层构成的构造导出范畴就是 $k$ 向量空间的有界导出范畴，即
\[
\on{Sh}(\Sp\kappa)^b\simeq\on{Vect}_k^b.
\]
显然，这是因为 $\Sp\kappa$ 没有有趣的 étale 覆叠；这也应合读者对拓扑或复几何中「点上的
通常层」的直觉。强调 $k$ 起的是层的系数域的作用，而 $\kappa$ 提供层所依托的几何对象
$\Sp\kappa$。特别地，$\kappa$ 与 $k$ 之间不必有任何关系。

然而，对一般的域 $\kappa$，$\Sp\kappa$ 确有有趣的 étale 覆叠。例如给定域的可分扩张
$\kappa\to\kappa'$，映射
\[
\Sp\kappa'\to\Sp\kappa
\]
是 étale 的，而任何有限 étale 覆叠都是这类映射的不交并。因此，记 $\kappa_s$ 为 $\kappa$ 的
可分闭包、$H:=\on{Gal}(\kappa_s/\kappa)$ 为其 Galois 群，则 $\Sp\kappa$ 上的构造层与 $H$ 在
$k$ 向量空间上的有限维连续表示的有界导出范畴之间有典范等价，即
\[
\on{Sh}(\Sp\kappa)^b\simeq\on{Rep}_k(H)^b.
\]
特别地，étale 层范畴捕捉了域 $\kappa$ 的算术信息。通常拓扑中与 $\Sp\kappa$ 类比的对象不是点，
而是（射有限）群 $H$ 的分类空间 $\on{pt}/H$。在这一类比下，人们把
$\Sp\kappa_s\to\Sp\kappa$ 想成类似于 $\on{pt}\to\on{pt}/H$。

\label{ss:shvpt}

\ssszh 给定概形的映射 $f:X\to Y$，有推前与拉回函子\footnote{此外还有 $!$-推前与 $*$-拉回，
但我们用不到。}
\[
f_*:\on{Sh}(X)^b\to\on{Sh}(Y)^b\qquad\text{和}\qquad
f^!:\on{Sh}(Y)^b\to\on{Sh}(X)^b.
\]
这里 $f_*$ 是通常的层推前，$f^!$ 是所谓的例外逆像。

若 $f$ 是紧合映射（回忆这大致是说映射的纤维是射影簇），则 $f_*$ 左伴随于 $f^!$。特别地，
给定 $X$ 的闭点 $x$（剩余域 $\kappa_x$），有伴随对
\[
i_{x,*}:\on{Sh}(\Sp\kappa_x)^b\rightleftarrows\on{Sh}(X)^b:i_x^!.
\]
具体地说，$i_{x,*}$ 给出支撑在 $x$ 处的摩天楼层，其伴随 $i_x^!$ 把层送到它在 $x$ 处支撑的
（导出）截面。

\ssszh 仍设 $X$ 是 $\Fq$ 上的簇。现在可以描述所许诺的映射
\[
\on{FF}:\on{Sh}(X)^b\to\Fun(X(\Fq)),
\]
即 Grothendieck 的函数-层对应。对每个有理点 $x:\Sp\Fq\to X$，需要给出映射
\[
\on{Sh}(X)^b\to k,
\]
即相应函数在 $x$ 处的值。提醒读者：如第 \ref{s:ladic} 节，$k$ 记 $\ell$-进数的代数闭包
$\overline{\mathbb{Q}}_\ell$。为定义所求的求值，先过渡到 $x$ 处的 $!$-茎，即考虑复合
\[
\on{Sh}(X)^b\xrightarrow{i_x^!}\on{Sh}(\Sp\Fq)^b\simeq
\on{Rep}_k(\on{Gal}(\overline{\mathbb{F}}_q/\Fq))^b.
\]
由此回忆，$\on{Gal}(\overline{\mathbb{F}}_q/\Fq)$ 就是整数环的射有限完备化
$\widehat{\mathbb{Z}}$，其拓扑生成元是 Frobenius 自同构 $\on{Fr}$。

特别地，取 Frobenius 的迹，便得到所求的映射
\[
\on{Sh}(X)^b\xrightarrow{i^!_x}
\on{Rep}_k(\on{Gal}(\overline{\mathbb{F}}_q/\Fq))^b
\xrightarrow{\on{tr}(\on{Fr},-)} k.
\]
具体地说，它把层 $\sS$ 送到数
\[
\underset{j}{\Sigma}\ (-1)^j\cdot\tr(\on{Fr},H^j(i^!_x\sS)).
\]

\begin{re} 若读者觉得这太遥远，让我们接着第 \ref{ss:shvpt} 节的类比讲下去。把完备化
$\widehat{\mathbb{Z}}$ 换成通常的整数 $\mathbb{Z}$，注意 $\on{pt}/\mathbb{Z}$ 就是圆周
$S^1$。也就是说，可以把 $\Sp\Fq$ 想成某种圆周。

于是可以把 $n$ 维簇 $X\to\Fq$ 想成类似于 $2n+1$ 维实流形 $\pi:M\to S^1$（纤维化在圆周
之上），而其有理点则是一族嵌入的圆周，给出 $\pi$ 的截面。Frobenius 迹的类比会把 $M$ 上的
局部系统送到它绕这些圆周的单值性的迹。
\end{re}

\ssszh 函数-层对应的一个基本性质是：它把函数的 $*$-推前与层的 $*$-推前、以及相应的
$!$-拉回对应起来。对拉回这是按定义，对推前则需要 étale 上同调中的 Lefschetz 不动点公式。

此外，由于 $i^!$ 把 distinguished 三角形送到 distinguished 三角形，\footnote{回忆这只是说：
$\on{Sh}(X)^b$ 中的映射序列 $A\to B\to C$，在同构意义下来自复形的短正合列。更内蕴地说，
$A$ 是 $B\to C$ 的同伦核，等价地 $C$ 是 $A\to B$ 的同伦余核，即锥。} 我们得到 $k$ 向量
空间的映射
\[
\on{FF}:K(\on{Sh}(X)^b)\underset{\mathbb{Z}}{\otimes}k\to\Fun(X(\Fq)),
\]
其中 $K(-)$ 记 Grothendieck 群；因此可以合理地把 $\on{Sh}(X)^b$ 看作函数空间的
「范畴化」。这是个有用的观点，不过我们要略微提醒：上述箭头通常并不是同构。

\ssszh 下面描述这一对应的两个基本例子。

\begin{ex} 记 $\underline{k}$ 为 $X$ 上的常层。若 $X$ 光滑且 $n$ 维，则 $\on{FF}(\underline{k})$
是 $X(\Fq)$ 上取值 $q^n$ 的常函数。
\end{ex}

\begin{re} 由于我们用的是 $!$-茎，在任何簇 $X$ 上「提升」常函数 $1$ 的层不是常层，而是
对偶化层 $\omega_X$。按定义，它由沿投影 $\pi:X\to\Sp\mathbb{F}_q$ 拉回
$\on{Gal}(\overline{\mathbb{F}}_q/\mathbb{F}_q)$ 的平凡表示 $k$ 得到，即
$$\omega_X\simeq\pi^!(k).$$ 读者应当知道，许多作者在建立函数-层对应时改用 Verdier 对偶的
约定（即用 $*$-茎），此时与常函数 $1$ 对应的会是常层，而不是对偶化层。
\end{re}

\begin{ex} 给定有理点 $x:\Sp\Fq\to X$，考虑摩天楼层 $i_{x,*}k$，其中把 $k$ 看作平凡的
$\on{Gal}(\overline{\mathbb{F}}_q/\Fq)$-模，集中在同调次数零。则有
\[
\on{FF}(i_{x,*}k)=\delta_x,
\]
即该摩天楼层「提升」了该点的示性函数。
\end{ex}

尽管这些简单的例子希望已把基本想法传达给读者，我们还应提到：远为有趣的数与函数，例如
Gauss 和、Kloosterman 和、有限李型群的特征标，都有层论的对应物。另一个丰富的例子
（它催生了整体几何朗兰兹理论）是：函数域上非分歧 Hecke 特征函数，提升为有限域上曲线的
$G$-丛模空间上的 Hecke 特征层。

\ssszh 更多细节，读者可参阅 \cite{milnenotes}、\cite{milneetalebook}、
\cite{freitagbook}。

\subsection{Riemann--Hilbert 对应}

\ssszh 从 $\Fun(X(\Fq))$ 走到 $\on{Sh}(X)^b$ 之后，我们已能理解后者对任意基域上的 $X$
（不限于 $\Fq$）的意义。

在复数上，改用关系密切的构造解析层范畴
\[
\on{Sh}(X_{\on{an}})^b
\]
会更方便。这里 $X_{\on{an}}$ 记 $X$ 的解析化，即其闭点集 $X(\CC)$ 配上度量拓扑。具体地说，
$\on{Sh}(X_{\on{an}})^b$ 是 $k$ 向量空间层的导出范畴中，由「沿 $X$ 的某个 Zariski 局部闭
分层局部常值」的层生成的全子范畴。

\begin{re} 范畴 $\on{Sh}(X_{\on{an}})^b$ 是 étale 层相应范畴的扩大，本质上来自 Artin
比较定理。例如取 $X=\mathbb{G}_m$，étale 范畴中的秩一局部系统，其单值性必须落在代数整数
的单位群 $\overline{\mathbb{Z}}_\ell^{\times}$ 中；而在解析范畴中，单值性可以是
$\overline{\mathbb{Q}}_\ell^\times$ 的任意元素。
\end{re}

\ssszh 最后，让我们得到与 D-模的联系。先固定一个同构 $\overline{\mathbb{Q}}_\ell\simeq\CC$，
以便把 $\on{Sh}(X_{\on{an}})^b$ 与它在复向量空间系数下的对应物等同起来。

回忆 $\on{Sh}(X_{\on{an}})^b$ 包含支撑在 $X$ 的子簇上的局部系统这样的对象。通过把「在微分
方程组与其局部解之间过渡」这一操作精细化，它们可以转化到 D-模的全子范畴中，如下。

\ssszh 先考虑构造导出范畴中由 $X$ 上局部系统构成的子范畴。假定 $X$ 光滑。此时有范畴的等价
\begin{equation} \label{e:vivaflatsections}
\{\text{$X_{\on{an}}$ 上有限秩局部系统}\}\simeq
\{\text{$X_{\on{an}}$ 上的平坦全纯联络}\},
\end{equation}
其中右边记复几何意义下带全纯平坦联络 $\nabla$ 的全纯向量丛 $\sE$ 构成的范畴。

我们来回顾这一等价如何运作。从右到左，把联络 $(\sE,\nabla)$ 送到其平坦截面层
$\sE^\nabla$。要点是：对任意 $x\in X$、纤维 $\sE_x$ 中的任意向量，都在某个小的解析球内
唯一地延拓为平坦截面。这里既用到联络的平坦性，也用到允许解析解（而不只是多项式解），从而
本质上化归为常微分方程的情形。从左到右，把局部系统 $\sL$ 送到全纯丛
\[
\sL\underset{\ck}{\otimes}\OO_{X_{\on{an}}},
\]
其联络由 $\OO_{X_{\on{an}}}$ 上的显然联络诱导。

\ssszh 为与 $X$ 上的代数 D-模建立联系，注意用
$-\otimes_{\OO_X}\OO_{X_{\on{an}}}$ 作用（即不仅允许多项式截面，也允许全纯截面）给出自然
的解析化函子
\begin{equation} \label{e:gaga}
\{\text{$X$ 上的平坦代数联络}\}\to
\{\text{$X_{\on{an}}$ 上的平坦全纯联络}\}.
\end{equation}
强调：左边考虑的是代数几何性质的对象，右边是复解析性质的对象。

若 $X$ 紧合，则函子 \eqref{e:gaga} 是等价。这是所谓的 GAGA（即 géométrie algébrique et
géométrie analytique）定理的推论。

然而，若 $X$ 非紧合，函子 \eqref{e:gaga} 不再是等价。它仍然本质满，但通常会把许多不同的
平坦代数联络送到同一个全纯联络。原因很简单，在秩一联络上就已可见——通常全纯函数远多于
多项式函数（与 $X$ 紧合的情形不同），而它们满足的微分方程会在解析化之后导致不同代数联络
之间的等同。

\begin{ex} 设 $X=\mathbb{A}^1$，坐标 $z$，考虑左代数 D-模
\[
D_X/(D_X\cdot\partial_z)\qquad\text{和}\qquad D_X/(D_X\cdot(\partial_z-1)).
\]
具体地说，两者的底 $O_X$-模都通过作用在 $1$ 在 $D_X$ 中的像上，与 $O_X$ 有典范等同，并分别
带有联络
\[
\nabla=\partial_z\qquad\text{和}\qquad\nabla=\partial_z-1.
\]
作为代数 D-模它们不同构，但其解析化同构。事实上，这来自 $e^z$ 是全纯而非多项式的函数，
且在不计常数倍的意义下是
\[
\partial_z u=u
\]
的唯一解，即第二个联络的平坦截面。
\end{ex}

补救 \eqref{e:gaga} 之多对一现象的，是下面这个非平凡的观察。对给定的平坦全纯丛
$\sE_{\on{an}}$，在映到它的种种代数联络 $\sE$ 之中，存在唯一的 $\sE$，其平坦截面没有
「无穷远处的本质奇点」。更精确地说，$X$ 上的代数联络 $(\sE,\nabla)$ 称为{\em 正则的}，
若对 $X$ 的一个光滑紧化 $$j:X\to\overline{X}$$（其边界
$\overline{X}\setminus X$ 是简单正规交叉除子），$j_*\sE$ 上的联络形式 $j_*\nabla$ 在边界
除子上大致只有单极点；更仔细的表述见 \cite{hottatakeuchitanisaki}。可以证明这与
$\overline{X}$ 的选取无关。

\begin{ex} 在上面的例子中，可取紧化
\[
j:\mathbb{A}^1\to\mathbb{P}^1,
\]
即添上无穷远点。在无穷远附近用坐标 $w=z^{-1}$，两个联络分别形如
\[
\nabla=\partial_w\qquad\text{和}\qquad\nabla=\partial_w+w^{-2}.
\]
特别地，第二个联络有二阶极点，因此不正则。这等价于函数 $e^z$ 在 $\infty$ 处有本质奇点。
\end{ex}

当限制到正则联络时，解析化函子
\[
\{\text{$X$ 上的正则代数联络}\}\to
\{\text{$X_{\on{an}}$ 上的平坦全纯联络}\}
\]
是等价。由此与局部系统复合所得的等价
\begin{equation} \label{e:rhconn}
\nabla:\{\text{$X$ 上的正则代数联络}\}\simeq
\{\text{$X_{\on{an}}$ 上的有限秩局部系统}\}
\end{equation}
显式地由「取平坦解析截面」给出。这称为联络的 Riemann--Hilbert 对应。

\ssszh 最后，让我们描述 \eqref{e:rhconn} 到 $X$ 上一切构造层的推广。直观上，我们现在不仅
要互换局部系统与正则联络，还要互换支撑在子簇上的那些对象。

第一个猜测或许是：构造层与某个适当的 D-模范畴之间有阿贝尔范畴的等价。毕竟上面关于联络的
对应匹配的是阿贝尔范畴。例如，人们或会期待把常层 $\CC_X$ 与正则函数 $\OO_X$ 互换，并对
$X$ 的任意闭点 $x$ 把秩一摩天楼层 $\CC_x$ 与 delta D-模 $\delta_x$ 互换。
\label{s:cannotlit}
然而，这不可能字面成立。事实上，若 $X$ 的维数 $d>0$，看阿贝尔范畴中的同态，有
\begin{align*}
\on{Hom}_{\on{Sh}(X_{\on{an}})^\heartsuit}(\CC_X,\CC_x)\simeq\CC\quad
&\text{但}\quad\on{Hom}_{\on{D-mod}(X)^\heartsuit}(\OO_X,\delta_x)\simeq 0.\\
\intertext{为修正这一点，注意过渡到导出范畴时，有}
\on{Hom}_{\on{Sh}(X_{\on{an}})}(\CC_X,\CC_x)\simeq\CC\quad
&\text{和}\quad\on{Hom}_{\on{D-mod}(X)}(\OO_X,\delta_x)\simeq\CC[-d].
\end{align*}
建议读者验证这两条断言，例如对仿射直线 $X=\mathbb{A}^1$。无论如何，我们看到所求的推广必须
涉及上同调平移，因而必须在导出范畴的层面表述。

\ssszh 有鉴于此，我们得到下面的定义。正则完整 D-模的范畴，记作
$\on{D-mod}^{\on{rh}}(X)$，是 $\on{D-mod}(X)$ 中由「光滑子簇 $Z\subset X$ 上正则代数联络的
$*$-扩张」生成的满三角子范畴。

有一个（导出）平坦解析截面的函子，定义在全部 D-模的范畴上
\begin{equation} \label{e:fllat}
\nabla:\on{D-mod}(X)\to\{\text{$\CC$ 向量空间的层}\},
\end{equation}
我们稍后会细说。这里的主要结果是：\eqref{e:fllat} 限制到正则完整 D-模后，通过构造导出范畴
分解，并给出等价
\begin{equation}
\nabla:\on{D-mod}^{\on{rh}}(X)\simeq\on{Sh}(X_{\on{an}})^b.
\end{equation}
这称为 D-模的 Riemann--Hilbert 对应。

\ssszh 内蕴地，\eqref{e:fllat} 定义如下。全纯典范丛 $\omega_{X_{\on{an}}}$ 自然地是右
D-模，\eqref{e:fllat} 由与之作张量得到：
\[
\nabla:=\omega_{X_{\on{an}}}\underset{\sD_X}{\otimes}-.
\]
与李代数同调的情形一样，$\omega_{X_{\on{an}}}$ 的一个典范射影分解给出下面的表示。记
$\Omega^i_{\on{an}}$（$1\leqslant i\leqslant d$）为 $X$ 上全纯 $i$-形式层。特别地，注意
$$\Omega^d_{\on{an}}=\omega_{X_{\on{an}}}.$$ 于是 \eqref{e:fllat} 把集中在同调次数零的
D-模 $\sM$ 送到层的 de Rham 复形
\begin{equation} \label{e:flatsec}
\sM\to\sM\underset{\OO_X}{\otimes}\Omega^1_{\on{an}}\to\cdots\to
\sM\underset{\OO_X}{\otimes}\Omega^d_{\on{an}},
\end{equation}
其中 $\sM$ 放在同调次数 $-d$。具体地说，对 $\sM$ 的局部截面 $m$ 与 $X$ 上的局部坐标
$z_1,\ldots,z_d$，第一个微分把
\[
m\mapsto\underset{i}{\Sigma}\ (\partial_i\cdot m)\otimes dz_i,
\]
更高阶的微分类似。D-模的复形则送到相应 de Rham 双复形的全化。这一构造显然尊重拟同构，
并给出所求的函子 \eqref{e:fllat}。

\ssszh 下面以一些基本例子与性质结束这段讨论。

\begin{ex} Riemann--Hilbert 对应把结构层送到常层的一个平移，即
\begin{align*}
\nabla(\OO_X)&\simeq\CC_X[d_X].\\
\intertext{它把 delta D-模 $\delta_x$ 送到相应的摩天楼层，即}
\nabla(\delta_x)&\simeq\CC_x.\\
\intertext{读者不妨直接验证，并特别注意它们与第 \ref{s:cannotlit} 节讨论的相容性。}
\intertext{\hspace{1em}更一般地，对维数为 $d_Z$ 的光滑子簇 $i:Z\to X$ 与 $Z$ 上以 $\sL$
为平坦解析截面的正则联络 $\sE$，有}
\nabla(i_*\sE)&\simeq i_*\sL[d_Z].
\end{align*}
\label{ex:exrhs}
\end{ex}

\ssszh 也可以处理奇异簇 $X$ 上的 D-模。对任何嵌入 $X\hookrightarrow S$（$S$ 光滑），该范畴
典范地等同于 $S$ 上 D-模中集合论地支撑在 $X$ 上的全子范畴。Riemann--Hilbert 对应对奇异簇
同样成立，可由光滑情形推出。

例 \ref{ex:exrhs} 中的公式则是下述事实的特例：对任意代数簇的映射 $X\to Y$，
Riemann--Hilbert 对应互换 D-模与构造层的 $*$-推前与 $!$-拉回。

\ssszh 最后讨论阿贝尔范畴。子范畴 $\on{D-mod}^{\on{rh}}(X)$ 对 $\on{D-mod}(X)$ 上的标准
截断函子封闭，并典范地等同于相应正则完整 D-模阿贝尔范畴的导出范畴。由例 \ref{ex:exrhs}
可见，这个阿贝尔范畴并不与构造层的标准阿贝尔范畴互换，而是与反常层的阿贝尔范畴等同。

\ssszh 更多细节，读者可参阅 \cite{arkhipovnotes}、\cite{bernsteinnotes}、
\cite{hottatakeuchitanisaki}。

\vspace{.5cm}

\noindent{\bf 利益冲突声明。} 无利益冲突。


\end{document}
